\documentclass[11pt,a4paper]{article}

\usepackage[utf8]{inputenc}
\usepackage[T1]{fontenc}
\usepackage{lmodern}
\usepackage[margin=2.4cm]{geometry}
\usepackage{amsmath,amssymb,amsthm,mathtools}
\usepackage{booktabs,array,longtable,tabularx}
\usepackage{enumitem}
\usepackage{microtype}
\usepackage{xcolor}
\usepackage{listings}
\usepackage[
  colorlinks=true,
  linkcolor=blue!45!black,
  citecolor=blue!45!black,
  urlcolor=blue!45!black,
  pdftitle={Scoped O.1 on the Prescribed Cone, v12},
  pdfauthor={Ngo Cheung}
]{hyperref}

\newcommand{\RR}{\mathbb R}
\newcommand{\CC}{\mathbb C}
\newcommand{\ZZ}{\mathbb Z}
\newcommand{\NN}{\mathbb N}
\newcommand{\dd}{\mathrm d}
\newcommand{\ii}{\mathrm i}
\newcommand{\ka}{\kappa}
\newcommand{\eps}{\varepsilon}
\newcommand{\WW}{\mathcal W}
\newcommand{\Fr}{\mathrm F}
\newcommand{\CR}{\mathrm{CR}}
\newcommand{\ApTr}{\operatorname{Tr}_{\mathrm{ap}}}
\newcommand{\ResUV}{\operatorname{Res}_{\mathrm{UV}}}
\newcommand{\id}{\operatorname{id}}
\newcommand{\Tr}{\operatorname{Tr}}
\newcommand{\rank}{\operatorname{rank}}
\newcommand{\diag}{\operatorname{diag}}
\newcommand{\supp}{\operatorname{supp}}
\newcommand{\TV}{\operatorname{TV}}
\newcommand{\abs}[1]{\lvert#1\rvert}
\newcommand{\norm}[1]{\lVert#1\rVert}
\newcommand{\pair}[2]{\langle#1,#2\rangle}
\newcommand{\ket}[1]{\lvert#1\rangle}
\newcommand{\bra}[1]{\langle#1\rvert}
\newcommand{\proj}[1]{\lvert#1\rangle\langle#1\rvert}
\newcommand{\loc}[1]{\nolinkurl{#1}}
\newcommand{\OPEN}{\textsf{OPEN}}
\newcommand{\CONDITIONAL}{\textsf{CONDITIONAL}}
\newcommand{\SUPPLIED}{\textsf{SUPPLIED}}
\newcommand{\PARTIAL}{\textsf{PARTIAL}}
\newcommand{\CPIR}{\textsf{CLOSED\_PENDING\_}\allowbreak\textsf{INDEPENDENT\_REVIEW}}
\newcommand{\NOTCLOSED}{\texttt{NOT\_CLOSED}}
\newcommand{\Ccal}{\mathcal C}
\newcommand{\Dcal}{\mathcal D}
\newcommand{\Span}{\mathrm{Span}}
\newcommand{\Hreg}{(\mathrm H_{\mathrm{reg}})}
\newcommand{\HI}{(\mathrm H_I)}

\newcolumntype{L}[1]{>{\raggedright\arraybackslash}p{#1}}
\newcolumntype{Y}{>{\raggedright\arraybackslash}X}

\theoremstyle{plain}
\newtheorem{theorem}{Theorem}[section]
\newtheorem{proposition}[theorem]{Proposition}
\newtheorem{lemma}[theorem]{Lemma}
\newtheorem{corollary}[theorem]{Corollary}
\theoremstyle{definition}
\newtheorem{definition}[theorem]{Definition}
\newtheorem{assumption}[theorem]{Assumption}
\newtheorem{inputstatement}[theorem]{Supplied input}
\newtheorem{example}[theorem]{Example}
\theoremstyle{remark}
\newtheorem{remark}[theorem]{Remark}

\numberwithin{equation}{section}
\setlist{itemsep=3pt,topsep=5pt}
\setcounter{tocdepth}{2}
\emergencystretch=2em
\allowdisplaybreaks

\newcommand{\scopebox}[2]{%
  \par\medskip\noindent
  \fbox{\parbox{\dimexpr\linewidth-2\fboxsep-2\fboxrule\relax}{%
  \textbf{#1}\par\smallskip #2}}\par\medskip
}
\newcommand{\nonsource}{%
  \par\smallskip\noindent
  \emph{Host boundary.}
  The background and realization are prescribed.
  A heat coefficient, conservation identity, or trace identity does
  not imply that the corresponding object sources the cone.%
}

\lstdefinestyle{python}{
  language=Python,
  basicstyle=\ttfamily\scriptsize,
  numbers=left,
  numberstyle=\tiny\color{gray},
  keywordstyle=\color{blue!60!black},
  commentstyle=\color{green!35!black},
  stringstyle=\color{red!45!black},
  breaklines=true,
  breakatwhitespace=false,
  columns=fullflexible,
  keepspaces=true,
  showstringspaces=false,
  frame=single,
  upquote=true,
  tabsize=4
}
\lstset{style=python}

\title{\textbf{Scoped O.1 on the Prescribed Cone}\\[5pt]
\Large Apex Heat Moments, a Constant-Profile
Slavnov--Taylor Identity, and Domain--Readout Boundaries\\[5pt]
\large Version 12; revision of \texttt{XVI-O1-CONE-v11.tex}:\\
closed mixed mode functional, derived D04 pairing, 40-digit \(M_{10}\), and wall certificates}
\author{Ngo Cheung}
\date{19 October 2026}

\begin{document}
\maketitle

\begin{abstract}
This is Version~12 of the scoped O.1 cone audit; it revises
XVI-O1-CONE-v11. The cone and the CR and F realizations remain
prescribed inputs. New in v12 (Section~\ref{sec:whatv12}; hashes in
Appendix~\ref{app:v12artifacts}):

\emph{(1) A closed form for the mixed \(r^2\) mode functional}
(identified, not derived): with \(h=(\mu+\nu)/2\) and \(d=\mu-\nu\),
\(F=\tfrac{h^3}3+(\tfrac{d^2}6-\tfrac13)h+d^2(d^2-4)\Psi(h)\), where
\(\Psi=\tfrac{h^3}{48}+\tfrac{h^2}{96}-\tfrac{5h}{288}-\tfrac1{96}-\tfrac{h^2(h^2-1)}{48}\psi_1(h)\).
It agrees with all 414 independent quadratures to
\(5\times10^{-28}\) relative. It makes the large-\(k\) asymptotics
exact: only odd powers occur, and the P0 pole is \(-\alpha/40\).

\emph{(2) D04 pairing and regulator.} The frame rotation
\(e^{\pm\ii\theta}\) shifts \(k\) by \(\pm\alpha\), which is the
channel offset itself. The pairing is therefore the same-\(k\) rule,
and the v11 unit-shift alternative is excluded. The regulators
\(R_k,R_h,R_\mu,R_\nu\) differ by an explicit local counterterm,
verified to \(10^{-39}\). XVI-P's prescription fixes local terms by
normalization, so D04 is unique once \(R_h\) is declared.
\(M_{10}\) is given to 40 digits:
\(M_{10}(\tfrac3{10})=0.11574388259581850963848639109243\), and
\(M_{10}^{R_h}(1)=(1-40\log2-20\gamma)/800\).

\emph{(3) Exact \(\mathbb Q(\ii)\) certificates at the chamber walls}
\(\alpha=\tfrac13,\tfrac12,\tfrac23\): the ranks are 27, 33, 27 (CR)
and 27, 33, 30 (F). The exceptional set in \((\tfrac14,1)\) is
exactly the three walls.

Still \OPEN: the weight-\(\pm1,\pm2\) traces and their log/finite
split, the 84 mixed weight-0 sums, the jet dictionary, the 20 jet
Ward combinations, the 34-/18-dim projections, \(e_5,e_6\), and the
transverse loop momentum (feasible, not started). O.1 on the cone
remains \NOTCLOSED; \texttt{bulk\_symmetry\_restored} remains false.
\end{abstract}

\tableofcontents

%======================================================================
\section{Question, evidence, and revision contract}
\label{sec:scope}
%======================================================================

\subsection{The scoped question}

The target is transport of the separately scoped XVI-Q two-point
record to a prescribed conical product. The fixed scope is:
\begin{enumerate}[label=\textup{(\roman*)}]
\item one loop and coupling order \(O(\ka^2)\);
\item the two-point sector with vanishing declared scalar mass;
\item frozen de Donder gauge;
\item dimensional reduction, with quasi-\(d\) loop momenta and
quasi-four-dimensional vertex indices;
\item the XVI-P p22 contract and the declared graph conventions;
\item a comparison of CR and F realizations without selecting
either by a finite readout.
\end{enumerate}
The later p23 excerpt is not silently substituted for p22.
The frozen XVI-P, XVI-Q, and XVI-R records are source-qualified
inputs and comparisons \cite{CheungP,CheungQ,CheungR}.
The immediate sources are the supplied v6 manuscript and its
script collection \cite{CheungConeV6,CheungConeScripts}.

The central analytical question is not merely whether a conical
heat coefficient exists. It is whether a complete differentiated
graph, on its specified operator domain and in its specified
meromorphic prescription, reduces to that coefficient and
satisfies the required interacting identities.

\subsection{Evidence classes}

\begin{description}[style=nextline,leftmargin=0pt]
\item[Analytical proof.]
A statement proved below from the displayed definitions and
explicit mathematical inputs. Its proof does not depend on a
claim that a program was run.

\item[Declared catalogue.]
The explicit arrays in Section~\ref{sec:catalogue}. Algebraic
identities of these arrays are proved locally. Identification
with complete cone graph residues is conditional on
Assumption~\ref{ass:reduction}.

\item[Historical executed record.]
The v6 source reports executions, archived logs, and SHA-256
comparisons. Those reports are preserved as attribution to v6,
not represented as a new execution in this revision.

\item[Embedded verification program.]
Appendix~\ref{app:code} contains a consolidated, corrected program
whose purpose and acceptance conditions are stated explicitly.
It is supplied as source, not as an executed certificate.
In v8 it was extracted unchanged and executed
(Section~\ref{sec:v8response}).

\item[Executed in v8.]
A statement whose evidence is a v8 run on the I-B audit box:
a named script under \loc{code/v8/}, its return code, and a log
under \loc{logs/v8/}, all hashed in Appendix~\ref{app:v8artifacts}.
Exact runs use \(\mathbb Q(\ii)\) arithmetic; numerical runs are
labelled corroboration. An executed statement about the declared
catalogue remains conditional on Assumption~\ref{ass:reduction}.

\item[Documentary comparison.]
A version-qualified statement from another manuscript, used to
keep models and targets distinct. Such a citation is not an
additional graph hypothesis unless explicitly declared as one.
\end{description}

\scopebox{What is self-contained here}{%
The heat-moment proofs, the closed Mellin calculation, the
catalogue identities, the finite-channel separations, and the
conditional graph-reduction theorem are self-contained at the
level stated. The frozen vertex arrays and raw historical logs
are not embedded in the material supplied for this revision.
Consequently this single \LaTeX\ file is not a self-contained
rebuild of the complete loop calculation.
The source listings make routing and reduction logic inspectable;
they do not manufacture the missing input arrays.}

\subsection{What v7 changes}

The revision makes four substantive improvements.
\begin{enumerate}[label=\textup{(\alph*)}]
\item It places the audit on the I-B host and gives a standing
domain ledger.
\item It proves the D01 Mellin value in closed form, without
requiring a symbolic simplifier or numerical quadrature.
\item It corrects several overstrong or ambiguous statements:
reflection-even invariance, uniqueness of extension, mixed-key
coverage, modular certificates, and numerical evidence.
\item It promotes the continuum/preparation/instrument split and
the I-C A--F separations into the release boundary.
\end{enumerate}

The original coefficient values are retained. No new complete
cone graph value is claimed.

\subsection{What v8 changes}
\label{sec:whatv8}

Version~8 keeps the v7 structure and every v7 proof. It adds
executed evidence, and every executed statement names its
artifact (Appendix~\ref{app:v8artifacts}).
\begin{enumerate}[label=\textup{(\alph*)}]
\item The v7 listing was extracted and run. Each of the fifteen
v7 corrections receives a verdict
(Section~\ref{sec:v8response}).
\item An exact, complete span certificate replaces the earlier
modular and random-rank evidence. It fixes the count of 26 Ward
combinations and the inventory dimension, each at named
\((\alpha,\text{branch})\) (Section~\ref{sec:v8cert}).
\item The 34-candidate basis is projected onto its twelve
constant-profile coordinates. The \(hh\) part of the
18-dimensional engine map is made explicit on six of its 18
coordinates, and four of the six exact values are fixed
(Section~\ref{sec:v8map}).
\item Hidden mixed-channel pairs are found. As a result the v6
\(hh\) branch difference becomes conditional. A Duhamel lemma
and a reduction formula for mixed traces are proved
(Section~\ref{sec:v8mixed}).
\item Evanescent pole coefficients are computed under a declared
slice-DRED prescription (Section~\ref{sec:v8evan}).
\item The open-obligation table and the release theorem are
re-stated with the labels \OPEN, \PARTIAL, \CONDITIONAL{} and
\(\texttt{CLOSED\_PENDING\_INDEPENDENT\_REVIEW}\).
\end{enumerate}
No v6 number changes value. One v6 statement changes status, from
unconditional on the declared catalogue to conditional, as
described in (d).

\subsection{Programme non-transfer}

The cone and CR branch remain inputs. No Einstein dynamics,
physical tensor Hessian, spacetime emergence, gravitational Page
curve, or inherited quantum coefficient is constructed.
XVI-N is not discharged; the downstream XV-B admission exit is
not opened. The historical H2 ending is retained
version-specifically, the narrower I2 construction is not
withdrawn, and the full I3/J2 characteristic maps remain separate
\cite{CheungN,CheungCharacteristic,CheungXVB}.

The extended Hawk05 discussion supplied with the revision request
concerns finite access and Page proxies. It is relevant only to
the reporting boundary below; it is not imported as an apex
calculation or as gravitational dynamics
\cite{CheungHawk05}.

%======================================================================
\section{The prescribed I-B host and the typed domain ledger}
\label{sec:host}
%======================================================================

\subsection{Geometry}

The Euclidean spectral host is
\begin{equation}
 X_\alpha=C_\alpha\times\RR^2,\qquad
 \dd s_E^2=\dd r^2+\alpha^2r^2\dd\varphi^2
             +\dd y_0^2+\dd y_3^2,
 \quad 0\le\varphi<2\pi,\quad \alpha>0.
 \label{eq:Euclidean}
\end{equation}
The transverse measure is
\[
 \dd V_\perp=\alpha r\,\dd r\,\dd\varphi.
\]
The apex stratum in the completed product is
\[
 \WW=\{r=0\}\times\RR^2.
\]
For the deficit comparison we use \(0<\alpha\le1\), with deficit
\(2\pi(1-\alpha)\).

The Lorentzian regular host has \(r>0\) and metric
\begin{equation}
 \dd s_L^2=\dd t^2-\dd r^2-\alpha^2r^2\dd\varphi^2-\dd z^2.
 \label{eq:Lorentzian}
\end{equation}
This is the I-B prescribed conical host, with an overall
signature convention chosen to agree with the pole catalogue
\cite{CheungIB}. The causal geometry is unchanged by the
opposite overall signature used in some companion descriptions.

\begin{proposition}[The regular host is not globally hyperbolic]
\label{prop:nonGH}
The spacetime \eqref{eq:Lorentzian} with its apex removed is not
globally hyperbolic.
\end{proposition}

\begin{proof}
Fix \(r_*>0\), an angle, and a longitudinal position. Choose
\(a>r_*\), and let \(p\) and \(q\) be the corresponding points at
times \(-a\) and \(a\).
For any sequence \(r_j\downarrow0\), the points at time zero and
radius \(r_j\), with the same angle and longitudinal position,
lie in \(J^+(p)\cap J^-(q)\): the radial spatial distance is
less than \(a\).
They have no convergent subsequence in the regular host.
The causal diamond is therefore not compact.
\end{proof}

This is the same geometric obstruction recorded in I-B-G
\cite[\loc{prop:nonGH}]{CheungIBG}. A free inverse is accordingly
tied to an extension and a test domain. Ordinary globally
hyperbolic uniqueness must not be substituted for those data.

The apex is a worldsheet defect, not a derived horizon. There is
no state-dependent metric update.

\subsection{Standing host and domain ledger}

\begingroup
\small
\renewcommand{\arraystretch}{1.12}
\begin{longtable}{@{}L{.20\linewidth}L{.22\linewidth}
                      L{.26\linewidth}L{.25\linewidth}@{}}
\caption{Standing comparison ledger. Companion claims are used
only within the indicated source scope.}
\label{tab:domain}\\
\toprule
Input or comparison & Source & Permitted use & Excluded inference\\
\midrule
\endfirsthead
\toprule
Input or comparison & Source & Permitted use & Excluded inference\\
\midrule
\endhead
\bottomrule
\endfoot

Prescribed cone, deficit, and holonomy &
I-B \cite{CheungIB} &
Host geometry and installed parameters &
A horizon, metric update, or matter-driven sourcing law.\\

Scalar Friedrichs Green operator and finite speed &
I-B-G, \loc{thm:green}, \loc{thm:speed}
\cite{CheungIBG} &
Free inverse and support control on its declared domain &
An exact sequence on ordinary compact-spacetime tests;
\loc{rem:compact-exact} explicitly retains that gap.\\

Maxwell--Friedrichs photon with scalar ghost &
I-B-G-S, \loc{thm:spin1}
\cite{CheungIBGS} &
Source-tested free reference &
An interacting Slavnov--Taylor identity or discharge of a
named cone Ward combination.\\

CR source-test selection &
I-B-G-S, \loc{def:geo}, \loc{thm:selection},
\loc{thm:spin2} \cite{CheungIBGS} &
Declared mathematical comparison branch &
A physical smooth-core selection theorem.\\

Positive tensor-core exclusion of idealized CR &
I-B-G-S-B, \loc{cor:exclude}, \loc{thm:contrast},
\loc{prop:warddefect} \cite{CheungTensor} &
Source-qualified comparison for its declared tensor families &
An independently completed tensor proof in this paper, or a
choice made by the present ST row.\\

Friedrichs Wick square, stress, and holonomy-odd current &
I-B-VP and I-B-VP-1 \cite{CheungVP,CheungVP1} &
One-point and compact-test infrared comparison &
Absence of infrared poles in massless loop graphs, or a
two-point graph prescription.\\

Frozen-source Maxwell response &
I-B-VP-4 \cite{CheungVP4} &
Classical Maxwell comparison, if such a detector is declared &
An interacting Ward identity.\\

Smooth-core finite parts and scalar/Dirac selection &
I-B-VP-5-Ex1-B, Ex2, Ex2-B
\cite{CheungEx1B,CheungEx2,CheungEx2B} &
Qualifications of core limits and finite-part universality &
An interacting spin-2 CR graph calculation.\\
\end{longtable}
\endgroup

The CR/F row proved below is labelled against this table:
it is a constant-profile catalogue identity on two prescribed
branches. It is neither the free photon theorem nor a tensor-core
selection theorem.

\begin{remark}[Sourcing is a separate equation]
I-B's \loc{prop:not-source} separates symmetry, conservation, and
trace conditions from an Einstein sourcing equation
\cite{CheungIB}. For example, on a flat chart a nonzero constant
symmetric conserved traceless tensor need not equal the
Einstein tensor of that chart.
The same refusal applies to every heat coefficient below.
\end{remark}

\subsection{Self-adjointness is not destructive access}

For a radial Bessel expression
\[
 L_\nu=-\partial_r^2-r^{-1}\partial_r+\nu^2r^{-2}
\]
in \(L^2((0,\infty),r\,\dd r)\), the endpoint pairing is
\begin{equation}
 \mathfrak b(u,v)=
 \lim_{r\downarrow0}r
 \bigl(\overline u\,v'-\overline{u'}\,v\bigr).
 \label{eq:boundary}
\end{equation}

\begin{proposition}[Closed endpoint dynamics]
\label{prop:unitary}
A self-adjoint realization \(A\) generates a unitary
Schr\"odinger group \(e^{-\ii tA}\). On sufficiently regular
domain solutions its total endpoint probability flux vanishes.
A bound eigenstate is not an absorbing output channel.
\end{proposition}

\begin{proof}
The spectral theorem gives unitarity. Equivalently,
\[
 \frac{\dd}{\dd t}\norm{u(t)}^2
 =2\operatorname{Re}\pair{u(t)}{-\ii Au(t)}=0.
\]
The endpoint current is proportional to
\(\operatorname{Im}(r\overline u\,u')\), which vanishes in the
total boundary pairing on a self-adjoint domain.
An eigenstate evolves by a phase.
\end{proof}

These are standard endpoint facts
\cite{DerezinskiRichard,KayStuder}.
They do not identify ghosts or gauge-fixed tensor components
with a positive physical probability space.

Self-adjointness alone does not imply positive-energy wave
stability: a negative eigenvalue \(-\sigma^2\) produces
\(\ddot a-\sigma^2a=0\). Conversely, reduced access can be lossy
while an enlarged complete dynamics remains unitary.
Neither an absorbing boundary nor an apex reservoir is introduced
here.

%======================================================================
\section{Twisted scalar channels and the normalized apex trace}
\label{sec:heat}
%======================================================================

\subsection{Channels and subtraction}

For \(b\in[0,1)\), the twisted Friedrichs transverse Laplacian
\(A_b^\Fr\) has radial orders
\[
 \nu_n=\frac{\abs{n+b}}{\alpha},\qquad n\in\ZZ.
\]
It chooses \(r^{\nu_n}\) at zero, with the no-logarithm condition
when \(\nu_n=0\).

Its diagonal heat kernel, relative to \(\dd V_\perp\), is
\begin{equation}
 K_b^\Fr(s;r,\varphi;r,\varphi)
 =\frac{e^{-z}}{4\pi\alpha s}
   \sum_{n\in\ZZ}I_{\abs{n+b}/\alpha}(z),
 \qquad z=\frac{r^2}{2s}.
 \label{eq:kernel}
\end{equation}
We use the Bessel heat representation and standard modified
Bessel identities \cite{DLMF,DerezinskiRichard}.

The infinite cone has infinite ordinary heat trace. The apex
functional is therefore defined by subtraction, not by writing an
unqualified \(L^2\)-trace.

\begin{definition}[Abel-subtracted apex functional]
\label{def:Abel}
For \(u,s>0\), let
\begin{align}
 T_b(u)
 &=\int_{C_\alpha}
 e^{-(\cosh u-1)r^2/(2s)}
 K_b^\Fr(s;x,x)\,\dd V_\perp,\\
 T_{\mathrm{bulk}}(u)
 &=\frac{\alpha}{2(\cosh u-1)}.
\end{align}
Define
\[
 \ApTr e^{-sA_b^\Fr}
 =\lim_{u\downarrow0}
       \bigl(T_b(u)-T_{\mathrm{bulk}}(u)\bigr)
\]
when the limit exists.
\end{definition}

\begin{theorem}[Normalized twisted apex constant]
\label{thm:C}
For every \(\alpha>0\), \(b\in[0,1)\), and \(s>0\),
\begin{equation}
 \ApTr e^{-sA_b^\Fr}
 =\Ccal_b^\Fr
 =\frac{1}{12}\left(\frac1\alpha-\alpha\right)
       -\frac{b(1-b)}{2\alpha}
 =\frac{B_2(b)}{2\alpha}-\frac{\alpha}{12},
 \label{eq:C}
\end{equation}
where \(B_2(b)=b^2-b+\tfrac16\).

For a compact proper-time superposition
\[
 f(\lambda)=\int_{s_-}^{s_+}w(s)e^{-s\lambda}\,\dd s,
 \qquad 0<s_-<s_+<\infty,\quad w\in L^1,
\]
the corresponding apex functional satisfies
\begin{equation}
 \ApTr f(A_b^\Fr)=\Ccal_b^\Fr f(0).
 \label{eq:Cfunctional}
\end{equation}
\end{theorem}

\begin{proof}
For \(\rho>-1\) and \(u>0\),
\[
 \int_0^\infty e^{-z\cosh u}I_\rho(z)\,\dd z
 =\frac{e^{-\rho u}}{\sinh u}.
\]
The nonnegative Bessel summands permit interchange with the
radial integral. Summing the geometric series gives
\begin{equation}
 T_b(u)=
 \frac{\cosh((\tfrac12-b)u/\alpha)}
      {2\sinh u\,\sinh(u/(2\alpha))}.
 \label{eq:Tclosed}
\end{equation}
Its Laurent expansion is
\[
 T_b(u)=\frac{\alpha}{u^2}
       +\frac{B_2(b)}{2\alpha}-\frac{\alpha}{6}
       +O(u^2),
\]
whereas
\[
 T_{\mathrm{bulk}}(u)
 =\frac{\alpha}{u^2}-\frac{\alpha}{12}+O(u^2).
\]
Subtraction proves \eqref{eq:C}. The change of radial variable
removed \(s\), so the result is independent of proper time.

On the compact proper-time interval the subtracted expressions
are uniformly bounded for sufficiently small \(u\). Dominated
convergence and linearity give
\[
 \ApTr f(A_b^\Fr)
 =\Ccal_b^\Fr\int_{s_-}^{s_+}w(s)\,\dd s.
\]
This is \eqref{eq:Cfunctional}.
\end{proof}
\nonsource

The untwisted result and the conical heat-kernel method are
classical \cite{Fursaev1994,FursaevMiele,Vassilevich}.
The calculation above fixes the normalization actually used here.
It is not an unrestricted ordinary-trace theorem for arbitrary
resolvent products or differentiated graph kernels.

\subsection{An image contribution is not a diffracted-front report}

\begin{example}[A pure-image apex contribution]
\label{ex:image}
At \(\alpha=\tfrac12\), \(b=0\),
\[
 \Ccal_0^\Fr=\frac18.
\]
The nonidentity image in the two-fold plane quotient contributes
\[
 \frac1{4\pi s}e^{-r^2/s}.
\]
Integrating over its wedge of angle \(\pi\) gives
\[
 \frac{\pi}{4\pi s}
 \int_0^\infty r e^{-r^2/s}\,\dd r=\frac18.
\]
Thus an apex heat constant can be nonzero in an image
representation without requiring a separately nonzero strictly
diffracted-front contribution.
\end{example}

The image-isolated windows and short-window deficit sensitivity
of I-B-ECHO concern another functional
\cite{CheungEcho}. They do not identify a window ratio with
\(\Ccal_b^\Fr\) or with a graph residue. This separation is part
of the release theorem, not merely an informal analogy.

%======================================================================
\section{Exceptional Bessel realizations and CR/F coefficients}
\label{sec:branches}
%======================================================================

\subsection{A one-channel replacement}

For \(0<\nu<1\), let \(A_\nu^+\) and \(A_\nu^-\) be the
nonnegative homogeneous regular and singular Bessel realizations
in \(L^2(r\,\dd r)\). Their radial heat kernels are
\[
 k_\nu^\pm(s;r,r')
 =\frac1{2s}e^{-(r^2+r'^2)/(4s)}
          I_{\pm\nu}\!\left(\frac{rr'}{2s}\right).
\]
The endpoint choices, rather than a periodic relabelling of an
ordinary fractional holonomy, distinguish them
\cite{DerezinskiRichard}.

\begin{theorem}[Relative diagonal heat shift]
\label{thm:flip}
For \(0<\nu<1\),
\begin{equation}
 \int_0^\infty
 \bigl(k_\nu^-(s;r,r)-k_\nu^+(s;r,r)\bigr)r\,\dd r
 =\nu.
 \label{eq:flip}
\end{equation}
The relative diagonal integral is absolutely convergent.
\end{theorem}

\begin{proof}
It equals
\[
 \frac12\int_0^\infty
 e^{-z}\bigl(I_{-\nu}(z)-I_\nu(z)\bigr)\,\dd z.
\]
At zero the worst behavior is \(z^{-\nu}\), which is integrable.
At infinity use
\[
 I_{-\nu}(z)-I_\nu(z)
 =\frac{2\sin(\pi\nu)}{\pi}K_\nu(z).
\]
The resulting integrand decays exponentially.
With the Abel factor, the integral is
\[
 \frac{e^{\nu u}-e^{-\nu u}}{2\sinh u}
 =\frac{\sinh(\nu u)}{\sinh u}.
\]
Dominated convergence gives \(\nu\).
\end{proof}
\nonsource

\subsection{Geometric helicities}

For geometric helicity \(S\), the Friedrichs fractional label is
\[
 b_\Fr(S)=\{S\alpha\}.
\]
In the comparison chamber \(1/2<\alpha<1\), put
\begin{equation}
 \nu=2-\alpha^{-1}\in(0,1).
 \label{eq:nu}
\end{equation}
For each exceptional helicity \(S=\pm2\), CR replaces one regular
channel of order \(\nu\) by the singular realization.

The per-helicity coefficients are
\begin{align}
 c_0&=\frac{1-\alpha^2}{12\alpha},&
 c_1&=\frac{5\alpha^2-6\alpha+1}{12\alpha},\nonumber\\
 c_2^\Fr&=\frac{23\alpha^2-36\alpha+13}{12\alpha},&
 c_2^\CR&=\frac{23\alpha^2-12\alpha+1}{12\alpha}.
 \label{eq:cvalues}
\end{align}
They satisfy
\begin{align}
 c_1-c_0&=\frac{\alpha-1}{2},&
 c_2^\CR-c_0&=2\alpha-1,\nonumber\\
 c_2^\Fr-c_0&=\frac{(\alpha-1)(2\alpha-1)}{\alpha},&
 c_2^\CR-c_2^\Fr&=\nu.
 \label{eq:cdifferences}
\end{align}
These follow from Theorem~\ref{thm:C} at
\(b=0,\alpha,2\alpha-1\), and from Theorem~\ref{thm:flip}.
The polynomial substitution \(b=2\alpha\) reproduces
\(c_2^\CR\), but is justified here only through that specified
channel replacement.

At \(\alpha=\tfrac12\), continuous evaluation gives
\[
 c_0=c_2^\Fr=c_2^\CR=\frac18,\qquad c_1=-\frac18.
\]
For \(0<\alpha<\tfrac12\), the Friedrichs weight-two label is
\(2\alpha\), not \(2\alpha-1\); the displayed F polynomial must
not be extended unchanged across that chamber boundary.

\begin{remark}[The CR flat limit is not an endpoint realization]
\label{rem:flat}
As \(\alpha\uparrow1\),
\[
 c_0,c_1,c_2^\Fr\longrightarrow0,\qquad c_2^\CR\longrightarrow1.
\]
This is a one-sided polynomial limit. At \(\nu=1\),
\(r^{-1}\notin L^2(r\,\dd r)\) near zero, so the deficient
singular branch used in the proof is no longer available in
the same form. No CR operator value at the exactly smooth
endpoint is inferred. Spin-dependent nonuniform smooth limits
also occur in the established conical heat-kernel literature
\cite{FursaevMiele}.
\end{remark}

%======================================================================
\section{The weighted apex moment}
\label{sec:weighted}
%======================================================================

\begin{definition}[Weighted Abel subtraction]
For \(u,s>0\), define
\begin{align*}
 T_b^{(2)}(u)
 &=\int_{C_\alpha}r^2
 e^{-(\cosh u-1)r^2/(2s)}
 K_b^\Fr(s;x,x)\,\dd V_\perp,\\
 T_{\mathrm{bulk}}^{(2)}(u)
 &=\frac{\alpha s}{(\cosh u-1)^2}.
\end{align*}
Their subtracted limit defines
\(\ApTr[r^2e^{-sA_b^\Fr}]\).
\end{definition}

\begin{theorem}[Normalized \(r^2\)-weighted coefficient]
\label{thm:D}
For \(\alpha>0\), \(b\in[0,1)\), and \(s>0\),
\begin{equation}
 \ApTr[r^2e^{-sA_b^\Fr}]=s\,\Dcal_b,
 \qquad
 \Dcal_b=
 -\frac{11\alpha}{180}
 +\frac{B_2(b)}{3\alpha}
 -\frac{B_4(b)}{6\alpha^3},
 \label{eq:D}
\end{equation}
where
\(B_4(b)=b^4-2b^3+b^2-\tfrac1{30}\).
Equivalently,
\begin{equation}
 \Dcal_b=
 \frac{1+10\alpha^2-11\alpha^4
       -60\alpha^2b(1-b)-30b^2(1-b)^2}
      {180\alpha^3}.
 \label{eq:Dexpanded}
\end{equation}
For the proper-time superpositions of
Theorem~\ref{thm:C},
\begin{equation}
 \ApTr[r^2f(A_b^\Fr)]=-\Dcal_b f'(0).
 \label{eq:Dfunctional}
\end{equation}
\end{theorem}

\begin{proof}
The radial variable \(z=r^2/(2s)\) gives
\[
 \frac{T_b^{(2)}(u)}{s}
 =-\frac{2}{\sinh u}\frac{\dd}{\dd u}T_b(u).
\]
To expose the constant term, write
\[
 T_b(u)=\frac{\alpha}{u^2}+t_0+t_2u^2+O(u^4),
\]
where expansion of \eqref{eq:Tclosed} yields
\[
 t_2=\frac{B_4(b)}{24\alpha^3}
     -\frac{B_2(b)}{12\alpha}
     +\frac{7\alpha}{360}.
\]
Therefore
\[
 \frac{T_b^{(2)}(u)}s
 =\frac{4\alpha}{u^4}-\frac{2\alpha}{3u^2}
   +\frac{7\alpha}{90}-4t_2+O(u^2).
\]
The bulk expression is
\[
 \frac{\alpha}{(\cosh u-1)^2}
 =\frac{4\alpha}{u^4}-\frac{2\alpha}{3u^2}
   +\frac{11\alpha}{180}+O(u^2).
\]
The pole terms cancel and the constant difference is
\eqref{eq:D}. Expanding the Bernoulli polynomials gives
\eqref{eq:Dexpanded}.

Dominated convergence on the compact proper-time support gives
\[
 \ApTr[r^2f(A_b^\Fr)]
 =\Dcal_b\int s\,w(s)\,\dd s=-\Dcal_b f'(0).
\]
\end{proof}
\nonsource

\begin{proposition}[Weighted singular-channel shift]
\label{prop:Dflip}
Replacing one regular channel of order \(0<\nu<1\) by its
singular realization changes the weighted coefficient by
\begin{equation}
 \Delta\Dcal=\frac{2\nu(1-\nu^2)}3.
 \label{eq:Dflip}
\end{equation}
In the weight-two comparison chamber,
\[
 \Dcal_{S=2}^\Fr=\Dcal_{2\alpha-1},
 \qquad
 \Dcal_{S=2}^\CR
 =\Dcal_{2\alpha-1}+\Delta\Dcal
 =\Dcal_{2\alpha}.
\]
The last expression is polynomial evaluation, not an assertion
that \(2\alpha\) is a fractional twist.
\end{proposition}

\begin{proof}
For one channel the weighted relative Abel integral divided by
\(s\) is
\[
 -\frac1{\sinh u}\partial_u
       \left(\frac{2\sinh(\nu u)}{\sinh u}\right).
\]
Its limit is \eqref{eq:Dflip}. The weighted diagonal relative
integral is convergent at zero because its Bessel integrand is
\(O(z^{1-\nu})\), and it is exponentially convergent at infinity.
Substitution in \eqref{eq:Dexpanded} proves the final equality.
\end{proof}

In particular,
\[
 \Dcal_0
 =\frac{(1-\alpha^2)(1+11\alpha^2)}{180\alpha^3},
 \quad
 \Dcal_{S=1}=\Dcal_\alpha,
 \quad
 \Dcal_{S=2}^\Fr=\Dcal_{2\alpha-1},
 \quad
 \Dcal_{S=2}^\CR=\Dcal_{2\alpha}.
\]
All four tend to zero as \(\alpha\uparrow1\).
Thus the weighted CR polynomial has no flat remnant.
This vanishing does not itself establish convergence of the
singular operator family at the endpoint.

\begin{theorem}[Exact orbifold and twisted-orbifold identities]
\label{thm:orbifold}
For every integer \(N\ge2\) and \(j\in\{0,\ldots,N-1\}\),
\begin{equation}
 \Dcal_{j/N}(1/N)
 =\frac1{4N}\sum_{k=1}^{N-1}
       \frac{\cos(2\pi jk/N)}{\sin^4(\pi k/N)}.
 \label{eq:orbifold}
\end{equation}
Consequently,
\begin{equation}
 \Dcal_{j/N}(1/N)
 =
 \frac{(N^2-1)(N^2+11)
       -60j(N-j)-30j^2(N-j)^2}{180N}.
 \label{eq:orbifoldclosed}
\end{equation}
\end{theorem}

\begin{proof}
On the rotation quotient, the \(k\)-th nonidentity image has
diagonal kernel
\[
 (4\pi s)^{-1}
       e^{-r^2\sin^2(\pi k/N)/s}.
\]
The character contributes the corresponding holonomy phase.
Pairing conjugate images replaces the phase by its cosine.
The wedge has angle \(2\pi/N\), and direct integration gives
\[
 \frac{2\pi/N}{4\pi s}
 \int_0^\infty r^3
 e^{-r^2\sin^2(\pi k/N)/s}\,\dd r
 =\frac{s}{4N\sin^4(\pi k/N)}.
\]
The identity image is precisely the bulk subtraction.
The character sector has orders \(\abs{Nn+j}\), hence
\(\alpha=1/N\), \(b=j/N\).
Theorem~\ref{thm:D} proves \eqref{eq:orbifold}.
Substitution in \eqref{eq:Dexpanded} proves
\eqref{eq:orbifoldclosed}.
\end{proof}

Examples include
\[
 \Dcal_0(1/2)=\frac18,\qquad
 \Dcal_{1/3}(1/3)=-\frac4{27},\qquad
 \Dcal_{1/4}(1/4)=-\frac1{16},\qquad
 \Dcal_{2/5}(1/5)=-\frac{16}{25}.
\]

\begin{remark}[Corrected literature normalization]
\label{rem:Fursaevnorm}
Fursaev uses the full opening angle \(\theta=2\pi\alpha\).
His coefficient \(C_4^{\mathrm{Fur}}(\theta)\), in
equation (2.15) of \cite{Fursaev1994}, is
\[
 C_4^{\mathrm{Fur}}(2\pi\alpha)
 =\frac{(\alpha^{-2}-1)(\alpha^{-2}+11)}{90}.
\]
The normalization relation is
\[
 \boxed{\ \Dcal_0=\frac{\alpha}{2}
                  C_4^{\mathrm{Fur}}(2\pi\alpha).\ }
\]
This replaces the memory-qualified identification in v6.
It follows also by inserting \(f=r^2\) in that reference's
weighted heat expansion. No numerical spin-two coefficient is
imported from the literature.
\end{remark}

%======================================================================
\section{The D01 derivative obstruction and its Mellin value}
\label{sec:Mellin}
%======================================================================

\subsection{The actual obstruction}

Let \(\nu\) be \eqref{eq:nu}. The exceptional relative radial
kernel is
\[
 \Delta k_s(r,r')
 =\frac{\sin(\pi\nu)}{\pi s}
   e^{-(r^2+r'^2)/(4s)}
   K_\nu\!\left(\frac{rr'}{2s}\right).
\]
Consider the raising expression
\[
 X=\partial_r-\frac{\nu}{r},
 \qquad
 \mathcal R_s(r)=
 \left[X_rX_{r'}\Delta k_s(r,r')\right]_{r'=r}.
\]

\begin{theorem}[Local obstruction]
\label{thm:obstruction}
For every \(0<\nu<1\),
\begin{equation}
 \mathcal R_s(r)
 =a_0\,r^{-2\nu-2}+O(r^{-2\nu}),
 \qquad
 a_0=\frac{2\nu^2(4s)^\nu}
               {s\,\Gamma(1-\nu)}>0.
 \label{eq:obstruction}
\end{equation}
Thus the ordinary diagonal integral
\(\int_0^\epsilon r\mathcal R_s(r)\,\dd r\) diverges.
The expression \(X\) does not map the full singular Bessel
domain into the next \(L^2\) channel.
\end{theorem}

\begin{proof}
The leading singular part of \(\Delta k_s\) is
\[
 \frac{(4s)^\nu}{2s\,\Gamma(1-\nu)}(rr')^{-\nu}.
\]
Applying \(X_r\) and \(X_{r'}\) multiplies it by
\(4\nu^2/(rr')\), proving \eqref{eq:obstruction}.
Moreover
\[
 X(r^{-\nu})=-2\nu r^{-\nu-1},
 \qquad
 \int_0^\epsilon r\abs{r^{-\nu-1}}^2\,\dd r=\infty.
\]
This excludes the unrestricted domain-preserving intertwining
statement. It does not exclude all densely restricted operators
or every possible distributional extension.
\end{proof}

\subsection{A closed Mellin formula}

\begin{theorem}[Exact Mellin continuation]
\label{thm:Mellin}
For \(\operatorname{Re}z>2\nu\), define
\[
 M_s(z)=\int_0^\infty r^{1+z}\mathcal R_s(r)\,\dd r.
\]
With \(t=z/2\), its meromorphic continuation is
\begin{equation}
 \boxed{
 M_s(z)=
 \frac{\sin(\pi\nu)}{2\sqrt\pi}\,
 s^{t-1}\,
 \frac{t\,\Gamma(t+\nu+2)\Gamma(t-\nu)}
      {\Gamma(t+\tfrac32)}.
 }
 \label{eq:Mclosed}
\end{equation}
It is regular at \(z=0\) for \(0<\nu<1\), and
\begin{equation}
 M_s(0)=0,\qquad
 M_s'(0)=-\frac{\nu+1}{2s}.
 \label{eq:Mzero}
\end{equation}
Possible poles are contained in
\[
 \{2\nu-2j:j\in\NN_0\}
 \ \cup\
 \{-2\nu-4-2j:j\in\NN_0\}.
\]
Some listed poles can cancel.
\end{theorem}

\begin{proof}
Put \(x=r^2/(2s)\) and
\(A=\sin(\pi\nu)/(\pi s)\).
Direct differentiation gives
\[
 \mathcal R_s(r)=\frac{A}{2s}e^{-x}H_\nu(x),
\]
where
\[
 H_\nu=
 xK-2xK'+K'+xK''
       -2\nu(K'-K)+\frac{\nu^2}{x}K,
 \qquad K=K_\nu(x).
\]
The Bessel equation
\(xK''+K'=(x+\nu^2/x)K\) reduces this to
\[
 H_\nu=2\left[(x+\nu)K-(x+\nu)K'
                         +\frac{\nu^2}{x}K\right].
\]

For \(\operatorname{Re}\mu>\abs{\operatorname{Re}\nu}\), let
\begin{equation}
 I_\nu(\mu)
 =\int_0^\infty x^{\mu-1}e^{-x}K_\nu(x)\,\dd x
 =\frac{\sqrt\pi\,\Gamma(\mu+\nu)\Gamma(\mu-\nu)}
        {2^\mu\Gamma(\mu+\tfrac12)}.
 \label{eq:IK}
\end{equation}
For completeness, use
\[
 K_\nu(x)=\frac12\int_0^\infty
 y^{\nu-1}e^{-x(y+y^{-1})/2}\,\dd y.
\]
Integrating in \(x\) first gives
\[
 I_\nu(\mu)
 =2^{\mu-1}\Gamma(\mu)
   \int_0^\infty
      \frac{y^{\mu+\nu-1}}{(1+y)^{2\mu}}\,\dd y.
\]
The beta integral and the gamma duplication formula give
\eqref{eq:IK}. This also identifies the \(a=1\) limit of the
standard Laplace--Bessel formula in DLMF 10.43.22
\cite{DLMF}.

In the initial convergence half-plane, integration by parts
has no endpoint contribution. It yields
\[
 M_s(z)=A(2s)^t
 \bigl[(t+1)I_\nu(t+1)+\nu(\nu+t)I_\nu(t)\bigr].
\]
Since
\[
 \frac{I_\nu(t+1)}{I_\nu(t)}
 =\frac{t^2-\nu^2}{2t+1},
\]
the bracket is
\[
 \frac{t(t+\nu)(t+\nu+1)}{2t+1}I_\nu(t).
\]
Substitution of \eqref{eq:IK} proves \eqref{eq:Mclosed}
on a nonempty convergence half-plane and hence by meromorphic
continuation.

At \(t=0\), the gamma factors are finite for \(0<\nu<1\),
and the explicit factor \(t\) gives \(M_s(0)=0\).
Using
\[
 \Gamma(\nu+2)\Gamma(-\nu)
 =-\frac{\pi(\nu+1)}{\sin(\pi\nu)}
\]
gives the derivative in \eqref{eq:Mzero}.
The pole statement follows from the numerator gamma functions.
\end{proof}

\subsection{What is and is not unique}

\begin{remark}[Correction of the extension claim]
\label{rem:extension}
The \emph{Mellin prescription} is unambiguous once its initial
integral and meromorphic continuation are specified.
On smooth rotationally invariant transverse tests, the local
Taylor powers are even. The singular heads
\(r^{-2\nu-2+2j}\) therefore have no resonance at \(z=0\)
for \(0<\nu<1\). In particular, \(\nu=\tfrac12\) is not a
resonance of this even radial test class.

This does \emph{not} prove that the full heat-time-dependent
diagonal has a unique scale-covariant extension.
At fixed \(s\), \(\mathcal R_s\) is not a homogeneous function
of \(r\). If \(\delta_{\mathrm{rad}}\) denotes the endpoint
evaluation functional relative to radial measure, then adding
\[
 \frac{\chi(\nu)}s\,\delta_{\mathrm{rad}}
\]
preserves parabolic dimensional covariance but changes the
integrated finite part.
Thus scale covariance alone does not select the Mellin value.

The earlier script comment ``unique iff \(2\nu\) is not an
integer'' refers to a different unrestricted half-line Taylor
class and must not be used as the uniqueness theorem for smooth
two-dimensional radial tests.
\end{remark}

\begin{corollary}[The remaining D01 prescription obligation]
\label{cor:D01}
The forbidden composite is exactly Mellin-null under the
prescription in Theorem~\ref{thm:Mellin}. If the complete D01
graph reduction identifies this composite and the frozen
prescription selects that continuation without additional
contact terms, it supplies no correction to the declared
constant-profile catalogue.

The italicized identification and prescription statements remain
part of the graph-specific contract. They are not proved by
\(M_s(0)=0\).
\end{corollary}

No uniform domain theorem at \(\alpha=\tfrac12\) or
\(\alpha=1\) follows from this interior-chamber Mellin identity.
Continuity of a selected polynomial is not uniform meromorphic
control of the joint endpoint family.

%======================================================================
\section{Conditional reduction of a complete graph}
\label{sec:reduction}
%======================================================================

\subsection{Constant profiles are local jets}

A constant transverse profile means a local jet constant near
the apex, with longitudinal Fourier dependence. It is not a
globally square-integrable constant section of an infinite cone.

The helicity-zero tensor restriction is
\[
 h_{11}=h_{22}=f_0,
\]
together with the longitudinal symmetric block.
The individual coordinate directions \(h_{11}\), \(h_{22}\),
and \(h_{12}\) are useful array diagnostics, but are not
separately globally admissible constant helicity-zero tensor
profiles on a generic cone.

A closed transverse loop has one spectral trace. Once that trace
has been reduced to a function of one identified operator, it
receives one coefficient \(c_S\), not a product of one heat
coefficient for each propagator.

\begin{assumption}[Graph-specific reduction contract]
\label{ass:reduction}
For a specified complete graph \(G\), external jet class,
routing, and branch, require:
\begin{enumerate}[label=\textup{(R\arabic*)}]
\item \textbf{Complete graph data.}
All vertices, propagator numerators, source conventions, phases,
statistics signs, symmetry factors, and independent-coordinate
normalizations are specified and matched to the frozen contract.

\item \textbf{Operator reduction.}
The fully contracted transverse expression is reduced to a finite
sum
\[
 f_{G,S}(A_S;l,p),
\]
on identified channel domains, plus a remainder proved to have
zero apex UV residue.

\item \textbf{Domain and trace legitimacy.}
Differentiated intertwiners, integrations by parts, and trace
rearrangements are valid on those domains.
Every boundary or finite-part defect is included or proved absent
in the complete permitted sum.

\item \textbf{UV/IR separation.}
An auxiliary infrared prescription distinguishes the UV residue
from massless infrared poles. A scaleless integral is not assigned
its UV pole merely by discarding its infrared part.

\item \textbf{Meromorphic interchange.}
The apex heat identity extends to the regulated spectral
functions, and the transverse subtraction, longitudinal
integration, cutoff removal, and Laurent-residue operation
commute in the stated prescription.
\end{enumerate}
\end{assumption}

\begin{theorem}[Conditional spectral apex reduction]
\label{thm:reduction}
Under Assumption~\ref{ass:reduction},
\begin{equation}
 \ResUV G_{\mathrm{ap}}(p)
 =
 \sum_S c_S\,
 \ResUV\left[
 \mu^{4-d}\int
 \frac{\dd^{\,d-2}l_\parallel}{(2\pi)^{d-2}}
 f_{G,S}(0;l_\parallel,p)\right].
 \label{eq:graphreduction}
\end{equation}
The factors, multiplicities, phases, and statistics are included
in \(f_{G,S}\).
\end{theorem}

\begin{proof}
Retain the auxiliary proper-time and infrared cutoffs.
By (R2), the complete graph is a finite sum of functions of
identified channel operators, modulo a residue-free remainder.
Apply Theorem~\ref{thm:C} and the channel replacement of
Theorem~\ref{thm:flip} to each regulated heat superposition.
This replaces its transverse apex functional by
\(c_S f_{G,S}(0;l,p)\).

Condition (R3) ensures that no boundary or trace defect has been
lost. Condition (R5) permits the stated integration and residue
operations. The remainder has zero residue by (R2), giving
\eqref{eq:graphreduction}.
\end{proof}

Zero transverse external momentum does not prove (R2).
A mixed-domain bubble or a domain-changing differential operator
need not be a function of one \(A_S\).

The v6 source reports that the archived engine data discharge
(R1), and that parts of (R3) were checked
\cite{CheungConeV6}. The supplied script listings allow inspection
of the declared routing and prefactors. Without the frozen input
files, their completeness and byte identity cannot be
independently re-established from the v7 material alone.
v8 replays the frozen engine (read-only copies, tree hashes in
Section~\ref{sec:v8hash}) but does not certify completeness.
In all cases, (R2), (R4), and (R5) remain hypotheses.

\subsection{Longitudinal UV normalization}

Write
\[
 \eps=4-d,\qquad D=d-2=2-\eps.
\]

\begin{lemma}[Angular coefficient of a simple UV pole]
\label{lem:UV}
Suppose an infrared-regular Euclidean longitudinal integrand has
\[
 I(R\omega;p)=R^{-2}g(\theta;p)+O(R^{-2-\delta}),
 \qquad \delta>0,
\]
uniformly in angle, and the dimensional angular continuation is
regular at \(D=2\).
Then its \(1/\eps\) coefficient is
\begin{equation}
 \frac1{(2\pi)^2}\int_0^{2\pi}g(\theta;p)\,\dd\theta.
 \label{eq:UV}
\end{equation}
A prescribed Wick phase multiplies this coefficient separately.
\end{lemma}

\begin{proof}
Above a fixed large radius,
\[
 \int_{R_0}^\infty R^{D-1}R^{-2}\,\dd R
 =\int_{R_0}^\infty R^{-1-\eps}\,\dd R
 =\frac{R_0^{-\eps}}{\eps}.
\]
The remainder is UV integrable. At \(D=2\), the measure supplies
\((2\pi)^{-2}\dd\theta\).
\end{proof}

An \(O(\eps)\) angular or numerator correction does not change
this simple-pole coefficient \emph{provided there is no hidden
higher pole, singular projection, or evanescent mixing that
invalidates that premise}. It can change finite parts.
Consistent dimensional reduction is not automatic preservation
of a proposed boundary identity \cite{Stockinger}.

%======================================================================
\section{The declared constant-profile catalogue}
\label{sec:catalogue}
%======================================================================

\subsection{Coordinate and normalization conventions}

Use the independent-coordinate order
\[
\begin{split}
0&\leftrightarrow00,\quad
1\leftrightarrow01,\quad
2\leftrightarrow02,\quad
3\leftrightarrow03,\quad
4\leftrightarrow11,\\
5&\leftrightarrow12,\quad
6\leftrightarrow13,\quad
7\leftrightarrow22,\quad
8\leftrightarrow23,\quad
9\leftrightarrow33.
\end{split}
\]
For longitudinal momentum,
\[
 p_\mu=(q_0,0,0,q_3),\qquad
 p^\mu=(q_0,0,0,-q_3),\qquad
 P^2=q_0^2-q_3^2.
\]
Let \(c_2\) denote either \(c_2^\Fr\) or \(c_2^\CR\), and put
\begin{equation}
 A=c_0+2c_1+4c_2,\quad B=-c_0+2c_1,\quad
 W=\frac12(c_0+2c_1+c_2).
 \label{eq:ABW}
\end{equation}

The hat in \(\widehat B_G\) denotes a declared pole record.
It does not remove Assumption~\ref{ass:reduction}.
The historical statement that the catalogue matches the v4
spectral logs is attributed to the v6 execution record
\cite{CheungConeV6}.

\subsection{Scalar sector}

\begin{center}
\begin{tabular}{@{}ll@{}}
\toprule
Graph & Declared coefficient of \(1/(4-d)\)\\
\midrule
D04 & \(\displaystyle-\frac{\ka^2}{4\pi}c_0P^2\)\\[4pt]
D05a & \(\displaystyle-\frac{\ka^2}{4\pi}(c_2-c_0)P^2\)\\[4pt]
Sum & \(\displaystyle-\frac{\ka^2}{4\pi}c_2P^2\)\\
\bottomrule
\end{tabular}
\end{center}

The supplied Q02 comparison gives symmetry factors
\(1\) for D04 and \(1/2\) for D05a and D12
\cite{CheungQ,CheungConeV6}. Those source values are not
re-fitted to the displayed poles.

\subsection{Tensor bubbles}

Every entry in Table~\ref{tab:hh} is multiplied by
\(\ka^2/(16\pi)\). Transposes are understood, together with the
\(1\leftrightarrow2\) images.

\begingroup\small
\begin{longtable}{@{}L{.28\linewidth}L{.63\linewidth}@{}}
\caption{Declared tensor-bubble entries. The \((5,5)\) entries
were added to the printed catalogue in v6.}
\label{tab:hh}\\
\toprule
Graph and key & Coefficient with \(\ka^2/(16\pi)\) removed\\
\midrule
\endfirsthead
\toprule
Graph and key & Coefficient with \(\ka^2/(16\pi)\) removed\\
\midrule
\endhead
\bottomrule
\endfoot

D11: \((4,4),(4,7),(7,7)\) &
\(\tfrac12c_0P^2\)\\[3pt]

D03: \((0,0),(0,9),(9,9)\) & \(-c_0P^2\)\\
D03: \((3,3)\) & \(4c_0P^2\)\\
D03: \((0,4)\) & \(c_0(q_0^2-3q_3^2)\)\\
D03: \((3,4)\) & \(4c_0q_0q_3\)\\
D03: \((4,9)\) & \(-c_0(3q_0^2-q_3^2)\)\\
D03: \((4,4),(7,7)\) & \(-(c_0+4c_1)P^2\)\\
D03: \((4,7)\) & \(-c_0P^2\)\\
D03: \((5,5)\) & \(-8c_1P^2\)\\[3pt]

D01: \((0,0),(0,9),(9,9)\) & \(c_0P^2\)\\
D01: \((3,3)\) & \(-4c_0P^2\)\\
D01: \((4,4),(7,7)\) & \((c_0-4c_1+c_2)P^2\)\\
D01: \((4,7)\) & \((c_0+c_2)P^2\)\\
D01: \((3,4)\) & \(-4(c_0+2c_2)q_0q_3\)\\
D01: \((4,9)\) & \(Aq_0^2-Bq_3^2\)\\
D01: \((0,4)\) & \(-Bq_0^2+Aq_3^2\)\\
D01: \((5,5)\) & \(-8c_1P^2\)\\
\end{longtable}
\endgroup

On the physical constant scalar jet,
\[
 B_{f_0f'_0}=B_{44}+B_{47}+B_{74}+B_{77}.
\]
No extra off-diagonal multiplicity is applied to an array already
differentiated in independent coordinates.

\subsection{The D02 and D12 tadpoles}

The nonzero D02 coordinate entries, including their transposes,
are the following, with common factor \(\ka^2W/\pi\):
\begin{align*}
 &(01)(01),(02)(02):q_3^2,
 &&(13)(13),(23)(23):q_0^2,\\
 &(12)(12):P^2,
 &&(11)(22):-\tfrac12P^2,\\
 &(01)(13),(02)(23):-q_0q_3,
 &&(03)(11),(03)(22):q_0q_3,\\
 &(00)(11),(00)(22):-\tfrac12q_3^2,
 &&(11)(33),(22)(33):-\tfrac12q_0^2.
\end{align*}
There are five diagonal and eighteen off-diagonal entries,
hence 23 nonzero coordinate entries.
The declared D12 array is
\[
 \widehat B_{\mathrm{D12}}=0.
\]
These are catalogue statements, not a claim that every massless
scaleless prescription has been independently audited.

\subsection{Ghost and source--ghost arrays}

Define
\begin{equation}
 R=3c_0+4c_1,\qquad
 T=4c_0+3c_1,\qquad
 V=c_1-c_0.
 \label{eq:RTV}
\end{equation}
For rows and columns \(0,3\),
\begin{equation}
 \widehat B_{\mathrm{D06}}
 =\frac{\ii\ka^2}{8\pi}RP^2
   \begin{pmatrix}1&0\\0&1\end{pmatrix}.
 \label{eq:D06}
\end{equation}
For source rows \(00,03,33\) and upper ghost columns \(0,3\),
\begin{equation}
 \widehat\Gamma_{\mathrm{D13}}
 =\frac{\ii\ka^2}{16\pi}
 \begin{pmatrix}
 -2Tq_0&2Vq_3\\
 -Rq_3&Rq_0\\
 -2Vq_0&2Tq_3
 \end{pmatrix}.
 \label{eq:D13}
\end{equation}
The transverse source rows \(11,22\), for these two ghost
columns, are declared zero in the v6 reduction.
This is the precise source-trace datum used below.

\subsection{Vector-type keys are excluded, not zero}

The supplied spectral script logs and skips mixed-channel
monomials. The v6 assessment locates them on:
\begin{itemize}
\item D01 vector-type tensor keys involving
\(01,02,13,23\);
\item D06 transverse ghost keys \(11,22\);
\item D13 keys \((01)c^1,(02)c^2,(13)c^1,(23)c^2\).
\end{itemize}

\scopebox{Mandatory import boundary}{%
Mixed-helicity monomials occur on vector-type keys that are
excluded from the released cone catalogue, not set to zero.
Any later sheet importing this catalogue must preserve this
sentence. A partial same-channel sum on an excluded key is not
a complete graph value.}

There is a programming distinction here: the historical
\loc{continue} statement skips a \emph{monomial}, not necessarily
the whole affected key. The corrected replay in
Appendix~\ref{app:code} quarantines the entire key whenever a
nonzero mixed-channel pole is found.

\subsection{Endpoint checks}

\begin{proposition}[Catalogue checks]
\label{prop:endpoints}
Within the comparison chamber and its stated continuous
endpoints:
\begin{enumerate}[label=\textup{(\roman*)}]
\item all displayed F entries vanish at \(\alpha=1\);
\item CR and F entries agree at \(\alpha=\tfrac12\);
\item the displayed rational coefficients are continuous at
\(\alpha=\tfrac12,\tfrac23,\tfrac45\);
\item D06 and D13 are branch independent;
\item D04+D05a has the scalar sum printed above.
\end{enumerate}
\end{proposition}

\begin{proof}
At \(\alpha=1\), \(c_0=c_1=c_2^\Fr=0\).
At \(\alpha=\tfrac12\), \(c_2^\CR=c_2^\Fr\).
There is no zero denominator at the three specified openings.
The ghost arrays contain no \(c_2\), and the scalar sum is
immediate.
\end{proof}

These checks are not uniform domain estimates and do not prove
\(\Hreg\).

%======================================================================
\section{Reflection-even invariant structure of the \(hh\) array}
\label{sec:structures}
%======================================================================

\subsection{The symmetry class}

Let \(h,h'\) be the longitudinal symmetric \(2\times2\) blocks,
and let \(f_0,f'_0\) be the transverse trace coordinates.
Write
\[
 \operatorname{tr}h=\eta^{ab}h_{ab},\qquad
 h\cdot h'=h^{ab}h'_{ab},\qquad
 php=p^ap^bh_{ab},
 \quad a,b\in\{0,3\}.
\]
The reflection-even invariant structures quadratic in \(p\) are
\begin{align}
 S_1&=P^2\operatorname{tr}h\operatorname{tr}h',&
 S_2&=P^2h\cdot h',\nonumber\\
 S_3&=(php)\operatorname{tr}h'
             +\operatorname{tr}h(ph'p),\nonumber\\
 S_6&=P^2(f_0\operatorname{tr}h'
             +\operatorname{tr}h\,f'_0),&
 S_7&=(php)f'_0+f_0(ph'p),\nonumber\\
 S_8&=P^2f_0f'_0.
 \label{eq:structures}
\end{align}

\begin{proposition}[Six-dimensional reflection-even sector]
\label{prop:six}
A slot-symmetric, reflection-even \(O(1,1)\)-invariant bilinear
of the displayed fields, quadratic in \(p\), is a unique linear
combination of the six structures in \eqref{eq:structures}.
\end{proposition}

\begin{proof}
Reflection evenness excludes a single Levi-Civita tensor; a pair
reduces to metric contractions. For two longitudinal symmetric
tensors, the metric contractions are generated by
\[
 P^2\operatorname{tr}h\operatorname{tr}h',\quad
 P^2h\cdot h',\quad
 (php)\operatorname{tr}h'
       +\operatorname{tr}h(ph'p),
\]
and the symmetrized contraction \(p\,hh'\,p\).
For the mixed matrices \(H=\eta^{-1}h\), \(H'=\eta^{-1}h'\),
the polarized Cayley--Hamilton identity is
\[
 HH'+H'H
 =(\operatorname{tr}H)H'
  +(\operatorname{tr}H')H
  +\bigl(\operatorname{tr}(HH')
          -\operatorname{tr}H\operatorname{tr}H'\bigr)I.
\]
It removes the last contraction.
A longitudinal tensor paired with a scalar gives \(S_6,S_7\);
two scalars give \(S_8\).
Independence follows by comparing the coefficients of the
independent field coordinates and \(q_0,q_3\).
\end{proof}

\begin{remark}[Why \(SO(1,1)\) alone is insufficient]
Without reflection evenness, pseudoscalar terms such as
\[
 \bigl(\epsilon^{ab}p_a h_{bc}p^c\bigr)f'_0
 +f_0\bigl(\epsilon^{ab}p_a h'_{bc}p^c\bigr)
\]
are allowed by proper longitudinal Lorentz transformations.
Thus the unqualified assertion that six structures classify
\emph{all} \(SO(1,1)\)-invariant bilinears is too strong.
The frozen metric-contraction sector used here is
reflection even.
\end{remark}

\subsection{Coefficients}

Write the restricted graph bilinear as
\[
 \widehat B_G(h,f_0;h',f'_0)
 =\frac{\ka^2}{\pi}\sum_i\beta_i^G S_i.
\]
The coefficient table is
\begin{center}
\scriptsize
\setlength{\tabcolsep}{4pt}
\begin{tabular}{@{}lcccccc@{}}
\toprule
Graph & \(\beta_1\)&\(\beta_2\)&\(\beta_3\)&
\(\beta_6\)&\(\beta_7\)&\(\beta_8\)\\
\midrule
D01&
\(-c_0/16\)&\(c_0/8\)&\(0\)&
\(-c_0/8-c_1/4-c_2/2\)&
\(c_0/4+c_2/2\)&
\(c_0/4-c_1/2+c_2/4\)\\
D02&\(0\)&\(0\)&\(0\)&\(W\)&\(-W\)&\(-W\)\\
D03&
\(c_0/16\)&\(-c_0/8\)&\(0\)&
\(3c_0/8\)&\(-c_0/4\)&\(-c_0/4-c_1/2\)\\
D11&\(0\)&\(0\)&\(0\)&\(0\)&\(0\)&\(c_0/8\)\\
D12&\(0\)&\(0\)&\(0\)&\(0\)&\(0\)&\(0\)\\
\midrule
Sum&\(0\)&\(0\)&\(0\)&
\(3(c_0+c_1)/4\)&\(-c_0/2-c_1\)&
\(-3c_0/8-2c_1-c_2/4\)\\
\bottomrule
\end{tabular}
\end{center}

\begin{proposition}[Catalogue decomposition and branch content]
\label{prop:decomposition}
The table is an exact identity of the declared coordinate arrays,
for arbitrary formal \(c_0,c_1,c_2\).
For the complete \(hh\) sum:
\begin{enumerate}[label=\textup{(\roman*)}]
\item the purely longitudinal structures cancel;
\item
\[
 \frac{\ka^2}{\pi}(\beta_6+\beta_7)
 =-\frac{\ka^2V}{4\pi};
\]
\item the CR-minus-F difference is
\begin{equation}
 \Delta\widehat B^{hh}
 =-\frac{\nu\ka^2}{4\pi}P^2f_0f'_0.
 \label{eq:branchdifference}
\end{equation}
\end{enumerate}
\end{proposition}

\begin{proof}
Substitute the definitions of \(S_i\) and compare with
Table~\ref{tab:hh} and the D02 entries. This is finite polynomial
coefficient comparison; no operator reduction is used.

Summing the rows gives the final row of the table.
Its first three coefficients vanish, and
\[
 \beta_6+\beta_7=\frac14(c_0-c_1).
\]
Only the total \(\beta_8\) depends on \(c_2\), with coefficient
\(-1/4\). Since \(c_2^\CR-c_2^\Fr=\nu\),
\eqref{eq:branchdifference} follows.
\end{proof}

For D01 alone the branch difference has additional \(S_6,S_7\)
parts. The D02 difference is
\[
 \frac{\nu\ka^2}{2\pi}(S_6-S_7-S_8),
\]
and cancels those parts. D12 is zero in the declared catalogue.
Thus the complete branch comparison includes the tadpoles.

The difference \eqref{eq:branchdifference} is gauge invariant on
this restricted lattice because the permitted longitudinal gauge
variation has \(\delta f_0=0\). It is a two-point branch
difference, not a stress tensor or a physical core selection.

%======================================================================
\section{Ghost equation and the constant-profile ST row}
\label{sec:ST}
%======================================================================

\subsection{The ghost equation}

For a symmetric source tensor \(Y\), define the longitudinal
de Donder contraction
\begin{align}
 (\mathsf F_LY)_0&=
 \frac{q_0}{2}(Y_{00}+Y_{33})-q_3Y_{03},\nonumber\\
 (\mathsf F_LY)_3&=
 q_0Y_{03}-\frac{q_3}{2}(Y_{00}+Y_{33}).
 \label{eq:FL}
\end{align}
The full contraction is
\[
 (\mathsf FY)_\rho
 =p^\mu Y_{\mu\rho}
  -\frac12p_\rho\eta^{\mu\nu}Y_{\mu\nu},
\]
so
\[
 (\mathsf FY)_0=(\mathsf F_LY)_0
       +\frac{q_0}{2}(Y_{11}+Y_{22}),
 \quad
 (\mathsf FY)_3=(\mathsf F_LY)_3
       +\frac{q_3}{2}(Y_{11}+Y_{22}).
\]

\begin{theorem}[Declared projected ghost identity]
\label{thm:ghost}
The arrays \eqref{eq:D06}--\eqref{eq:D13} satisfy
\begin{equation}
 \widehat B_{\mathrm{D06}}
 =-2\eta_\parallel\mathsf F_L
       \widehat\Gamma_{\mathrm{D13}},
 \qquad \eta_\parallel=\diag(1,-1).
 \label{eq:ghostidentity}
\end{equation}
With the declared zero transverse D13 source rows for columns
\(0,3\), \(\mathsf F=\mathsf F_L\) on this matrix.
\end{theorem}

\begin{proof}
Since \(T+V=R\), direct multiplication gives
\[
 \mathsf F_L\widehat\Gamma_{\mathrm{D13}}
 =\frac{\ii\ka^2}{16\pi}RP^2
   \begin{pmatrix}-1&0\\0&1\end{pmatrix}.
\]
Multiplication by \(-2\eta_\parallel\) proves the identity.
The full-contraction statement uses exactly the two declared
transverse source rows.
\end{proof}

The identification with a graph antighost equation additionally
uses the source normalization, Fourier convention, and
Assumption~\ref{ass:reduction}. This is not a test of the full
interacting basis.

\subsection{Independent-coordinate weights}

The ghost is upper-indexed, and
\[
 \delta_ch_{\mu\nu}=p_\mu c_\nu+p_\nu c_\mu,
 \qquad c_\nu=\eta_{\nu\nu}c^\nu.
\]
For an invariant quadratic action written in independent
coordinates,
\begin{equation}
 \frac{\partial S_{\mathrm{inv}}}{\partial h_{(mn)}}
 =(2-\delta_{mn})E^{mn}(h).
 \label{eq:coordinateweight}
\end{equation}
Thus the lower-index \(03\) component carries
\[
 2E^{03}=-2E_{03}.
\]
The sign is not an optional positive multiplicity.

The supplied v4 scan varied only
\(w\in\{\tfrac12,1,2\}\), so it omitted \(w=-2\)
\cite{CheungConeScripts}. This diagnoses its reported failure,
without interpreting that historical residual as an anomaly.

\begin{inputstatement}[Tree normalization, not loop fitting]
\label{input:tree}
In the frozen conventions, the tree source is
\[
 \frac{\partial S}{\partial K'_{(mn)}}
 =\ii A_{\mathrm{GH}}
      (p_mc_n+p_nc_m),
 \qquad A_{\mathrm{GH}}=\frac12,
\]
and the invariant quadratic normalization relative to the
displayed \(E\) is \(N=1\).
Accordingly,
\begin{equation}
 \lambda_{\mathrm{pred}}
 =\frac{N}{\ii A_{\mathrm{GH}}}=-2\ii.
 \label{eq:lambdapred}
\end{equation}
These values are attributed to the frozen tree data and the v6
normalization check, not derived from the loop row
\cite{CheungP,CheungConeV6}.
\end{inputstatement}

The optional tree routine in Appendix~\ref{app:code} checks the
quadratic normalization on nonzero de Donder configurations if
the frozen engine is supplied. It is not represented as run here.

\subsection{An explicit algebraic closure theorem}

On the axisymmetric lattice write
\[
 h=(h_{00},h_{03},h_{33},f_0),\qquad
 c=(c^0,c^3),\qquad
 p\cdot c=q_0c^0+q_3c^3.
\]
Use the lower-index Einstein expression
\begin{equation}
\begin{split}
 E_{\mu\nu}(h)
 ={}&P^2h_{\mu\nu}
 -p_\mu p^\rho h_{\rho\nu}
 -p_\nu p^\rho h_{\rho\mu}
 +p_\mu p_\nu\operatorname{tr}_4h\\
 &+\eta_{\mu\nu}
 \bigl(p^\rho p^\sigma h_{\rho\sigma}
                   -P^2\operatorname{tr}_4h\bigr).
\end{split}
\label{eq:Einstein}
\end{equation}
This normalization is fixed for the present row. It is not a
physical tensor-core Hessian theorem.

Let
\begin{align}
 \mathcal B(c,h)
 &=\widehat B^{hh}(\delta_ch,h),\nonumber\\
 \mathcal G(c,h)
 &=\sum_{(mn),\rho}
 \widehat\Gamma^{\mathrm{D13}}_{(mn),\rho}c^\rho
          (2-\delta_{mn})E^{mn}(h),
 \label{eq:BG}
\end{align}
where
\[
 \widehat B^{hh}
 =\mathrm{D01+D02+D03+D11+D12}
\]
is the declared lattice array.

\begin{theorem}[Constant-profile \(hh/Kc\) identity]
\label{thm:ST}
For arbitrary formal \(c_0,c_1,c_2\),
\begin{equation}
 \mathcal B(c,h)
 =-\frac{\ka^2V}{2\pi}P^2f_0(p\cdot c),
 \qquad
 \mathcal G(c,h)
 =\frac{\ii\ka^2V}{4\pi}P^2f_0(p\cdot c).
 \label{eq:BGexplicit}
\end{equation}
Hence
\begin{equation}
 \mathcal B+\lambda_{\mathrm{pred}}\mathcal G=0.
 \label{eq:ST}
\end{equation}
For CR and F with \(1/2<\alpha<1\) and nonzero coupling,
\(-2\ii\) is the unique constant \(\lambda\) for which
\(\mathcal B+\lambda\mathcal G\) vanishes as a polynomial in all
the declared fields, ghosts, and momenta.
\end{theorem}

\begin{proof}
The longitudinal gauge variation satisfies
\[
 \delta_cf_0=0,\qquad
 \operatorname{tr}(\delta_ch)=2(p\cdot c),
 \qquad
 p(\delta_ch)p=2P^2(p\cdot c).
\]
By Proposition~\ref{prop:decomposition}, only
\(\beta_6+\beta_7=(c_0-c_1)/4\) contributes.
This proves the formula for \(\mathcal B\).

In two longitudinal dimensions the longitudinal part of
\eqref{eq:Einstein} vanishes when \(f_0=0\).
Its contributions from the transverse trace are
\[
 E_{00}=-2q_3^2f_0,\quad
 E_{03}=-2q_0q_3f_0,\quad
 E_{33}=-2q_0^2f_0.
\]
After raising and independent-coordinate weighting they become
\[
 E^{\mathrm{coord}}_{00}=-2q_3^2f_0,\quad
 E^{\mathrm{coord}}_{03}=4q_0q_3f_0,\quad
 E^{\mathrm{coord}}_{33}=-2q_0^2f_0.
\]
Multiplying by \eqref{eq:D13} and using \(R-T=V\) gives the
formula for \(\mathcal G\).
Equation~\eqref{eq:ST} follows from \eqref{eq:lambdapred}.

For either physical branch,
\(V=(\alpha-1)/2\ne0\) in the open comparison chamber.
The polynomial \(P^2f_0(p\cdot c)\) is not identically zero.
Thus the scalar equation
\(-1/2+\lambda\ii/4=0\) forces \(\lambda=-2\ii\).
\end{proof}

This proof establishes the identity of the declared arrays
without fitting \(\lambda\) or invoking an unavailable log.
Its complete-graph interpretation remains conditional.

\begin{corollary}[Matter row and UNIT]
The declared matter row D11+D12 is transverse on this lattice.
For the UNIT substitution \(c_0=c_1=c_2=1\),
both \(\mathcal B\) and \(\mathcal G\) vanish separately.
UNIT therefore cannot determine \(\lambda\).
\end{corollary}

\begin{proof}
D11 contributes only \(S_8\), and D12 is zero.
For UNIT, \(V=0\), so \eqref{eq:BGexplicit} gives the result.
\end{proof}

The normalized quadratic convention is important:
if \(Q(h)=\tfrac12\widehat B^{hh}(h,h)\), then
\(\delta Q=\widehat B^{hh}(\delta h,h)\).
An extra factor two would change the printed convention constant.

\subsection{The three identity classes}

\begin{center}
\small
\begin{tabularx}{\linewidth}{@{}L{.25\linewidth}Y@{}}
\toprule
Class & Status and target\\
\midrule
Free source-tested &
The Maxwell--Friedrichs identity
\(C_1G^\pm_{(1)}=G^\pm_{(0)}C_1\) is a free-domain comparison
from I-B-G-S \cite{CheungIBGS}. It is not counted as a cone
interacting-basis test.\\[3pt]

Constant-profile interacting catalogue &
The ghost row and Theorem~\ref{thm:ST} are exact identities of
the stated restricted arrays, with graph identification
conditional. No certified map assigns them to a named member
of the programme's 26 combinations.\\[3pt]

Full interacting basis &
The 34-dimensional projection and 26 Ward combinations are the
O.1 target. They have not been constructed or evaluated here.
No missing combination is reported as zero.\\
\bottomrule
\end{tabularx}
\end{center}

The historical UNIT expectation also needs care. An identity
holding only after full loop integration need not survive
restriction to \(l_\perp=0\). For rapidly decreasing \(F\),
\[
 \int\dd l\,\dd^2k\,
       \partial_{k_1}(k_1F(l,k))=0,
\]
but its restriction at \(k=0\) is \(F(l,0)\), whose longitudinal
integral need not vanish.
Theorem~\ref{thm:ST} avoids this shortcut by direct algebra.

%======================================================================
\section{Heat moments, one-point fields, and missing two-point data}
\label{sec:missing}
%======================================================================

\subsection{The object ledger}

\begin{center}
\small
\begin{tabularx}{\linewidth}{@{}L{.23\linewidth}L{.31\linewidth}Y@{}}
\toprule
Object & Definition or source & What it does not establish\\
\midrule
Apex heat moment &
\(\Ccal_b,\Dcal_b\), Theorems~\ref{thm:C} and \ref{thm:D} &
A stress tensor, Landauer heat, or complete differentiated graph.\\
One-point field observable &
Friedrichs Wick square, stress, and current in I-B-VP/I-B-VP-1
\cite{CheungVP,CheungVP1} &
A massless loop UV/IR theorem or two-point normalization.\\
Two-point graph pole &
A graph satisfying (R1)--(R5) &
A physical core selection or an Einstein source law.\\
Two-point branch difference &
Equation~\eqref{eq:branchdifference} &
A one-point stress or a tensor-core limit.\\
\bottomrule
\end{tabularx}
\end{center}

\subsection{Constant-profile restriction is not injective}

\begin{proposition}[A transverse-jet kernel]
\label{prop:noninjective}
On a smooth radial external-jet class containing constant and
quadratic transverse jets, restriction of apex bilinear
functionals to constant profiles need not be injective.
\end{proposition}

\begin{proof}
Consider
\[
 L(F,G)=\int_\WW
       \bigl(F\Delta_\perp G+G\Delta_\perp F\bigr)\,\dd^2y.
\]
It vanishes for two constant transverse jets.
For \(F\) constant near the apex and \(G=r^2g(y)\),
\(\Delta_\perp r^2=4\), so suitable longitudinal tests give a
nonzero value.
Thus the restriction has a nontrivial kernel.
\end{proof}

This does not declare a complete physical counterterm basis.
It proves that the displayed restriction cannot determine every
transverse-derivative coefficient.

\subsection{The second scalar value is partial}

For a quadratic jet, \eqref{eq:Dfunctional} gives a weighted
spectral moment only after an operator reduction has justified
that insertion.
Within that reduction:
\begin{itemize}
\item the supplied D05a weighted longitudinal expression has
large-radius behavior \(O(R^{-4})\) in \(D=2\), hence no
logarithmic UV term under Lemma~\ref{lem:UV};
\item D04 requires mixed-channel traces with transverse position
operators between different channel realizations. It is not
determined by a function of one \(A_S\).
\end{itemize}
Accordingly the retained status is
\[
 \text{D05a part}=0\ \text{in the stated reduced model},
 \qquad
 \text{D04 part}=\OPEN.
\]
The second scalar basis value is not identified with its D05a
part. Neither the weighted heat theorem nor the jet/curvature
language of the tensor-core comparison computes D04.
v8 status: still \OPEN; see Section~\ref{sec:v8map} and
Lemma~\ref{lem:v8duhamel} for the ingredients supplied.

\subsection{Evanescent data and the engine map}

The simple-pole observation after Lemma~\ref{lem:UV} is a
pole-level statement under its stated hypotheses.
The \(\tilde g\)-leg arrays needed to determine evanescent
coefficients are not supplied. Those coefficients remain open.
v8 computes them under a declared slice-DRED prescription
(Section~\ref{sec:v8evan}); that is \CONDITIONAL, not an identification
with the frozen export.

The six reflection-even structures classify the executed
constant-profile restriction. They do not supply its convention
map into the programme's 18-dimensional engine basis, which also
sees transverse derivatives and non-axisymmetric data.
That map and the unseen components are open, not zero.
v8 supplies the six constant-profile coordinates of the \(hh\) block
exactly (Section~\ref{sec:v8map}); twelve remain open.

\subsection{What an exact span certificate requires}

The 34-dimensional basis and associated counts are retained as
supplied metadata, not reconstructed from the coordinate table.
A complete certificate needs the declared external-jet class,
the full matrix, its coefficient field, column completeness,
and the appropriate relation quotient.
v8 supplies these for the declared inventory at named angles and
branches (Section~\ref{sec:v8cert}).

\begin{lemma}[What modular rank can certify]
\label{lem:modular}
Let \(M\) be a specified rational matrix, and let \(p\) avoid all
denominators. Then
\[
 \rank_{\mathbb F_p}(M\bmod p)\le\rank_{\mathbb Q}M.
\]
A nonzero modular minor is an exact certificate that the
corresponding rational minor is nonzero.
If this lower bound equals an independently established upper
bound, it certifies the exact rational rank.
\end{lemma}

\begin{proof}
Clear denominators by an integer factor not divisible by \(p\).
An integer determinant nonzero modulo \(p\) is nonzero as an
integer. Apply this to a maximal nonzero modular minor.
\end{proof}

Thus modular arithmetic is not intrinsically incapable of an
exact certificate. The unresolved issue here is the absence of
the complete specified matrix and its completeness and upper-rank
arguments, not a general prohibition on modular methods.
For rational-function matrices, a good specialization can prove
generic nonvanishing, but does not settle exceptional parameter
values without additional analysis.

I-B-G's compact-test exactness warning is a useful free-field
comparison: an identity on one spectral test domain is not
automatically exactness on another test space
\cite[\loc{rem:compact-exact}]{CheungIBG}.
That comparison is not a proof of a loop span statement.

%======================================================================
\section{The continuum--preparation--instrument boundary}
\label{sec:readout}
%======================================================================

\subsection{Three objects and three targets}

Any proposed finite report of an apex coefficient must separately
specify
\begin{equation}
 \boxed{\text{defined continuum coefficient}}
 \ \longrightarrow\
 \boxed{\text{acquisition or supplied preparation}}
 \ \longrightarrow\
 \boxed{\text{complete finite instrument}}.
 \label{eq:pipeline}
\end{equation}
The arrows are obligations, not constructions supplied by the
diagram. This paper studies the first object.
A program printing a rational expression is not an apparatus
acquiring an unknown continuum amplitude.

Following the distinction recorded in the supplied Oai290-A
companion, the reporting target must also be typed
\cite{CheungOai290}:
\begin{enumerate}[label=\textup{(\roman*)}]
\item recovery of an arbitrary unknown quantum code, including
references;
\item reconstruction or estimation of one named observable;
\item consistency with a separately declared geometric model.
\end{enumerate}
A supplied finite report of \(\Dcal_b\) is not automatically even
an acquisition theorem for target (ii), and is not target (i).
No external catalogue family supplies an apex graph premise.

The no-acquisition warning attributed to X-J-ST-2, and the
window/report distinction in I-B-ECHO, preserve precisely this
boundary \cite{CheungXJ2,CheungEcho}.

\subsection{The complete destructive instrument}

Let \(Q\) be a finite measured register with orthonormal basis
\(\{\ket{\tau_j}\}\), and let \(A\) contain every retained quantum
register. Define
\[
 D_j=\bra{\tau_j}_Q\otimes I_A,\qquad
 \mathcal M(X)=\sum_j\proj j_R\otimes D_jXD_j^\dagger.
\]

\begin{proposition}[Record and retained residual]
\label{prop:instrument}
The map \(\mathcal M\) is completely positive and trace preserving,
and
\[
 \Tr_R\mathcal M(X)=\Tr_QX.
\]
These identities hold with every untouched finite reference.
The output contains \(A\), not a returned quantum \(Q\)-factor.
Positive-probability branches retain their actual normalized
conditional states; zero branches have no assigned normalized
state.
\end{proposition}

\begin{proof}
The recorded Kraus operators are \(\ket j_R\otimes D_j\), and
\[
 \sum_jD_j^\dagger D_j=I_{QA}.
\]
Writing
\(X=\sum_{j,k}\ket{\tau_j}\bra{\tau_k}\otimes X_{jk}\)
gives output blocks \(X_{jj}\), whose sum is \(\Tr_QX\).
The same expansion includes a reference.
\end{proof}

This is the complete I-B/En09-A instrument, restated locally
\cite{CheungIB,CheungEn09}. Consumption refers to the declared
interface, not physical annihilation of all apparatus degrees
of freedom.

\subsection{I-C A and D: a classical report is not code recovery}

For channels use half-diamond distance
\[
 d_\diamond(\Phi,\Psi)=\frac12\norm{\Phi-\Psi}_\diamond,
 \qquad
 e_{\mathrm{rec}}(\mathcal C)
 =\inf_{\mathcal D}d_\diamond(\mathcal D\mathcal C,\id).
\]

\begin{theorem}[Sharp record-only obstruction]
\label{thm:classical}
Let \(C\simeq\CC^d\), \(d\ge2\), and suppose the complete accessible
output is classical, with no input-correlated quantum side
resource available to the decoder:
\[
 \mathcal C(X)=\sum_j\Tr(F_jX)\proj j,\qquad
 F_j\ge0,\quad\sum_jF_j=I_C.
\]
For every deterministic CPTP decoder \(\mathcal D\),
\begin{equation}
 F_e(\mathcal D\mathcal C)\le\frac1d,\qquad
 d_\diamond(\mathcal D\mathcal C,\id_C)\ge1-\frac1d.
 \label{eq:classical}
\end{equation}
The distance bound is sharp over this class.
If the implemented complete channel is within half-diamond
distance \(\delta\) of such a classical channel, its optimal
recovery error is at least
\(\max\{0,1-d^{-1}-\delta\}\).
\end{theorem}

\begin{proof}
The decoded channel has measure-and-prepare form
\[
 \mathcal D\mathcal C(X)=\sum_j\Tr(F_jX)\sigma_j.
\]
Its normalized Bell-input state is separable.
For product density operators \(a,b\),
\[
 \Tr[\Phi_d(a\otimes b)]
 =\frac1d\Tr(a^{\mathsf T}b)\le\frac1d.
\]
Convexity gives the fidelity bound.
Testing the Bell projector gives the trace-distance lower bound,
and the diamond distance dominates this referenced test.

For sharpness, computational dephasing is
\[
 \Delta_d=\frac1d\sum_{k=0}^{d-1}\operatorname{Ad}_{Z^k},
 \quad Z=\sum_{j=0}^{d-1}e^{2\pi\ii j/d}\proj j.
\]
Separating its identity term gives the matching upper bound
\(1-1/d\). The approximate statement follows from contraction
under the decoder and the triangle inequality.
\end{proof}

This is the entanglement-breaking mechanism
\cite{HSR}; its operational forms are recorded in I-C,
En07-A/B/C, and the supplied Hawk05 sources
\cite{CheungIC,CheungEn07,CheungHawk05}.
Appending classical heat labels or classical postprocessing
does not change the premise.
A retained input-dependent quantum register does change it.

The distinct Haar-average pure-state infidelity bound is
\((d-1)/(d+1)\), using
\(F_{\mathrm{av}}=(dF_e+1)/(d+1)\)
\cite{NielsenAverage}. It must not replace the whole-code
diamond constant.

\subsection{I-C B and F: logs are records, not a Page history}

For a complete retained classical history,
\[
 H(Y_{\le n+1})
 =H(Y_{\le n})+H(Y_{n+1}\mid Y_{\le n})
 \ge H(Y_{\le n}).
\]
For a classical--quantum output,
\[
 S(YK)=H(Y)+\sum_y p_yS(\rho_{K\mid y}).
\]
If only the classical record survives, its accessible observable
algebra is commutative.

These are the relevant I-C B/F distinctions
\cite{CheungIC}. A particular list of SHA-256 digests is
deterministic integrity data unless a probability model for that
list is separately introduced. It is not an entropy of radiation,
and an append-only execution transcript is not a Page history
of the loop.

The finite Page means in the Hawk05 sources have fixed-input
ensemble quantifiers. Whole-code recovery has arbitrary
referenced-input quantifiers. No Page mean is used to infer
\eqref{eq:classical}, a Ward identity, or a gravitational endpoint
\cite{CheungHawk05}.

\subsection{I-C C: spectral heat is not thermodynamic heat}

The symbols \(\Ccal_b\) and \(\Dcal_b\) denote spectral
functionals. They are not Landauer costs.

If a complete non-bath system \(M\), a fixed bath Hamiltonian
\(H_B\), inverse temperature \(\beta>0\), an initially independent
Gibbs bath \(\tau_B\), and a joint unitary are separately declared,
then
\begin{equation}
 \beta Q_B
 =S_{\mathrm n}(\rho_M)-S_{\mathrm n}(\omega_M)
  +I_{\mathrm n}(M:B)_\omega
  +D_{\mathrm n}(\omega_B\Vert\tau_B).
 \label{eq:thermal}
\end{equation}
Here entropies use natural logarithms and
\(Q_B=\Tr H_B(\omega_B-\tau_B)\).
The equality follows from unitary entropy conservation and the
definition of Gibbs relative entropy
\cite{ReebWolf}. It does not require an entropy-decreasing reset
merely to be valid. A strictly positive entropy decrease of the
complete non-bath system implies positive heat because the last
two terms are nonnegative.

No such bath, Hamiltonian, or reset is declared in this loop
calculation. Thus neither the equality's physical hypotheses nor
a positive reset cost is supplied by an apex heat coefficient.

\subsection{I-C E: a destructive prefix selects no apex domain}

A finite destructive prefix does not choose between CR, F,
a positive tensor-core family, or a different boundary model.
These are operator-domain specifications. Theorem~\ref{thm:ST}
closing on both branches is not a physical branch-selection rule.

Likewise, fixed-resolution neutral-core report convergence does
not imply equality or non-injectivity of exact continuum
functionals. The I-C/I-B-ECHO comparison concerns fixed reports,
specified continuous preparations, and complete finite
instruments \cite{CheungIC,CheungEcho}.

\scopebox{Readout release boundary}{%
Only the first object in \eqref{eq:pipeline} is an analytical
target of this paper. A shredder instrument neither acquires
the missing coefficient nor selects its domain.
An echo window is not an apex residue.
No readout, entropy, or thermal comparison discharges O.1.}

%======================================================================
\section{Static debugging and artifact integrity}
\label{sec:integrity}
%======================================================================

\subsection{Corrections identified from the supplied listings}

\begingroup\small
\begin{longtable}{@{}L{.31\linewidth}L{.62\linewidth}@{}}
\toprule
Issue & Correction in this revision\\
\midrule
\endfirsthead
\toprule
Issue & Correction in this revision\\
\midrule
\endhead
\bottomrule
\endfoot

Six ``\(SO(1,1)\)'' structures &
Require reflection evenness; otherwise Levi-Civita
pseudoscalars are allowed.\\

Unique scale-covariant D01 extension &
Replace by uniqueness of the specified meromorphic
continuation. Parabolic covariance alone allows endpoint
contact terms; see Remark~\ref{rem:extension}.\\

Mellin value established by a routine name &
Replace with the explicit formula and proof
\eqref{eq:Mclosed}.\\

Mixed monomial skipped by \loc{continue} &
Quarantine the entire affected physical-branch key.
Keep a coverage status and the nonzero mixed terms.\\

UNIT mixed terms &
A genuine UNIT momentum-slice replay must retain all channel
pairs. On the released lattice, where the reported mixed terms
are absent, the algebraic UNIT closure is unchanged.\\

Missing log key interpreted as computed zero &
Require an evaluated-key/coverage ledger.
Legacy absence supports at most a comparison with a declared
zero, not proof that the key was evaluated.\\

\loc{weighted_checks.py} uses numerical
\loc{nsimplify} for an ``exact'' check &
Use exact arithmetic in a cyclotomic quotient. Rational
reconstruction from a numerical approximation is not an exact
identity proof by itself.\\

Printed diagnostics without assertions &
New acceptance checks raise exceptions on mismatch.
A printed \loc{ALL_MATCH} is not a substitute for an acceptance
predicate and process status.\\

\loc{eval} on serialized component keys &
Use literal key parsing and a restricted arithmetic-expression
parser. This is an input-integrity improvement, not a graph proof.\\

``Radial'' Ward scan &
The historical point \((X,Y)=(\texttt{rr},0.2)\) has
\(r=\sqrt{\texttt{rr}^2+0.04}\ge0.2\), so it does not approach
the apex. A true radial scan must scale both coordinates at
fixed angle. No apex asymptotic theorem is inferred from that
scan.\\

Finite resonance or random-rank scans &
Distinguish exact arithmetic on a specified finite matrix from
proof of candidate completeness and polynomial identities on
the full jet class.\\

Modular rank described as only heuristic &
Apply Lemma~\ref{lem:modular}; exactness fails here because the
complete matrix/certificate is absent, not because modular
minors can never certify rank.\\

Memory-qualified \(C_4\) normalization &
Use the explicit relation in
Remark~\ref{rem:Fursaevnorm}.\\
\end{longtable}
\endgroup

Numerical quadrature without interval or rigorous truncation
bounds remains corroboration, not a certified error enclosure.
Converting high-precision numbers to ordinary Python
\loc{float} or \loc{complex} also limits the precision actually
retained.


\subsection{v8 execution of the v7 program and response to its corrections}
\label{sec:v8response}

The listing of Appendix~\ref{app:code} was extracted byte for byte
as \loc{code/v8/cone_v7.py} and run unchanged. The results were:
\begin{itemize}
\item \loc{selftest}: return code 0. All assertions passed,
covering the cyclotomic orbifold checks for \(N=2,\dots,8\), the
catalogue algebra, the ST row, and rank 6 for the six-structure
decomposition.
\item legacy log comparison: PASS on all 18 graph--branch pairs.
Keys absent from the legacy logs are listed rather than read as
zero.
\item tree normalization: \(N=1\), \(A_{\rm GH}=1/2\),
\(\lambda_{\rm pred}=-2\ii\).
\item frozen-engine replay: 27 runs (graph \(\times\)
\{CR, F, UNIT\}), all with return code 0. The coverage
comparison passed on all 18 physical pairs, and the physical
branch values differ from \loc{logs/v4} by exactly zero.
\end{itemize}
The replay quarantines mixed keys. These are D01 \((1,1)\),
\((1,6)\), \((2,2)\), \((2,8)\), \((6,1)\), \((6,6)\), \((8,2)\),
\((8,8)\); D06 \((1,1)\), \((2,2)\); and D13 \((1,1)\), \((2,2)\),
\((6,1)\), \((8,2)\).
On the UNIT lattice, those keys gain the mixed-pair terms that the
historical \loc{continue} discarded. For example, D01 \((1,1)\)
gains \(-\ka^2q_3^2/\pi\), D06 \((1,1)\) gains
\(\ii\ka^2P^2/(2\pi)\), and D13 \((1,1)\) gains
\(\ii\ka^2q_0/(4\pi)\).
Table~\ref{tab:v8verdict} gives the point-by-point verdicts.

\begingroup\small
\begin{longtable}{@{}p{.04\linewidth}L{.30\linewidth}L{.16\linewidth}L{.42\linewidth}@{}}
\caption{Verdicts on the v7 static-debugging corrections
(executed in v8).}
\label{tab:v8verdict}\\
\toprule
\# & v7 correction & Verdict & Executed check and effect on v6 numbers\\
\midrule
\endfirsthead
\toprule
\# & v7 correction & Verdict & Executed check and effect\\
\midrule
\endhead
\bottomrule
\endfoot
1 & Six structures classify only the reflection-even sector &
wording, agreed & The rank-6 assertion passes. No number changes.\\
2 & Contact term and uniqueness of extension & wording, agreed &
No number changes.\\
3 & Mellin value should not rest on a routine name & agreed &
The closed form \eqref{eq:Mclosed} is checked exactly; the value is
unchanged.\\
4 & Mixed monomial skipped by \loc{continue} & real defect
(\loc{spectral_apex.py}, line 134) & No historical physical value
was affected, because the v4 logs hold no partial values on mixed
keys. v8 strengthens the point: the quarantine must test pair
densities, not pole totals (Section~\ref{sec:v8mixed}).\\
5 & UNIT records omit mixed terms & real defect & The UNIT records
change on 14 keys. The UNIT statements in v6 were lattice-only, so
they are unaffected.\\
6 & A missing key read as a computed zero & real defect &
The v7 program's explicit \loc{key_status} coverage was run on
every key. No number changes.\\
7 & \loc{nsimplify} in \loc{weighted_checks.py} & real defect &
The exact cyclotomic check passes. No number changes.\\
8 & Unasserted diagnostics & real defect & The v7 program asserts
them, and every assertion passes. No number changes.\\
9 & \loc{eval} on keys & real defect (\loc{st_scan},
\loc{st_test_const}, \loc{st_v6}) & A literal parser gives
identical results.\\
10 & Ward scan not at a radial point & real defect & The true
radial scan was run. The geometric residual is \(\le10^{-28}\) for
all \(r\ge0.003\). The Friedrichs residual grows from \(0.0039\)
at \(r=1\) to \(0.101\) at \(r=0.003\), whereas the historical
reading was a plateau near \(0.019\). v6 printed no Ward-scan
number.\\
11 & Finite resonance and random-rank scans & real limitation &
Replaced by the complete exact certificates of
Section~\ref{sec:v8cert}.\\
12 & Modular rank called only heuristic & addressed &
Every exact rank equals the historical modular rank, and
Lemma~\ref{lem:modular} applies.\\
13 & \(C_4\) normalization & algebra verified &
\((\alpha/2)(\alpha^{-2}-1)(\alpha^{-2}+11)/90=\Dcal_0\) holds
identically. The equation-number attribution to the source
remains from memory (\dag).\\
14 & Coarse channel marker & verified & \loc{projector_check}
passes on the frozen engine.\\
15 & Tree-hash ordering & valid in principle &
The locale is \loc{C.UTF-8}, and the shell digest is reproducible.
The frozen-tree digests are unchanged (Section~\ref{sec:v8hash}).\\
\end{longtable}
\endgroup

\noindent\textbf{v6 numbers changed: none.}
One v6 status changes: the \(hh\) branch difference becomes
conditional (Section~\ref{sec:v8mixed}).

\subsection{The coarse channel marker needs its own check}

The historical spectral code labels a tensor propagator entry
by whether its independent coordinates are longitudinal,
mixed, or transverse. That rule is not a general helicity
projector theorem.

For the four-dimensional de Donder numerator, the transverse
trace--trace block vanishes, while the transverse trace can
couple to the longitudinal scalar block. In that particular
case the coarse marker can agree with the projected
propagator. The corrected replay verifies
\[
 P_S\Pi=\Pi_S
\]
against explicit coordinate projectors before using the marker.
It does not export this four-dimensional check to unsupplied
evanescent-index arrays.

\subsection{The XVI-O1-ST sign witness}

The supplied comparison attributes to XVI-O1-ST-v5 a
cross-artifact witness: a typeset D01 sixth coefficient
\(-4/3\), a regression print \(+4/3\), and a residual
\((8/3)e_6\) identified there as a transcription discrepancy
\cite{CheungO1ST}.

The arithmetic is immediate:
\[
 \left(\frac43-\left(-\frac43\right)\right)e_6
 =\frac83e_6.
\]
Its identification as a transcription rather than a computational
error remains source-qualified: that manuscript's complete
artifacts are not supplied here.
The \(e_6\) coordinate belongs to that comparison's basis and
must not be silently identified with a present cone coordinate.

Selecting the source's preferred sign does not rerun a loop,
clear every predecessor sign audit, or update XVI-O.
A matching cone catalogue is not a universal cross-artifact
sign certificate.

\subsection{Historical hashes and new manifests}

The following are \emph{reported v6 labels}, not freshly
recomputed hashes:
\begin{center}
\scriptsize
\begin{tabular}{@{}ll@{}}
\toprule
Tree & Reported SHA-256\\
\midrule
XVI-P &
\texttt{ec965c15ae25cd3bb9cc2b3a134e966b}\\
&\texttt{19a1803c78d1fe5bf7554da2a8b3aaa6}\\
XVI-Q &
\texttt{0df9a0d10409d31ac0c2d0a47197e3f6}\\
&\texttt{e03c7ae258ec5fb775137f66cb0493c4}\\
XVI-R &
\texttt{57187bc6825f5c6ae88633b7279c6c0b}\\
&\texttt{e57f7c4918fdbc525da4aaf958331edb}\\
\bottomrule
\end{tabular}
\end{center}

The original shell tree hash depends on its filename ordering,
exclusions, and command semantics. The new Python manifest
routine records relative names and byte hashes without
\loc{find|xargs} filename ambiguity. Its tree digest uses a
different explicit serialization and is \emph{not} expected to
equal the historical shell digest.

A hash binds supplied bytes to an expected digest. It does not
prove a theorem, establish authorship, authenticate an
unavailable file, or show that a log was generated by a claimed
execution. Before/after equality also requires both actual
manifests, not merely two printed matching strings.


\subsection{v8: recomputed frozen-tree digests}
\label{sec:v8hash}

The historical shell serialization was recomputed in v8
(locale \loc{C.UTF-8}) on the read-only trees, both before
any v8 run and after the last run:
\begin{center}
\scriptsize
\begin{tabular}{@{}lll@{}}
\toprule
Tree & When & SHA-256 (shell serialization)\\
\midrule
XVI-P & before/after & \texttt{ec965c15ae25cd3bb9cc2b3a134e966b}\\
&&\texttt{19a1803c78d1fe5bf7554da2a8b3aaa6}\\

XVI-Q & before/after & \texttt{0df9a0d10409d31ac0c2d0a47197e3f6}\\
&&\texttt{e03c7ae258ec5fb775137f66cb0493c4}\\

XVI-R & before/after & \texttt{57187bc6825f5c6ae88633b7279c6c0b}\\
&&\texttt{e57f7c4918fdbc525da4aaf958331edb}\\
\bottomrule
\end{tabular}
\end{center}
They equal the reported v6 labels above, before and after
(\loc{logs/v8/tree_hashes_v8.json}). The v7 Python manifest
serialization gives different digests, as v7 predicted:
XVI-P \texttt{acd88f27}\ldots\texttt{5b3c7} (1209 files), XVI-Q \texttt{c13b12d1}\ldots\texttt{37586} (879 files), XVI-R \texttt{d8b9ed80}\ldots\texttt{f4015} (457 files).
No XVI-P, XVI-Q or XVI-R file was edited.

\subsection{Computational scope of the embedded program}

Appendix~\ref{app:code} consolidates the parts of the supplied
scripts needed by this paper:
\begin{itemize}
\item heat-moment algebra and exact orbifold checks;
\item declared catalogue construction;
\item invariant decomposition, ghost equation, and ST identity;
\item coverage-aware comparison with historical or new logs;
\item optional frozen-engine momentum-slice replay and tree
normalization;
\item byte-level manifests.
\end{itemize}

The older exploratory jet/cohomology, kernel, scratch, and
document-assembly scripts remain historical material in
\cite{CheungConeScripts}; they are not copied as redundant
submission machinery or claimed as completed certificates.
The consolidated program is a replacement verification
implementation, not a verbatim archive of every v1--v6 file.
Its inclusion does not change their historical hashes.



%======================================================================
\section{v8: an exact complete span certificate}
\label{sec:v8cert}
%======================================================================

\subsection{What is certified}

The apex functionals of the historical inventory
(\loc{apex_jets.py}) are bilinear in finitely many jet atoms and
polynomial in the external momentum. v8 ports them mechanically to
exact Gaussian-rational arithmetic (\loc{code/v8/apex_exact.py};
the port reproduces every historical modular count) and evaluates
them on \emph{all} relevant unit atom pairs over a complete
\(6\times6\) momentum grid, enough to determine every polynomial of
the occurring degree. Therefore the evaluated columns are complete,
and the row spaces computed are the exact row spaces.

\begin{theorem}[Exact span certificate; executed]
\label{thm:v8cert}
For each \((\alpha,\text{branch})\) in Table~\ref{tab:v8cert}, over
\(\mathbb Q(\ii)\) and for the functionals as implemented in the
declared inventory:
the gauge-invariant candidate ranks, the ST rank, \(\dim Z\),
\(\dim B\), \(B\subset Z\) and \(\dim H^0=\dim Z-\dim B\) are as
displayed. Each rank is certified by a nonzero minor (hash and
indices in \loc{apex_cert_*.json}) and each kernel by verification
on every complete column.
\end{theorem}

\begin{center}
\scriptsize
\begin{longtable}{@{}lllllllllll@{}}
\caption{Complete exact certificates (v8).}\label{tab:v8cert}\\
\toprule
\((\alpha,\text{br})\) & \(N\) & ST rows & ST rank & \(\dim Z\) & \(B\subset Z\) & \(\dim B\) & \(\dim H^0\) & Ward & on \(K\) & s\\
\midrule
\((3/10,\CR)\) & 34=18+8+6+2 & 8640 & 26 & 8 & yes & 6 & 2 & 26 & 6 & 292\\
\((3/10,\Fr)\) & 34=18+8+6+2 & 8640 & 26 & 8 & yes & 6 & 2 & 26 & 6 & 302\\
\((7071/10^4,\CR)\) & 34=18+8+6+2 & 8640 & 26 & 8 & yes & 6 & 2 & 26 & 6 & 316\\
\((7071/10^4,\Fr)\) & 38=21+8+7+2 & 14400 & 29 & 9 & yes & 6 & 3 & 29 & 6 & 554\\
\bottomrule
\end{longtable}
\end{center}

Here ``Ward'' is the number \(N-\dim Z\) of Ward-constrained
combinations; at the 34-candidate points this is the v6 count
\(26\), now exact. The ghost-number \(+1\) cohomology vanishes at
the sampled exact point \((7071/10^4,\CR)\) (26/26 images closed;
\(H^1=0\)).

\subsection{The count 34 is angle and branch dependent}

The historical modular inventory, which the exact port reproduces
wherever both exist, gives:
\begin{center}
\scriptsize
\begin{tabular}{@{}lllllll@{}}
\toprule
\(\alpha\) & branch & \(N\) & \(\dim Z\) & \(\dim B\) & \(\dim H^0\) & \(N-\dim Z\)\\
\midrule
3/10 & CR & 34 & 8 & 6 & 2 & 26\\
3/10 & F & 34 & 8 & 6 & 2 & 26\\
41/100 & CR & 34 & 8 & 6 & 2 & 26\\
41/100 & F & 34 & 8 & 6 & 2 & 26\\
1/2 & CR & 42 & 9 & 7 & 2 & 33\\
1/2 & F & 42 & 9 & 7 & 2 & 33\\
309/500 & CR & 34 & 8 & 6 & 2 & 26\\
309/500 & F & 38 & 9 & 6 & 3 & 29\\
2/3 & CR & 35 & 8 & 6 & 2 & 27\\
2/3 & F & 39 & 9 & 6 & 3 & 30\\
7071/10000 & CR & 34 & 8 & 6 & 2 & 26\\
7071/10000 & F & 38 & 9 & 6 & 3 & 29\\
4/5 & CR & 34 & 8 & 6 & 2 & 26\\
4/5 & F & 38 & 9 & 6 & 3 & 29\\
8501/10000 & CR & 34 & 8 & 6 & 2 & 26\\
8501/10000 & F & 38 & 9 & 6 & 3 & 29\\
\bottomrule
\end{tabular}
\end{center}
Thus \(N=34\), \(\dim Z=8\), \(N-\dim Z=26\) hold on CR at every
tested non-resonant angle and on F for \(\alpha<1/2\); F at
\(\alpha>1/2\) gives \(38\) with \(\dim H^0=3\) (an additional
F-branch singular exponent enters), and \(\alpha=1/2\), \(2/3\) are
resonant. The ``34-dimensional basis'' and ``26 Ward combinations''
of v6/v7 are therefore qualified statements; only the rows of
Table~\ref{tab:v8cert} are exact.

\subsection{Scope of the certificate}

The certificate concerns the declared inventory and its functionals,
at the listed rational angles. It does not prove a statement for
general \(\alpha\) (two points do not certify a rational function),
does not identify the inventory with the complete interacting basis,
and does not evaluate any graph. Within that scope it is
\CPIR.

%======================================================================
\section{v8: Ward combinations, the 34- and 18-dimensional maps, and \texorpdfstring{\(\phi\phi\)}{phi phi}}
\label{sec:v8map}
%======================================================================

\subsection{Constant-profile support of the Ward combinations}

Let \(K\) be the set of constant-profile coordinates, meaning jets
with \(m=0\). Let \(\pi_K\) be the coordinate projection onto
\(K\), and let \(\mathrm{Ann}_K\) be the annihilator of \(\pi_K Z\)
inside the span of the \(K\)-supported functionals. In every exact
run we found \(\abs K=12\), with HH~6, CB~2, CK~3 and PP~1. The
ST functionals supported on \(K\) span a space of dimension
\(6\). The \(m=0\) rows alone already span \(\mathrm{Ann}_K\).
Further, \(\rank\pi_KZ=6\) and \(\rank\pi_KB=5\). Also
\(\dim(Z\cap\ker\pi_K)=2\), except at \((7071/10^4,\Fr)\), where it
is \(3\).
The \(K\)-supported Ward combinations split by block support as
\[
 \underbrace{2}_{hh\text{ only}}+\underbrace{2}_{hh+\mathrm{CK}}
 +\underbrace{2}_{\mathrm{CB}+\mathrm{CK}}=6.
\]
No combination is supported on CB, CK or PP alone.
This decomposition is identical at every computed angle and
branch.

\subsection{The \texorpdfstring{\(hh\)}{hh} dictionary and the two exact Ward checks}

The engine's constant-profile \(hh\) coordinates are related to
the Wick-rotated longitudinal block by
\[
 x_{00}=-h_{\tau\tau},\quad x_{03}=\pm\ii\,h_{\tau z},\quad
 x_{33}=h_{zz},\quad f=\tfrac12(h_{11}+h_{22}),\qquad
 q_0=\pm\ii q_\tau,\ q_3=q_z.
\]
Both sign choices give a solvable linear system, and the two
give the same map. With
\(\Gamma_{\rm sym}=4\Gamma\), the six engine coordinates
\(y\) are
\[
 y=\bigl(-b_1,\,-2b_3,\,-b_2,\,b_6,\,b_7,\,-\tfrac14 b_8\bigr)
\]
in units of \(\ka^2/\pi\), where \(b_j\) are the coefficients of the
structures \eqref{eq:structures}
(from \(E=MS\) with \(4b=M^{\mathsf T}y\), recorded in
\loc{hh_map_*.json}).
The two \(hh\)-only Ward combinations are
\[
 W_1=b_1+2b_3,\qquad W_2=b_2 .
\]

\begin{proposition}[Exact \(hh\) Ward checks; executed]
\label{prop:v8hhward}
In exact arithmetic, at \(\alpha=3/10\) and at
\(\alpha=7071/10000\) and on both branches, the declared \(hh\) catalogue sum
satisfies \(W_1=W_2=0\).
The individual graphs do not. At \(\alpha=3/10\), for example,
\[
 (W_1,W_2)^{\rm D01}=\Bigl(-\tfrac{91}{5760},\ \tfrac{91}{2880}\Bigr)
 =-(W_1,W_2)^{\rm D03},
\]
and D02, D11 and D12 give zero.
\end{proposition}

\noindent
So the check is not trivially satisfied: it tests a cancellation
between D01 and D03. Like every catalogue statement, it is
conditional on Assumption~\ref{ass:reduction}.

\subsection{The remaining Ward combinations}

The two \(hh\)+CK and two CB+CK combinations contain the v7
ST row and the projected ghost equation restricted to \(m=0\)
(Theorems~\ref{thm:ghost} and \ref{thm:ST}). This identification
is \CONDITIONAL, because v8 builds no CK/CB coordinate dictionary.
The \(26-6=20\) combinations that involve \(m>0\) jets have not
been evaluated. They are \OPEN, not zero.

\subsection{The 18-dimensional engine map}

The HH sector of the inventory has rank 18 on the 34-candidate
branches, and rank 21 at \((7071/10^4,\Fr)\). The dictionary above
fixes the six constant-profile coordinates exactly, on all four
computed branches. The other twelve involve transverse
derivatives or non-axisymmetric jets, and their convention map is
\OPEN. The map is therefore \PARTIAL: 6 of 18 coordinates.

\subsection{The six exact values}

Write the non-PP part of a closed class as \(\sum_{i=1}^6e_iB_i\),
using the exact basis of \(B\) stored in
\loc{apex_cert_*_vectors.json}.
The \(hh\) constant-profile data give
\(\rank(\pi_{K\cap hh}B)=4\), and the resulting system is consistent:
\begin{center}
\small
\begin{tabular}{@{}lllll@{}}
\toprule
\((\alpha,\text{branch})\) & \(e_1\) & \(e_3\) & \(e_2+2e_5\) & \(e_4+e_5\)\\
\midrule
\((3/10,\CR)\) & \(7/40\) & \(0\) & \(7/120\) & \(-131/1920\)\\
\((3/10,\Fr)\) & \(7/40\) & \(0\) & \(7/120\) & \(63/640\)\\
\((7071/10^4,\CR)\) & \(2929/40000\) & \(0\) &
\(-32842877/282840000\) & \(-78439319/4525440000\)\\
\((7071/10^4,\Fr)\) & \(2929/40000\) & \(0\) & \(-32842877/282840000\) & \(-136599773/1508480000\)\\
\bottomrule
\end{tabular}
\end{center}
At both angles, \(e_1=(1-\alpha)/4\). This is observed at two
points, not proved for general \(\alpha\).
The CR/F difference appears only in \(e_4+e_5\), where it equals
\(\nu/8\) exactly: \(-1/6\) at \(\alpha=3/10\) and \(2071/28284\) at
\(\alpha=7071/10^4\), with \(\nu=4142/7071\).
This matches the \(hh\)-sum difference
\(b_8^{\CR}-b_8^{\Fr}=-\nu/4\), which is consistent with
\eqref{eq:branchdifference}.
The coefficients \(e_5\) and \(e_6\) are not determined. They need
CK/CB constant data, which in turn need a dictionary, and
\(m>0\) data, because \(\dim(Z\cap\ker\pi_K)=2\) (and \(3\) at
\((7071/10^4,\Fr)\)).
The status is \PARTIAL, and only in the stated basis.

\subsection{The \texorpdfstring{\(\phi\phi\)}{phi phi} values}

The first \(\phi\phi\) value is the catalogue PP coefficient
\(-c_2\ka^2P^2/(4\pi)\) on each branch. Its branch difference is
\(-\nu\ka^2P^2/(4\pi)\), and the engine coordinate is
\(y_{PP}=\ka^2c_2/(4\pi)\).
For the second value, the D05a part is \(0\) in the reduced model
of Section~\ref{sec:missing}.
The D04 part is still \OPEN, for the following reasons.
\begin{itemize}
\item At constant profile, the replay of D04 has a single key
\((0,0)\). It has no hidden mixed pair and no evanescent
coefficient, and its UNIT value is \(-\ka^2P^2/(4\pi)\).
\item The quadratic-profile part needs the \(r^2\)-weighted
mixed-channel trace. v8 proves only the ingredients for it
(Lemma~\ref{lem:v8duhamel}); the mode sum and the \(r^2\)
assembly were not done.
\end{itemize}

%======================================================================
\section{v8: mixed channels, hidden pairs, and a Duhamel lemma}
\label{sec:v8mixed}
%======================================================================

\subsection{Feynman densities of channel pairs}

For a bubble whose two lines carry channels \(s\) and \(t\), let
\(x\) be the Schwinger fraction of the \(t\)-line. The identity
\[
 \frac1{l^2(l-p)^2}=\int_0^1\frac{\dd x}{\bigl((l-xp)^2+x(1-x)P^2\bigr)^2}
\]
turns the pole of the slice numerator into a polynomial density
\(\rho_{(s,t)}(x)\). By construction,
\(\int_0^1\rho_{(s,t)}\,\dd x\) equals the circle pole of that pair.
This identity was checked exactly for every bubble pair of every
graph (\loc{slice_D*.json}, field \loc{rho_integrates_to_circle}).
The cone value of the pair is
\[
 \int_0^1\rho_{(s,t)}(x)\,C_{s,t}(x)\,\dd x,\qquad
 C_{s,t}(x)=\ApTr\bigl[e^{-(1-x)A_s}e^{-xA_t}\bigr].
\]
For \(s=t\) this is the circle pole times \(\ApTr e^{-A_s}\).

\begin{definition}[Hidden mixed pair]
A pair with \(s\ne t\) is \emph{hidden} if its circle pole vanishes
but \(\rho_{(s,t)}\not\equiv0\).
\end{definition}

A selection rule based on totals, \(\text{circle}\ne0\), never sees
a hidden pair. The historical lattice records value each pair as
its circle pole times one channel coefficient, so a hidden pair
receives the value zero.

\subsection{The complete list (executed)}

The slice run \loc{logs/v8/slice/hidden_mixed.json} found hidden
pairs in exactly two graphs.
\begin{itemize}
\item \textbf{D01, 13 keys.} Eight of them, \((1,1)\), \((1,6)\),
\((2,2)\), \((2,8)\), \((6,1)\), \((6,6)\), \((8,2)\), \((8,8)\),
were already quarantined; their pairs \((0,1)\) and \((1,0)\) have
\(\rho\propto(2x-1)\).
The other five are \textbf{new}: \((4,4)\), \((4,7)\), \((5,5)\),
\((7,4)\), \((7,7)\). Their pairs \((0,2)\) and \((2,0)\) have
\[
 \rho^{(4,4)}_{(0,2)}=-\frac{\ii(q_0^2-q_3^2)(6x^2-10x+3)}{8\pi},\qquad
 \rho^{(4,4)}_{(2,0)}=-\frac{\ii(q_0^2-q_3^2)(6x^2-2x-1)}{8\pi}.
\]
The \((5,5)\) densities are \(4\) times these, \((7,7)\) equals
\((4,4)\), and \((4,7)\), \((7,4)\) carry the opposite sign.
\item \textbf{D03, 8 keys.} These are the vector-type keys,
already excluded.
\end{itemize}
D02, D04, D05a, D06, D11, D12 and D13 have no hidden pairs.

\begin{proposition}[Status of the v6 branch difference; executed and proved]
\label{prop:v8hidden}
The new D01 keys enter the \(hh\) sum, because the indices 4, 5
and 7 are the transverse components 11, 12 and 22.
Since \(\rho_{(2,0)}(1-x)=\rho_{(0,2)}(x)\) and
\(C_{2,0}(x)=C_{0,2}(1-x)\) by cyclicity, their total cone value is
\[
 2\int_0^1\rho_{(0,2)}(x)\,C_{0,2}(x)\,\dd x
 =-2\int_0^1R(x)\,C_{0,2}'(x)\,\dd x .
\]
Here \(R(x)=\int_0^x\rho_{(0,2)}\), which is proportional to
\(x(x-1)(2x-3)\).
This total is generically nonzero, and on the CR and F branches it
is built from different channel spectra.
Therefore the printed v6 value
\(-\nu\ka^2P^2f_0f'_0/(4\pi)\) is unchanged as the lattice
evaluation of the declared catalogue. As a statement about the
cone graph, however, it is \CONDITIONAL{} on the vanishing, or the
separate evaluation, of this mixed-channel term.
\end{proposition}

\begin{proof}
The symmetry of the densities is read off the executed list. The
cyclicity follows from the cyclicity of the regulated apex trace.
Integration by parts uses \(R(0)=R(1)=0\), which holds because
\(\int_0^1\rho_{(0,2)}=0\).
\end{proof}

\subsection{A per-mode Duhamel identity}

\begin{lemma}[Duhamel identity for Bessel channels]
\label{lem:v8duhamel}
In one angular mode, let
\(A_\mu=-\partial_r^2-r^{-1}\partial_r+\mu^2r^{-2}\) on its
Friedrichs domain, with \(\mu,\nu>0\). With the bulk-subtracted
mode trace,
\[
 (\mu^2-\nu^2)\int_0^1\Tr\bigl[e^{-xA_\nu}\,r^{-2}\,e^{-(1-x)A_\mu}\bigr]\dd x
 =\frac{\mu-\nu}2 ,
\]
and \(C_{\mu,\nu}'(x)=(\mu^2-\nu^2)\Tr[e^{-(1-x)A_\mu}r^{-2}e^{-xA_\nu}]\).
In particular \(C(0)=\Tr e^{-A_\mu}\) and \(C(1)=\Tr e^{-A_\nu}\)
after the same subtraction.
\end{lemma}

\begin{proof}
Differentiate \(x\mapsto e^{-(1-x)A_\mu}e^{-xA_\nu}\) and use
\(A_\mu-A_\nu=(\mu^2-\nu^2)r^{-2}\). Integrate over \(x\in[0,1]\),
and use the per-mode value \(-\mu/2\) (up to a \(\mu\)-independent
constant) of the subtracted trace of \(e^{-A_\mu}\).
\end{proof}

\noindent\emph{Numerical corroboration.} We used Richardson
extrapolation in the radial cutoff. For \(\mu-\nu=1\) the result
was \(0.49993\) against \(0.5\), and for \(\mu-\nu=0.85\) it was
\(0.84990\) against \(0.85\)
(\loc{duhamel_mode_check2.json}). This is corroboration, not an
enclosure.

In the Cartesian frame, all channels share the angular phase set,
and the Bessel orders are \(\abs{k/\alpha+S}\). The CR branch flips
the \(k=-1\) mode of \(S=\pm2\). Lemma~\ref{lem:v8duhamel} reduces
\(C_{0,2}'\) to a mode sum of \(r^{-2}\)-sandwiched traces.
v8 does \emph{not} evaluate that mode sum, and does not do the
\(r^2\)-weighted D04 assembly. Both remain \OPEN.

%======================================================================
\section{v8: evanescent coefficients under a declared prescription}
\label{sec:v8evan}
%======================================================================

The frozen \(\tilde g\)-leg arrays are still not supplied. v8
therefore declares a prescription and computes under it.

\begin{definition}[Slice-DRED prescription]
On the zero-transverse-momentum slice, the UV degree \(-2\) angular
average uses \(\hat g_L=\eta_L-\tilde G\). The symbols
\(\tilde G_{00},\tilde G_{03},\tilde G_{33}\) are formal, and
\(d=2\) at pole level. The average is
\(\langle P\rangle=(l^2)^j\hat\Delta^jP/(2^jj!\prod_i(d+2i))\).
The \emph{evanescent coefficients} of a key are the coefficients of
the \(\tilde G\)-dependent part of its pole.
\end{definition}

Two checks were run for every pair of every graph. The
\(\tilde G=0\) part reproduces the circle pole of the replay, and
the bubble densities integrate to it.
After combining with \(c_s\) and excluding quarantined and
evanescent-mixed keys, the counts are as follows
(\loc{evanescent_combined.json}).
\begin{center}
\small
\begin{tabular}{@{}lccccccccc|c@{}}
\toprule
 & D01 & D02 & D03 & D04 & D05a & D06 & D11 & D12 & D13 & total\\
\midrule
CR & 21 & 0 & 26 & 0 & 0 & 4 & 25 & 0 & 10 & 86\\
F & 21 & 0 & 26 & 0 & 0 & 4 & 25 & 0 & 10 & 86\\
UNIT & 34 & 0 & 34 & 0 & 0 & 6 & 25 & 0 & 14 & 113\\
\bottomrule
\end{tabular}
\end{center}
Beyond the earlier quarantine, D01 also excludes its five
evanescent-mixed keys. On the CR branch, D03, D06, D11 and D13
carry an overall factor \((\alpha-1)\). For example,
\[
 \text{D13}\,(4,0):\quad
 \frac{\ii\ka^2(\alpha^2-1)}{192\pi\alpha}
 \bigl(\tilde G_{00}q_0-2\tilde G_{03}q_3+\tilde G_{33}q_0\bigr),
\]
so the coefficient is \(-c_0/16\) times the bracket.
\textbf{Status:} \CONDITIONAL{} on the declared prescription. These
values are not identified with the frozen engine's \(\tilde g\)
export. Agreement at \(\tilde G=0\) is a consistency check, not a
proof of identity preservation.

%======================================================================

%======================================================================
\section{v9: what is new, and the response to the priority list}
\label{sec:whatv9}
%======================================================================

v9 is built mechanically from the read-only v8 source
(\loc{code/build_v9.py}). Every v8 statement is kept unless this
section or Table~\ref{tab:open} says it has changed. v9 has new
evidence class marks: \emph{Executed in v9} means a script under
\loc{code/v9/}, a log under \loc{logs/v9/}, and a SHA-256 in
Appendix~\ref{app:v9artifacts}. The frozen trees XVI-P, XVI-Q and
XVI-R were hashed before and after the v9 runs. Both digests equal
the v8 values (Section~\ref{sec:v9hash}). Files v1--v8 were not
edited.

\begingroup\small
\begin{longtable}{@{}L{.05\linewidth}L{.27\linewidth}L{.20\linewidth}L{.40\linewidth}@{}}
\caption{v9 verdicts on the nine priority items.}\label{tab:v9verdict}\\
\toprule
\# & Item & Verdict & Evidence / obstruction\\
\midrule
\endfirsthead
\toprule
\# & Item & Verdict & Evidence\\
\midrule
\endhead
1 & Hidden mixed D01 terms & computed; v6 \emph{restored} &
Proposition~\ref{prop:v9TT}: pure TT, vanishes on the lattice.
\(\Delta J(\alpha)\) at 14 angles (Table~\ref{tab:v9dJ}).
IR log: Proposition~\ref{prop:v9log}.\\
2 & \(r^2\)-weighted mixed sum, D04 part of 2nd \(\phi\phi\) &
\OPEN, precise obstruction & Section~\ref{sec:v9D04}: the jet densities
are not exported, and same-\(k\) pairing is not IR-summable. No
closed form is claimed.\\
3 & CK/CB dictionary; \(e_5,e_6\); closed forms &
dictionary computed; \(e_5,e_6\) \emph{excluded with proof} at constant profile &
Theorem~\ref{thm:v9ghost}: all 6 constant-profile Ward functionals
hold; \(\ker\pi_KB\) is one-dimensional.\\
4 & 20 jet-level Ward combinations, full projection &
\OPEN, precise obstruction & Section~\ref{sec:v9jet}.\\
5 & 18-dim engine map beyond 6 & \OPEN, precise obstruction &
Section~\ref{sec:v9jet}.\\
6 & Evanescent vs frozen \(\tilde g\) & \emph{excluded}, precise obstruction &
Proposition~\ref{prop:v9evan}.\\
7 & Vector-type keys & log parts computed; finite parts \OPEN &
Proposition~\ref{prop:v9vector}.\\
8 & Generic-\(\alpha\) rank certificate &
computed on chambers; exceptional set \PARTIAL &
Theorem~\ref{thm:v9generic}.\\
9 & New sheets & read; one newer sheet & Section~\ref{sec:v9sheets};
\loc{new_oct18g/} written.\\
\bottomrule
\end{longtable}
\endgroup

\paragraph{Changes to v8 results.}
(a) The v8 status ``\CONDITIONAL{} on the hidden D01 term'' for the v6
\(hh\) branch difference is lifted. The value is unchanged and is
now proved on the lattice (Proposition~\ref{prop:v9TT}).
(b) The four CK/CB-supported Ward combinations, which v8
identified only conditionally, are now checked exactly on the
catalogue. This uses two fitted, \(\alpha\)-independent ghost
normalizations (Theorem~\ref{thm:v9ghost}).
(c) The v8 open item ``\(e_5\) and \(e_6\)'' is resolved negatively:
the constant-profile data determine only \(e_6-e_2\) and
\(e_5+e_6/2\) beyond the v8 combinations.
(d) The incomplete-beta reduction used in an internal v8/v9 draft
module for the F-branch \(k=-1\) mode (\loc{I_F}) contained an
exponent error (\(x^{1-\mu}\) for \(x^{\mu-1}\)). It was not used
in any printed v8 number. It was removed, and every v9 value uses
the closed-form quadrature \loc{I_gen}.
No other v8 number changes.

%======================================================================
\section{v9: the hidden D01 pairs}
\label{sec:v9hidden}
%======================================================================

\subsection{Closed form of the mixed trace}

Let \(A_\mu\) be the radial Bessel operator of order \(\mu\ge0\) on
the cone channel, and let
\[
 g(x;\mu,\nu)=\Tr\bigl[e^{-(1-x)A_\mu}\,r^{-2}\,e^{-xA_\nu}\bigr],
 \qquad 0<x<1,
\]
with \(d=\mu-\nu\) and \(h=(\mu+\nu)/2\).

\begin{proposition}[Mixed trace; derived and executed]
\label{prop:v9g}
\[
 g(x;\mu,\nu)=\frac{\Gamma(h)\Gamma(h+1)}{4\,\Gamma(\mu+1)\Gamma(\nu+1)}\;
 {}_2F_1\!\bigl(\tfrac d2+1,\tfrac d2;\mu+1;x\bigr)\;
 {}_2F_1\!\bigl(1-\tfrac d2,-\tfrac d2;\nu+1;1-x\bigr).
\]
In particular \(g(x;\mu,\mu)=1/(4\mu)\), and
\(\int_0^1g\,\dd x=1/(2(\mu+\nu))\).
\end{proposition}

\begin{proof}[Derivation]
Expand both Bessel kernels in their power series and integrate the
radial variables term by term. This gives
\(x^\mu(1-x)^\nu\) times
\(F_4(h,h+1;\mu+1,\nu+1;x^2,(1-x)^2)\). Since
\(h+(h+1)-(\mu+1)+1=\nu+1\), Bailey's reduction
\(F_4(a,b;c,a+b-c+1;X(1-Y),Y(1-X))={}_2F_1(a,b;c;X)\,{}_2F_1(a,b;a+b-c+1;Y)\)
applies with \(X=x\) and \(Y=1-x\). Euler's transformation of each
factor removes the prefactor \(x^\mu(1-x)^\nu\) and gives the
stated parameters. The integral identity is an independent check
of the v8 Duhamel lemma, Lemma~\ref{lem:v8duhamel}.
\loc{code/v9/g_formula_check.py} confirms both the closed form and
the integral against direct quadrature of the kernels. The
residual, \(4.7\times10^{-3}\), is the truncation tail of the
quadrature at \(R=60\), not a discrepancy.
\end{proof}

\subsection{The hidden pairs are pure TT}

With \(\hat\rho=6x^2-10x+3\) and
\(\hat R=x(1-x)(3-2x)\), we have \(\hat R'=\hat\rho\). For a mode
pair \((\mu,\nu)\), integration by parts gives
\[
 I(\mu,\nu)=\int_0^1\hat\rho\,f
 =-(\mu^2-\nu^2)\int_0^1\hat R\,g .
\]
\loc{code/v9/modefun.py} evaluates this in two independent ways
for \(\nu=\mu\pm2\): an exact Gamma series summed by
Euler--Maclaurin, and direct quadrature of
Proposition~\ref{prop:v9g}. The two agree to 30 digits.

\begin{proposition}[TT structure; proved]
\label{prop:v9TT}
On the D01 keys \((4,4),(4,7),(7,4),(7,7),(5,5)\), the hidden
\((0,2)\)/\((2,0)\) pairs contribute
\[
 -\frac{\kappa^2P^2J}{8\pi}\Bigl[(h_{11}-h_{22})(h'_{11}-h'_{22})
 +4h_{12}h'_{12}\Bigr],
\]
with \(J=\sum_{\rm modes}I\) in the stated mode list. On the v6 lattice
(\(h_{11}=h_{22}=f\), no \(h_{12}\)) the bracket is identically zero.
Hence the v6 \(hh\) branch difference
\(-\nu\kappa^2P^2f_0f_0'/(4\pi)\) holds without the v8 condition.
\end{proposition}

\begin{proof}
The key-by-key coefficients come from
\loc{logs/v8/slice/hidden_mixed.json} and the catalogue
normalization
\((P^2\kappa^2/16\pi)/(\ii(q_0^2-q_3^2)/8\pi)\). Key \((4,4)\)
receives \(-\kappa^2P^2J/(8\pi)\); keys \((4,7)\) and \((7,4)\) carry
the opposite sign, and key \((5,5)\) carries four times key
\((4,4)\). Collecting terms gives the bracket. The lattice
projection of \loc{invariant_data} uses
\(\{0{:}0,3{:}1,9{:}2,4{:}3,7{:}3\}\). It sets \(h_{11}=h_{22}\)
and drops index~5, so it maps the bracket to zero for every value
of \(J\).
\end{proof}

\subsection{Infrared structure and the finite branch difference}

\begin{proposition}[IR logarithm; executed asymptotics]
\label{prop:v9log}
As \(\mu\to\infty\),
\(I(\mu,\mu+2)=\tfrac13+\tfrac1{30\mu}-\tfrac1{10\mu^2}+\cdots\) and
\(I(\mu,\mu-2)=-\tfrac13+\tfrac1{30\mu}+\tfrac1{10\mu^2}+\cdots\)
(coefficients up to \(\mu^{-8}\) are in \loc{logs/v9/asym02.json}).
Under same-\(k\) pairing, the constant terms cancel between \(k\)
and \(-k\), but the \(1/\mu\) terms add to \(\alpha/(15k)\).
Therefore \(J\) diverges like \((\alpha/15)\log\Lambda\). In
general, a pair \((s,t)\) contributes
\(\tfrac{\alpha}{2}(s-t)^2\,x(x-1)\log\Lambda\) to \(C_{s,t}(x)\).
This is the frame/flux mismatch of a long-range \(r^{-2}L\)
perturbation.
\end{proposition}

So the absolute TT values are not pure apex coefficients. Their
finite part depends on the scheme (for example \(k\)-\(\zeta\)).
The CR and F branches differ only in the \(S=2\), \(k=-1\) mode,
whose order is \(-\nu_0\) (CR) or \(+\nu_0\) (F), with
\(\nu_0=2-1/\alpha\). The difference
\[
 \Delta J(\alpha)=I(1/\alpha,\,1/\alpha-2)-I(1/\alpha,\,2-1/\alpha)
\]
is therefore finite and does not depend on the regularization.
The corresponding catalogue difference is
\(-\kappa^2P^2\Delta J/(8\pi)\) times the TT bracket, and it is
invisible on the lattice.

\begin{center}
\scriptsize
\begin{longtable}{@{}lllll@{}}
\caption{The finite TT branch difference (executed, \loc{code/v9/hidden02_run.py}; CR mode by series and by quadrature, agreement \(<10^{-30}\)).}\label{tab:v9dJ}\\
\toprule
\(\alpha\) & \(\nu_0\) & \(I_{\rm CR}\) (\(k=-1\)) & \(I_{\rm F}\) (\(k=-1\)) & \(\Delta J\)\\
\midrule
\(7/20\) & \(-0.857143\) & \(-0.301772223719917\) & \(-0.316747429419663\) & \(0.014975205699746\)\\
\(2/5\) & \(-0.500000\) & \(-0.292350993934024\) & \(-0.321408664644675\) & \(0.029057670710650\)\\
\(9/20\) & \(-0.222222\) & \(-0.281195058597902\) & \(-0.300043868930418\) & \(0.018848810332516\)\\
\(1/2\) & \(0.000000\) & \(-0.268135533878013\) & \(-0.268135533878013\) & \(0.000000000000000\)\\
\(11/20\) & \(0.181818\) & \(-0.253008191060374\) & \(-0.233034968697703\) & \(-0.019973222362671\)\\
\(3/5\) & \(0.333333\) & \(-0.235654644431579\) & \(-0.198192074766226\) & \(-0.037462569665354\)\\
\(2/3\) & \(0.500000\) & \(-0.208792354336879\) & \(-0.154666895890071\) & \(-0.054125458446808\)\\
\(7071/10000\) & \(0.585773\) & \(-0.190296473359310\) & \(-0.130398179653525\) & \(-0.059898293705785\)\\
\(3/4\) & \(0.666667\) & \(-0.168755307135675\) & \(-0.106489500975910\) & \(-0.062265806159765\)\\
\(4/5\) & \(0.750000\) & \(-0.141051936421474\) & \(-0.080941692454300\) & \(-0.060110243967174\)\\
\(17/20\) & \(0.823529\) & \(-0.110436566041365\) & \(-0.057725471830814\) & \(-0.052711094210551\)\\
\(9/10\) & \(0.888889\) & \(-0.076793924011713\) & \(-0.036635399844723\) & \(-0.040158524166990\)\\
\(19/20\) & \(0.947368\) & \(-0.040015650826980\) & \(-0.017460622324418\) & \(-0.022555028502562\)\\
\(99/100\) & \(0.989899\) & \(-0.008266790265495\) & \(-0.003364124456341\) & \(-0.004902665809154\)\\
\bottomrule
\end{longtable}
\end{center}
At \(\alpha=1/2\), \(\nu_0=0\) and \(\Delta J=0\) exactly. The sign
of \(\Delta J\) is the sign of \(-\nu_0\) at every computed angle.
As \(\alpha\to1\), the CR mode degenerates (\(\mu+\nu\to0\)); the
values decrease toward zero (\(\Delta J(99/100)\approx-4.9\times10^{-3}\)).
PSLQ was not attempted. We claim no elementary
closed form, since each value is an integral of a product of two
\({}_2F_1\) factors.

%======================================================================
\section{v9: the CB/CK dictionary and the six constant-profile Ward functionals}
\label{sec:v9ghost}
%======================================================================

\subsection{Dictionaries}

We use the same method as the \(hh\) map (Section~\ref{sec:v8map}).
Unit \(m=0\) atoms \(L=(1,0),(0,1)\) for \(c,\bar c\), and
\(LL\)/\(PM\) atoms for \(K\), are evaluated by \loc{eval_sym}
(reflection only, so \(\Gamma_{\rm sym}=2\Gamma\)) on the complete
\(6\times6\) momentum grid. A dictionary is accepted only if the
catalogue bilinear (D06 for \(\langle\bar c\,c\rangle\), D13 for
\(\langle K c\rangle\)) lies \emph{exactly} in the span of the
evaluated \(K\)-candidates on all columns
(\loc{code/v9/cbck_map.py}, \loc{code/v9/ck_free.py}).

\begin{proposition}[CB and CK dictionaries; executed, exact]
\label{prop:v9dict}
With \(q_0=\pm\ii q_\tau\), \(q_3=q_z\):
\begin{enumerate}[label=(\roman*)]
\item CB: exactly the four sign choices with opposite \(\ii\)-signs
for \(\bar c_0\) and \(c_0\) are accepted. All four give
\(y_{\rm CB}=\ii(-R/8,\,0)\,\kappa^2/\pi\), with \(R=3z_0+4z_1\).
\item CK: with the \(K\) components left free, the solution is
unique: \(K^{00}=-LL_{\tau\tau}\), \(K^{03}=\mp2\ii LL_{\tau z}\), and
\(y_{\rm CK}=\ii\bigl((1-\alpha)/16,\,R/8,\,0\bigr)\kappa^2/\pi\).
The \(PM\) scale is undetermined because the catalogue's PM keys
are computed zeros.
\end{enumerate}
Both hold at \((3/10,\CR)\), \((3/10,\Fr)\), \((7071/10^4,\CR)\) and
\((7071/10^4,\Fr)\). Here
\(R=(17\alpha-7)(\alpha-1)/(12\alpha)\).
\end{proposition}

\subsection{The six functionals}

The annihilator of \(\pi_KZ\) is printed in
\loc{apex_cert_*_vectors.json} and is identical at all four
points. In the order
\(K=(0,1,2,5,6,12,18,19,26,27,29,32)\) it consists of
\(W_1=y_0+y_1\), \(W_2=y_2\), \(W_3=-2y_0+y_{19}\),
\(W_4=-2\ii y_0+4\ii y_5+4\ii y_6+y_{26}\),
\(W_5=-2\ii y_{18}+y_{27}\), and \(W_6=2\ii y_0+y_{29}\). (v8
quoted only the real parts of these vectors.)

\begin{theorem}[Ghost-sector Ward check; executed, exact]
\label{thm:v9ghost}
Put \(\lambda_{\rm CB}=-4\ii\) and \(\lambda_{\rm CK}=-8\), fixed
once and independent of \(\alpha\) and of the branch. Rescale the
CB and CK coordinates by these factors. Then all six \(W_i\)
vanish on the catalogue at all four certified points. Solving
\(\sum_ie_i\pi_KB_i=y_K\) on all twelve coordinates gives
\[
 e_3=0,\quad e_1=\tfrac{1-\alpha}4,\quad
 e_6-e_2=-\tfrac R2,\quad
 e_5+\tfrac{e_6}2=\frac{\alpha^2-1}{48\alpha},
\]
\[
 e_4-\tfrac{e_6}2=\frac{119\alpha^2-120\alpha+25}{192\alpha}\ (\CR),\qquad
 e_4-\tfrac{e_6}2=\frac{7(17\alpha-7)(\alpha-1)}{192\alpha}=\tfrac{7R}{16}\ (\Fr).
\]
The kernel of \(\pi_KB\) is exactly
\(\Span\{(0,1,0,\tfrac12,-\tfrac12,1)\}\). Therefore \(e_5\) and
\(e_6\) are not determined separately by constant-profile data.
\end{theorem}

\begin{proof}
\loc{code/v9/e_solve.py} carries out the linear algebra exactly; its output is
\loc{logs/v9/e_solve.json}. The checks have different strengths.
\(W_3\) and \(W_6\) hold because both of their entries vanish.
\(W_5\) defines \(\lambda_{\rm CB}\), since both of its terms are
proportional to \(R\).
\(W_4\) is a genuine check. Its \(hh\) part is
\(4\ii(b_6+b_7)=\ii(c_0-c_1)=\ii(1-\alpha)/2\), coming from
D01--D03 and D11--D12. Its CK part is
\(\lambda_{\rm CK}\,\ii(1-\alpha)/16\), coming from D13. The two
have the same \(\alpha\)-shape, so one constant
\(\lambda_{\rm CK}=-8\) works at every \(\alpha\) and on both
branches. The \(e\)-values follow from the twelve equations. The
closed forms are the rational identities displayed above,
checked at the four points.
\end{proof}

\noindent\emph{Scope.} The normalizations \(\lambda\) are fitted,
not derived from the catalogue's BRST conventions. The closed
forms hold on the certified points. For general generic \(\alpha\)
they follow from the closed forms of \(z_0,z_1,z_2\) together with
the observed \(\alpha\)-independence of \(\pi_KB\) and of the
annihilator, which was verified at four points but not proved.
Status: \PARTIAL{} for general \(\alpha\), exact on the four points.
The v8 identity \(e_1=(1-\alpha)/4\) is now obtained twice, from
the \(hh\) block and from the CK coordinate~26.

%======================================================================
\section{v9: chambers and a generic-\texorpdfstring{\(\alpha\)}{alpha} rank certificate}
\label{sec:v9chamber}
%======================================================================

The structural key of a jet \((C,n,m)\) is
\((C,n,\epsilon,j,\operatorname{sgn}(n/\alpha+s))\), where
\(m=\epsilon\nu+2j\). Two angles lie in the same \emph{chamber} when
every gh0 (\(T=4\)) and gh\(-1\) (\(T=3\)) sector has the same set
of structural candidates. In a chamber, every ST row entry is a
fixed rational function of \(\alpha\).

\begin{proposition}[Walls; executed and proved]
\label{prop:v9walls}
In \((1/4,1)\) the walls are exactly \(\{1/3,1/2,2/3\}\), on both
branches. CR has the same structural signature on all four
chambers. F has one signature on \((1/4,1/3)\cup(1/3,1/2)\) and a
different one, with \(N=38\) at the certified points, on
\((1/2,2/3)\cup(2/3,1)\).
\end{proposition}

\begin{proof}
A non-generic candidate requires \(\pm2n/\alpha\in\mathbb Z\) for
some admissible pair \((n,-n)\), or a sign or \(\epsilon\) change.
The jet budget gives \(\nu\le6\) and \(n\le4\). So a wall has the
form \(\alpha=2n/M\) with \(M\le12\), or \(\alpha=n/|s|\). Every
such rational has denominator at most 12.
\loc{code/v9/chamber_scan.py} evaluates the structural signature
at every rational with denominator at most 40 in \((1/4,1)\),
which is exhaustive by this bound. It finds only the three walls
listed (\loc{logs/v9/chamber_cache2.json}).
\end{proof}

\begin{theorem}[Generic rank over \(\mathbb Q(\alpha)\)]
\label{thm:v9generic}
On every CR chamber in \((1/4,1)\), and on the F chambers in
\((1/4,1/2)\), the ST rank over \(\mathbb Q(\alpha)\) is exactly
\(26\), with \(\dim Z=8\) generically. On the F chambers in
\((1/2,1)\), the generic rank \(r\) satisfies \(29\le r\le30\).
\end{theorem}

\begin{proof}
\emph{Lower bound.} The rank is lower semicontinuous under
specialization within a chamber. So it suffices to have one point
per chamber with a nonzero minor of the required size. Exact
\(\mathbb Q(\ii)\) minors exist at \(3/10\) and \(7071/10^4\)
(Theorem~\ref{thm:v8cert}). A nonzero minor mod \(p\) also implies
a nonzero rational minor, and the modular runs at \(41/100\),
\(309/500\), \(4/5\), \(8501/10^4\) and at the scan points of
Table~\ref{tab:v9scan} provide such minors.
\emph{Upper bound.} The identity \(\sigma_0\sigma_{-1}=0\) holds
identically, so \(B\subset Z\). We have \(s\phi=0\), so \(PP\subset Z\).
\(B\) has no PP components, and \(\dim B\ge6\) generically, by lower
semicontinuity from the certified points. Hence \(\dim Z\ge8\),
which gives \(r\le N-8\): \(26\) for \(N=34\) and \(30\) for
\(N=38\). The ninth cocycle seen at the F points with \(N=38\) is
not shown to persist generically, so \(29\le r\le30\) there.
\end{proof}

\begin{center}
\scriptsize
\begin{longtable}{@{}lll@{}}
\caption{Modular scan (\loc{code/v9/apex_scan.py}, an unchanged copy of the historical driver with a new output directory; prime \(10^9+9\), randomized complete-rank sampling). Every rational \(p/q\in(1/4,1)\) with \(q\le20\).}\label{tab:v9scan}\\
\toprule
branch & result & angles\\
\midrule
CR & \((1/4,1/3)\): \(N=34\), rank \(26\), \(\dim Z=8\), \(\dim B=6\), \(\dim H^0=2\) & 10 angles\\
CR & \((1/3,1/2)\): \(N=34\), rank \(26\), \(\dim Z=8\), \(\dim B=6\), \(\dim H^0=2\) & 21 angles\\
CR & \((1/2,2/3)\): \(N=34\), rank \(26\), \(\dim Z=8\), \(\dim B=6\), \(\dim H^0=2\) & 21 angles\\
CR & \((2/3,1)\): \(N=34\), rank \(26\), \(\dim Z=8\), \(\dim B=6\), \(\dim H^0=2\) & 41 angles\\
CR & wall \(1/3\): \(N=35\), rank \(27\), \(\dim Z=8\), \(\dim H^0=2\) & wall\\
CR & wall \(1/2\): \(N=42\), rank \(33\), \(\dim Z=9\), \(\dim H^0=2\) & wall\\
CR & wall \(2/3\): \(N=35\), rank \(27\), \(\dim Z=8\), \(\dim H^0=2\) & wall\\
F & \((1/4,1/3)\): \(N=34\), rank \(26\), \(\dim Z=8\), \(\dim B=6\), \(\dim H^0=2\) & 10 angles\\
F & \((1/3,1/2)\): \(N=34\), rank \(26\), \(\dim Z=8\), \(\dim B=6\), \(\dim H^0=2\) & 21 angles\\
F & \((1/2,2/3)\): \(N=38\), rank \(29\), \(\dim Z=9\), \(\dim B=6\), \(\dim H^0=3\) & 21 angles\\
F & \((2/3,1)\): \(N=38\), rank \(29\), \(\dim Z=9\), \(\dim B=6\), \(\dim H^0=3\) & 41 angles\\
F & wall \(1/3\): \(N=35\), rank \(27\), \(\dim Z=8\), \(\dim H^0=2\) & wall\\
F & wall \(1/2\): \(N=42\), rank \(33\), \(\dim Z=9\), \(\dim H^0=2\) & wall\\
F & wall \(2/3\): \(N=39\), rank \(30\), \(\dim Z=9\), \(\dim H^0=3\) & wall\\
\multicolumn{3}{@{}l}{Completed: 192 of 192 runs (an incomplete scan is reported as such, never extrapolated).}\\
\bottomrule
\end{longtable}
\end{center}

\noindent\emph{Exceptional set.} Inside a chamber, the rank drops
exactly on the common zero set of the maximal minors. That set is
finite, but the minor polynomial was not computed: it would need a
symbolic-\(\alpha\) elimination of unknown degree. So the
exceptional set is \(E=\{1/3,1/2,2/3\}\cup E_{\rm minor}\). The
scan shows that \(E_{\rm minor}\) contains no rational with
denominator at most 20, but \(E_{\rm minor}\) itself is \OPEN.

%======================================================================
\section{v9: precise obstructions}
\label{sec:v9obs}
%======================================================================

\subsection{The D04 part of the second \texorpdfstring{\(\phi\phi\)}{phi phi} value}
\label{sec:v9D04}

The D04 internal lines are a graviton, with channel weights
\(\mathrm{pw}\), and \(\phi\) (\(S=0\)). At constant profile only
\(z_0\) contributes, and the same-channel part is the
\(x\)-independent \(D\)-moment
\(\Tr[r^2e^{-A}]\). The second \(\phi\phi\) value needs the
\(r^2\)-weighted mixed traces
\(\Tr[e^{-(1-x)A_\mu}r^2e^{-xA_\nu}]\) with transverse-momentum
numerators and the \(H_I\) intertwining. There are two
obstructions:
\begin{enumerate}[label=(\alph*)]
\item the zero-transverse-momentum slice does not export these
jet-level densities, so the integrand is not available as data;
\item the frame/flux mismatch of Proposition~\ref{prop:v9log} carries
over with an extra \(r^2\). Same-\(k\) pairing is then not
IR-summable. The mismatch gives the log of the \(r^{-2}\) case
there, and power counting gives a power divergence with the
\(r^2\) weight (\emph{argued, not executed}).
\end{enumerate}
So a rational reconstruction in \(\alpha\) has nothing
well-defined to reconstruct until a scheme for the mixed
same-\(k\) sum is declared. Status: \OPEN; no closed form is
claimed, and no checks at \(\alpha=1\) or \(1/2\) are claimed.

\subsection{Jet-level Ward combinations and the engine map}
\label{sec:v9jet}

Twenty of the 26 Ward combinations, and twelve of the eighteen
\(hh\) engine coordinates, involve \(m>0\) jets (\(r^2\) profiles)
or transverse \(P\)/\(M\) components. The dictionary method of
Sections~\ref{sec:v8map} and \ref{sec:v9ghost} would need the
catalogue's values on \(m=2\) and weight-\(\pm1\) atoms. The
declared catalogue has no such values: it is the zero-transverse,
constant-profile slice. They would come from the
\(r^2\)/\(D_b\)-weighted traces, which are obstructed as in (b)
above. The full 34-dimensional projection is therefore available
only on the twelve \(K\) coordinates, where all six functionals now
hold (Theorem~\ref{thm:v9ghost}). That projection has rank 6, and
its angle and branch dependence enters only through \(z_0,z_1,z_2\).
Status: \OPEN.

\subsection{Evanescent coefficients versus the frozen \texorpdfstring{\(\tilde g\)}{g-tilde} rule}

\begin{proposition}[Excluded, with obstruction; executed]
\label{prop:v9evan}
The frozen engine (\loc{XVI-P/code/p5_entry/dred.py},
\loc{eta_trace_Gt}, read-only) admits exactly one evanescent
pattern, \(c\sum_r\eta^{rr}\tilde G_{rr}\to c(4-d)\), and declares
every other pattern open. Of the 285 v8 slice-DRED keys
(\loc{logs/v8/slice/evanescent_combined.json}), 164 contain
\(\tilde G\) monomials of degree 2 or 3. All of them contain only
the longitudinal \(\tilde G_{00},\tilde G_{03},\tilde G_{33}\), and
only 31 have a linear part of partial-trace form
\(c(\tilde G_{00}-\tilde G_{33})\)
(\loc{code/v9/evan_match.py}). The v8 coefficients therefore lie
outside the class the frozen rule can evaluate. A match would need
the transverse \(\tilde G_{11},\tilde G_{22}\) components, which
the slice does not export, and a rule for non-trace monomials.
\end{proposition}

\subsection{Vector-type keys}

\begin{proposition}[IR logs of the hidden keys; executed]
\label{prop:v9vector}
Let \(L_{\rm key}=\sum_{(s,t)}(s-t)^2\int_0^1\rho_{(s,t)}x(x-1)\)
over D01 and D03 (\loc{code/v9/vector_logs.py}), and put
\(c=\ii P^2/(30\pi)\). Then the vector keys \((1,6),(2,8),(6,1),(8,2)\)
have \(L=0\). The keys \((1,1)\) and \((2,2)\) have \(L=+c\), and
\((6,6)\) and \((8,8)\) have \(L=-c\). The TT keys give
\(-c,+c,+c,-c,-4c\) on \((4,4),(4,7),(7,4),(7,7),(5,5)\).
\end{proposition}

The vector-key log is
\(c\,P^2\sum_{a=1,2}\eta^{\mu\nu}h_{\mu a}h'_{\nu a}\) up to key
normalization. It does not cancel within D01\(+\)D03, so the
vector keys inherit the scheme dependence of
Proposition~\ref{prop:v9log}. Whether profile-dependent Ward
partners remove it cannot be decided from the slice. The finite
parts of the vector keys stay \OPEN{} and are excluded, not zero.

%======================================================================
\section{v9: frozen-tree digests and the new registry sheets}
\label{sec:v9hash}
\label{sec:v9sheets}
%======================================================================

The tree digests of XVI-P, XVI-Q and XVI-R, in the shell
serialization of Section~\ref{sec:v8hash}, were recomputed before
the first v9 run and after the last one:
\begin{center}\scriptsize
\begin{tabular}{@{}lll@{}}
\toprule
Tree & before / after & SHA-256\\
\midrule
XVI-P & equal & \texttt{ec965c15ae25cd3bb9cc2b3a134e966b}\\ && \texttt{19a1803c78d1fe5bf7554da2a8b3aaa6}\\
XVI-Q & equal & \texttt{0df9a0d10409d31ac0c2d0a47197e3f6}\\ && \texttt{e03c7ae258ec5fb775137f66cb0493c4}\\
XVI-R & equal & \texttt{57187bc6825f5c6ae88633b7279c6c0b}\\ && \texttt{e57f7c4918fdbc525da4aaf958331edb}\\
\bottomrule
\end{tabular}
\end{center}
They are equal to the v8 values. The v8 source and PDF were
hashed before and after as well
(\loc{logs/v9/v8_hashes_before.sha}, \loc{logs/v9/v8_hashes_after.sha}).

\paragraph{Registry.} The registry was fetched again on 18 October
2026. It lists 258 identifiers, one more than the 257 in
\loc{new_oct18f/}: the new identifier is \loc{I-C-Ex2}. It was
stored under \loc{new_oct18g/} with SHA-256 sums. The six sheets
of \loc{new_oct18f/} and \loc{I-C-Ex2} were read as cited inputs
(abstracts and all theorem statements; \loc{logs/v9/sheets_extract.txt}):
\begin{itemize}
\item \emph{I-B-G-S-B} \cite{CheungIBGSBv9} is the only sheet that
bears on a v9 input. For a declared class of nonnegative tensor
connection forms, it proves Mosco and strong-resolvent selection of
the componentwise Friedrichs realization, and states that for
neutral spin~2 and \(\alpha>1/2\) the limit excludes the
curvature-regular branch. This bears on the CR/F branch input: it
favours F \emph{within that declared core class}. It does not
derive the graviton operator, and the sheet itself stresses that a
smooth metric and its holonomy do not fix the perturbation
operator. So both branches remain inputs here, and the CR results
are kept as conditional on CR.
\item \emph{I-B-ECHO-v2} \cite{CheungEchoV2} gives the scalar
apex-diffracted Green operator and finite report instruments. It
is not an apex-residue input; compare Example~\ref{ex:image}.
\item \emph{X-J-ST-2} \cite{CheungXJST2} gives the free photon
reference (Maxwell--Friedrichs with a Friedrichs scalar ghost) and
charged scalar coefficients. It is a free, source-test reference
and imports no interacting graviton Ward identity.
\item \emph{Hawk05-ST, I-C-v2, I-C-Ex-v2, I-C-Ex2}
\cite{CheungHawk05ST,CheungICv2,CheungICExv2,CheungICEx2} are
finite access and recovery bounds (\(1-1/d\), \(1-k/d\)). They
belong to the instrument side of the boundary of
Section~\ref{sec:readout} and are not used for O.1.
\end{itemize}

%======================================================================
\section{What v10 adds, and what it does not}
\label{sec:whatv10}
%======================================================================

v10 keeps every v9 statement \cite{CheungConeV9}. The v9 statements that v10 amends are
listed in Section~\ref{sec:v10changes}. The new items, with their
status, are:

\begin{center}
\small
\begin{longtable}{@{}L{.30\linewidth}L{.62\linewidth}@{}}
\caption{v10 verdicts (all executed; hashes in Appendix~\ref{app:v10artifacts}).}\label{tab:v10verdict}\\
\toprule
Item & Status\\
\midrule
\(\Delta J\) closed form & CR mode \emph{derived}. F mode \emph{identified} by PSLQ and verified at 9 fresh angles (\(\le1.5\times10^{-52}\)), not derived. Section~\ref{sec:v10dJ}.\\
\(\lambda_{\rm CB},\lambda_{\rm CK}\) & \emph{Derived} from tree operators and the v9 dictionaries (Theorem~\ref{thm:v10lambda}).\\
Exceptional points inside chambers & \emph{None}: rank 26 on all \(N=34\) chambers, 29 on F above \(\tfrac12\). Conditional on the a priori degree bound (Theorem~\ref{thm:v10chamber}).\\
Ninth cocycle on F above \(\tfrac12\) & \emph{Resolved}. It is present at every angle; it is not a constant vector, while the other eight are.\\
Transverse-jet catalogue & \PARTIAL. Exact slice jets for all 9 graphs. Same-channel \(r^2\) part dressed in two declared schemes; mixed and weight \(\ne0\) parts \OPEN{} (Section~\ref{sec:v10jet}).\\
D04 part of second \(\phi\phi\) value & \PARTIAL. The same-channel part is exactly 0. The full jet term is mixed \((1,0)\), with constant density \(\kappa^2/4\pi\). Per-mode traces are computed; the mode sum is not executed.\\
20 jet Ward combinations; 34-/18-dim maps; \(e_5,e_6\) & \OPEN. Precise obstruction in Section~\ref{sec:v10conseq}. No residual is claimed, since none was evaluated.\\
Vector-key finite parts & Computed in a declared \(k\)-\(\zeta\) scheme at 4 angles (Table~\ref{tab:v10vec}).\\
Registry & Unchanged (258 ids); no new sheet directory written.\\
\bottomrule
\end{longtable}
\end{center}

%======================================================================
\section{v10: closed form of the finite TT branch difference}
\label{sec:v10dJ}
%======================================================================

\begin{proposition}[CR mode; derived]
\label{prop:v10ICR}
Let \(\mu=1/\alpha\) and \(\beta=\mu-1\). Then
\[
 I(\mu,\mu-2)=\frac{(\mu-1)(6\mu^2+3\mu-2)}6-(\mu-1)^3(\mu+1)\,\psi_1(\mu-1).
\]
The same expression gives every \(k\le-2\) mode with \(\mu=m/\alpha\).
\end{proposition}

\begin{proof}[Derivation]
In the Gamma series of \loc{code/v9/modefun.py}, use
\((n+1)!\,\Gamma(\mu+1)/\Gamma(\mu+1+n)=\mu(\mu-1)B(n+2,\beta)\).
Each sum becomes a log-Beta integral. Continue it analytically,
keeping the Laurent data at the coincident poles. With
\(T=\psi_1(\beta)\), the intermediate sums are
\(G_2=(T\beta-1)/\beta\),
\(G_3=(-T\mu^2+3T\mu-2T+\mu-\tfrac32)/\beta\) and
\(G_4=(6T\mu^3-36T\mu^2+66T\mu-36T-6\mu^2+27\mu-25)/(12\beta)\).
Collecting them gives the formula
(\loc{code/v10/cr_mode_closed.py}). At \(\alpha=3/5\),
\(10000/7071\), \(4/7\) and \(99/100\) it matches the 34-digit
series to \(\le6\times10^{-35}\).
\end{proof}

\begin{proposition}[F mode; identified, not derived]
\label{prop:v10IF}
\[
 I(\mu,2-\mu)=\frac{(\mu-1)(59\mu^3-57\mu^2-62\mu+48)}{36}-(\mu-1)^3(\mu+1)\frac{\pi^2}6 .
\]
\end{proposition}
PSLQ \cite{PSLQ} on the basis \(\{I,1,\pi^2\}\) found integer relations at
\(\mu\in\{6/5,5/4,4/3,7/5,3/2,8/5,5/3\}\). The values were computed
to 55 digits by direct quadrature of Proposition~\ref{prop:v9g}
(\loc{code/v10/f_mode_pslq.py}). The relations interpolate to the
displayed polynomial. That polynomial was then checked at \(7/4\)
and \(9/5\), which were not used in the fit.
This is an identification, and its status is \emph{verified, not
derived}.

\begin{corollary}[\(\Delta J\) in closed form]
\label{cor:v10dJ}
\[
 \Delta J(\alpha)=\frac{(\mu-1)(-59\mu^3+93\mu^2+80\mu-60)}{36}
 -(\mu-1)^3(\mu+1)\Bigl[\psi_1(\mu-1)-\frac{\pi^2}6\Bigr],\qquad \mu=\tfrac1\alpha .
\]
It vanishes at \(\alpha=\tfrac12\) (\(\psi_1(1)=\pi^2/6\), and the
polynomial has a zero at \(\mu=2\)). As \(\alpha\to1\),
\(\Delta J=-\tfrac12(\mu-1)+O((\mu-1)^2)\), consistent with the v9
value \(\Delta J(99/100)\approx-4.9\times10^{-3}\).
\end{corollary}

\begin{center}\scriptsize
\begin{longtable}{@{}llll@{}}
\caption{Fresh-angle verification (\loc{code/v10/dJ_closed_verify.py}; quadrature at 60 digits, and 80 digits for a precision check).}\label{tab:v10dJ}\\
\toprule
\(\mu=1/\alpha\) & \(|I_{\rm F}^{\rm quad}-I_{\rm F}^{\rm closed}|\) & \(|I_{\rm CR}^{\rm quad}-I_{\rm CR}^{\rm closed}|\) & \(\Delta J\) (closed)\\
\midrule
\(\sqrt2\) & \(1.0\times10^{-56}\) & \(2.5\times10^{-56}\) & \(-0.0598989753454146726\)\\
\(10000/7071\) & \(5.1\times10^{-57}\) & \(1.5\times10^{-56}\) & \(-0.0598982937057849765\)\\
\(137/100\) & \(2.3\times10^{-56}\) & \(2.8\times10^{-56}\) & \(-0.0616430256621234476\)\\
\(193/100\) & \(6.1\times10^{-56}\) & \(2.0\times10^{-56}\) & \(-0.0073607981128415757\)\\
\(\varphi\_{\rm g}\) & \(2.5\times10^{-57}\) & \(3.6\times10^{-56}\) & \(-0.0427973775860188253\)\\
\(\pi/2\) & \(1.0\times10^{-56}\) & \(7.6\times10^{-57}\) & \(-0.0476710589083126482\)\\
\(5/2 (\alpha=2/5)\) & \(1.6\times10^{-55}\) & \(1.5\times10^{-52}\) & \(0.02905767071065043833\)\\
\(29/10 (\alpha=10/29)\) & \(5.1\times10^{-56}\) & \(5.6\times10^{-53}\) & \(0.01105442640041502692\)\\
\(101/100\) & \(8.0\times10^{-59}\) & \(3.2\times10^{-58}\) & \(-0.0048550798217175446\)\\
\bottomrule
\end{longtable}
\end{center}
The 14 v9 values of Table~\ref{tab:v9dJ} agree with the closed form
to \(4.4\times10^{-27}\). The v9 sentence ``we claim no elementary
closed form'' is therefore superseded for the CR mode and qualified
for the F mode.

%======================================================================
\section{v10: the ghost normalizations from tree data}
\label{sec:v10lambda}
%======================================================================

\begin{theorem}[Derivation of \(\lambda_{\rm CK}\) and \(\lambda_{\rm CB}\)]
\label{thm:v10lambda}
Use the dictionaries of Sections~\ref{sec:v8map} and
\ref{sec:v9ghost}, with signs \(s_i=1\), \(s_c=-1\), \(s_{\bar c}=1\);
the \(K\)-dictionary \(K^{00}=-\mathrm{LL}_{\tau\tau}\),
\(K^{03}=-2\ii\,\mathrm{LL}_{\tau z}\),
\(K^{33}=\mathrm{LL}_{zz}\); and \(q_0=\ii q_\tau\), \(q_3=q_z\).
Then the following three tree-level operator identities hold on all
\(m=0\) atoms and all nine test momenta \((q_\tau,q_z)\in\{1,2,3\}^2\):
\begin{align*}
 \text{engine }K_2c&=\rho\,\delta_ch,& \delta_ch_{mn}&=p_mc_n+p_nc_m, & \rho&=-\ii,\\
 \text{engine }E(-q)h&=\epsilon\,(2-\delta_{mn})E^{mn}_{\rm cat}(h), && & \epsilon&=1,\\
 \text{engine }C_2^{\mathsf T}(-q)b&=\tau\,Y_b, & Y_b&=-2\eta F_L^{\mathsf T}b, & \tau&=-\ii/2 .
\end{align*}
With the eval normalizations of the engine (\(\Gamma_{\rm sym}=4\Gamma\)
on HH and \(2\Gamma\) on CB and CK), the \(hc\) row becomes
\(8\rho B+2\lambda_{\rm CK}\epsilon G=0\). Theorem~\ref{thm:ST}
gives \(B+\lambda_{\rm pred}G=0\), with
\(\lambda_{\rm pred}=N/(\ii A_{\rm GH})=-2\ii\)
(Input~\ref{input:tree}). Hence
\[
 \lambda_{\rm CK}=\frac{4\rho\lambda_{\rm pred}}{\epsilon}=-8 .
\]
The \(b c\) row becomes
\(2\lambda_{\rm CB}\Gamma_{\rm D06}+2\lambda_{\rm CK}\tau\Gamma_{\rm D06}=0\)
by the ghost identity
\(\hat B_{\rm D06}=-2\eta F_L\hat\Gamma_{\rm D13}\)
(Theorem~\ref{thm:ghost}). Hence
\[
 \lambda_{\rm CB}=-\tau\lambda_{\rm CK}=-4\ii .
\]
\end{theorem}

\begin{proof}
\loc{code/v10/lambda_derive.py} evaluates the engine operators of
\loc{code/v8/apex_exact.py} (\(K_2\) at \(q\); \(E\) and
\(C_2^{\mathsf T}\) at \(-q\), as in the row construction of
\loc{apex_certificate.py}). Each ratio is single-valued on every
nonzero component, with no mismatch
(\loc{logs/v10/lambda_derive.json}). The two row reductions are
linear algebra on these ratios.
\end{proof}

\begin{remark}
The derivation uses no loop value. It does use the dictionary
signs, which v9 selected by exact solvability on loop data. The
values \(-8\) and \(-4\ii\) are therefore no longer free fitted
parameters, but they are not independent of the v9 sign choice.
The relation \(\lambda_{\rm CK}=\lambda_{\rm pred}\lambda_{\rm CB}\)
noted in v9 is explained by the numerical coincidence
\(4\rho/\epsilon=-4\ii=\lambda_{\rm CB}\), i.e.\ \(-\tau\lambda_{\rm CK}=4\rho/\epsilon\) with
\(\lambda_{\rm CK}=(4\rho/\epsilon)\lambda_{\rm pred}\).
\end{remark}

%======================================================================
\section{v10: chamber ranks without exceptional points}
\label{sec:v10chamber}
%======================================================================

\loc{code/v10/minor_poly.py} picks, at one angle per chamber, three
nonsingular maximal minors with \(\alpha\)-free row tags. The atom
key is \((C,n,j,i)\) with \(j=(m-\varepsilon|a|)/2\). The program
then re-evaluates the minors at 11--15 angles per chamber
(\loc{logs/v10/minor/<tag>/}).

\begin{theorem}[Rank inside chambers]
\label{thm:v10chamber}
Assume the a priori bound (D): every ST row entry is a polynomial of
degree \(\le2\) in \(1/\alpha\) on each chamber. This is argued from
the construction: \(E\) carries two derivative factors, \(K_2\) one,
and each is linear in \(1/\alpha\). It is not separately
machine-checked. Then:
\begin{enumerate}[label=(\alph*)]
\item On CR \((\tfrac14,\tfrac13),(\tfrac13,\tfrac12),(\tfrac12,\tfrac23),(\tfrac23,1)\)
and on F \((\tfrac14,\tfrac13),(\tfrac13,\tfrac12)\) (\(N=34\)), all
three minors are the same constants at every sampled angle,
\(124515522497539473408\,\ii\),
\(-73477495818488684950134624092160000\) and
\(-4059104145057262382496562859212800\), and their entries are
\(\alpha\)-independent. The rank is therefore 26 at \emph{every}
\(\alpha\) in these chambers.
\item On F \((\tfrac12,\tfrac23)\) and \((\tfrac23,1)\) (\(N=38\)),
minor~2 equals \(-400569487999071945641108176896000000000\,\ii\) at
all 28 sampled angles. Its only \(\alpha\)-dependent entries lie in
one column, so it is affine in that column and constant by (D). Hence
the rank is \(\ge29\) everywhere. Minors 0 and 1 vary, with 13 and
15 distinct values, and are never zero.
\item On F above \(\tfrac12\), \(\dim Z(\alpha)=9\) at 16 angles
(\loc{code/v10/cocycle9.py}). Every entry fits an affine function of
\(1/\alpha\) on 8 points per chamber. Write the 9 Cramer kernel
vectors of minor~2 at \(t=1/\alpha\). The residual \(W(t)\) of the
remaining rows has degree \(\le6\) by (D) and vanishes at 8 points,
so it vanishes identically. Hence the generic rank is 29, and the
rank equals 29 at every \(\alpha\): at most 29 by semicontinuity,
at least 29 by (b).
\item The intersection \(\bigcap_\alpha Z(\alpha)\) over the 8
angles of each chamber is 8-dimensional. Eight cocycles are constant
vectors. The ninth exists at every angle, but its direction depends
on \(\alpha\) and is rational in \(\alpha\) by Cramer's rule.
\end{enumerate}
\end{theorem}
This replaces the v9 range ``between 29 and 30'' by exactly 29,
conditional on (D). It also answers the v9 open item ``exceptional
set inside chambers'': the set is empty, conditional on (D).

%======================================================================
\section{v10: the transverse-jet extension of the slice catalogue}
\label{sec:v10jet}
%======================================================================

\paragraph{Construction.}
\loc{code/v10/jet_replay.py} executes the numerator section of the
frozen v7/v8 replay with one change. The external transverse momenta
\(p_1,p_2\) are kept as symbols instead of being set to zero. The
bubble denominator gains \(-p_\perp^2\), and the loop momentum stays
on the longitudinal slice (\(l_\perp=0\)). This is a declared
coverage limitation, and the status remains
\textsf{ALGEBRAIC\_SLICE\_ONLY}. The UV pole is then an exact
polynomial in \(q_0,q_3,p_\pm\). It is stored per key and channel
pair, by transverse degree and helicity weight
(\loc{logs/v10/jet/jet_<G>.json}). No frozen file is edited.

\paragraph{Checks.}
(i) At order \(p_\perp^0\), \(\sum_s\mathrm{pole}^{(0)}_{ss}z_s\)
reproduces the v8 catalogue on every computed key: 684 key--branch
checks with 0 mismatches. (ii) In D04 the UNIT sum is
\(-\kappa^2(q_0^2-q_3^2-p_\perp^2)/4\pi\), which is flat
covariance restored by the jet terms.

\paragraph{Census.}
Summed over the 9 graphs, the jet entries are:
\begin{itemize}
\item same-channel \((2,0)\): 128;
\item mixed \((2,0)\): 84;
\item same \((2,\pm2)\): 152 each;
\item mixed \((2,\pm2)\): 70 each;
\item same \((1,\pm1)\): 252 each;
\item mixed \((1,\pm1)\): 206 each.
\end{itemize}
D12 has none.

\begin{definition}[Declared jet dressing and schemes]
A same-channel, weight-0 coefficient of \(p_+p_-\) on channel \(s\)
is dressed by the \(r^2\) moment \(d_s\). In scheme A (the Abel
subtraction of Theorem~\ref{thm:D}),
\(d_0=\Dcal_0\), \(d_1=\Dcal_\alpha\),
\(d_2^{\rm F}=\Dcal_{2\alpha-1}\) and
\(d_2^{\rm CR}=\Dcal_{2\alpha}\). In scheme Z, each mode's
\(\Tr[r^{2w}e^{-sA_\mu}]\) is continued from \(\Re w<-\tfrac12\) to
\(w=1\), and the modes are summed by Hurwitz zeta. This mirrors the
role of \(\Ccal_b\) at constant profile. It is a reduction contract
in the sense of Assumption~\ref{ass:reduction}, not a theorem.
\end{definition}

\begin{proposition}[Scheme Z; executed]
\label{prop:v10Z}
Per mode, \(\Tr[r^2e^{-sA_\mu}]^Z=s\,\mu(\mu^2-1)/3\), from the
Mellin closed form
\((2s)^w\Gamma(w+1+\mu)\Gamma(-\tfrac12-w)/(2^{w+2}\sqrt\pi\,\Gamma(\mu-w))\).
This form matches quadrature for \(w<-\tfrac12\) to \(10^{-17}\) or
better. Its only poles are at half-integer \(w\), so \(w=1\) is
regular. The mode sum gives
\[
 \Dcal^Z_b=\Dcal_b+\frac{11\alpha}{180}\quad\text{(exactly, all }b),
\]
so the scheme dependence is a bulk term linear in \(\alpha\). Flipping one mode
\(\nu\mapsto-\nu\) reproduces \(\Delta\Dcal=2\nu(1-\nu^2)/3\) of
Proposition~\ref{prop:Dflip}.
\end{proposition}

\begin{proposition}[Dressed same-channel jets; executed]
\label{prop:v10dressed}
In scheme A, all dressed same-channel values vanish at \(\alpha=1\)
(F), and CR\(=\)F at \(\alpha=\tfrac12\), on every key
(\loc{logs/v10/jet/jet_analysis.json}). Examples:
D05a \((0,0)\) gives
\(\kappa^2(\alpha-1)(2\alpha-1)/(6\pi\alpha)\) (CR) and
\(\kappa^2(\alpha-1)(2\alpha-1)(\alpha^2-3\alpha+1)/(6\pi\alpha^3)\) (F),
each times \(p_+p_-\). D11 \((1,1)\) gives
\(\kappa^2(\alpha^2-1)(11\alpha^2+1)/(2880\pi\alpha^3)\) in scheme A
and \(-\kappa^2(10\alpha^2+1)/(2880\pi\alpha^3)\) in scheme Z. Of the
74 keys with weight-0 \(r^2\) content (CR), 38 are complete. The
other 36 (all of D01, D06, and 9 of D03) also need mixed pairs and
are \PARTIAL.
\end{proposition}
The D05a value is a slice-jet coefficient under the declared
dressing. It is not compared with the earlier retained statement
``D05a part \(=0\)''. That statement came from a different reduced
model (a large-radius count without the slice), and no
supersession is claimed.

\paragraph{The mixed \(r^2\) trace.}
Term-by-term integration gives
\[
 T_1(x;\mu,\nu)=\tfrac14(4x(1-x))^2\tfrac{\Gamma(h+2)\Gamma(h+1)}{\Gamma(\mu+1)\Gamma(\nu+1)}x^\mu(1-x)^\nu
 F_4(h+2,h+1;\mu+1,\nu+1;x^2,(1-x)^2)
\]
with \(h=(\mu+\nu)/2\). The double series has power divergences
\(N^{1/2-j}\) and no logarithms. Write \(F_4(h+2,h+1)=(h+1+\theta)(h+\theta)F_4(h,h+1)/(h(h+1))\), with \(\theta\) the
Euler operator, and apply Bailey's reduction \cite{BaileyF4} off the boundary,
\(\lambda=1-t^2\). The finite part is then
\(h(h+1)\phi_0-2(h+1)\phi_2+2\phi_4\) in the Taylor coefficients
of \(\Phi(t)\) (\loc{code/v10/t1_fp.py}). It agrees with an
independent shell-by-shell zeta continuation to
\(10^{-19}\)--\(10^{-24}\) (\loc{logs/v10/jet/t1_fp_check.json}). It
reproduces \(\mu(\mu^2-1)/3\) exactly at \(\mu=\nu\) for every
\(x\), and its endpoint values are \(T_1(0)=\mu(\mu^2-1)/3\) and
\(T_1(1)=\nu(\nu^2-1)/3\). After \(x\)-integration, the shell
continuation is \emph{not} reliable, because of non-uniform
endpoint asymptotics (\loc{logs/v10/jet/mixed_r2.json}). It is used
nowhere.

%======================================================================
\section{v10: consequences of the jets}
\label{sec:v10conseq}
%======================================================================

\paragraph{D04.}
The constant density for the jet term is
\(\rho_{(1,0)}=\kappa^2/4\pi\) (\loc{code/v10/jet_density_D04.py}).
Under the declared rule, the D04 part of the second \(\phi\phi\)
value is
\(\tfrac{\kappa^2}{4\pi}M_{10}\), where
\(M_{10}=\sum_{k\ge0}\varphi(k/\alpha)+\sum_{j\ge1}\varphi(j/\alpha-1)\)
is \(k\)-\(\zeta\)-regulated and \(\varphi(n)=\int_0^1T_1(x;n+1,n)\,\mathrm dx\).
It is branch independent, since only channels 0 and 1 enter. We
have \(\varphi\) at 27 points (for example
\(\varphi(0)=-0.0578297132876329558555\) and
\(\varphi(3/2)=2.32785224733468215021\);
\loc{logs/v10/jet/phi/}). No closed form was found in the basis
\(\{n^j,n^j\psi_1(n),n^j\psi(n)\}\) (\loc{code/v10/phi_fit.py}). The
large-\(n\) runs needed for the mode sum were stopped for time.
\(M_{10}\) is therefore \emph{not evaluated}, and the D04 part is
\PARTIAL. A further caveat: the bare \(r^2\) insertion between
channels 1 and 0 omits the frame intertwining \(H_I\), which the
spin-weight-one channel may require. The candidate is therefore a
declared-rule quantity, not a derived value.

\paragraph{The 20 jet-level Ward combinations, the 34- and 18-dimensional
projections, and \(e_5,e_6\).}
These need catalogue values on the engine's \(m=2\) and weight
\(\pm1\) atoms. The jet catalogue now supplies the weight-0
same-channel part (Proposition~\ref{prop:v10dressed}). It does not
supply the following:
\begin{enumerate}[label=(\alph*)]
\item the mixed weight-0 pairs (84 entries), which need \(M_{st}\)
mode sums as for D04;
\item weight \(\pm1\) (412 mixed, 504 same entries). These need
\(k\to k\pm1\) traces with an \(r\,e^{\pm\ii\theta}\) insertion, at
\(w=\tfrac12\), where scheme Z has a pole: a genuine
logarithm;
\item weight \(\pm2\), which needs \(k\to k\pm2\) traces;
\item a dictionary from jet keys to engine \(m=2\) atoms. Such a
dictionary was not constructed.
\end{enumerate}
No residual was evaluated, and none is reported as zero. Status:
\OPEN, with the obstructions (a)--(d).

%======================================================================
\section{v10: vector-key finite parts in a declared scheme}
\label{sec:v10vector}
%======================================================================

\begin{proposition}[Executed]
\label{prop:v10vector}
Per mode, take scheme Z with \(\Tr_{\rm mode}e^{-A_\mu}=-\mu/2\); this
is consistent with \(\int g=1/(2(\mu+\nu))\). Then for a pair density
\(\rho\) with \(R=\int_0^x\rho\),
\(I_\rho=-R(1)\nu/2-(\mu^2-\nu^2)\int_0^1Rg\). Pair the modes at
the same \(k\) (channel 0: \(|k|/\alpha\); channel 1:
\(|k+\alpha|/\alpha\)), and regulate the sum by \(|k|^{-\varepsilon}\)
(\loc{code/v10/vector_fp.py}, \loc{vector_combine.py}). Then, for
D01\(+\)D03 in the v7 normalization:
\begin{enumerate}[label=(\alph*)]
\item the \(1/\varepsilon\) residues equal the v9 logarithms
\((\alpha/2)L_{\rm key}\) of Proposition~\ref{prop:v9vector} to
\(\le1.1\times10^{-20}\) at all four angles. For keys \((1,1)\) and
\((2,2)\) the pole is \(-\alpha\kappa^2P_L^2/(60\pi\varepsilon)\);
\item the off-diagonal keys \((1,6),(2,8),(6,1),(8,2)\) have finite
part 0, because D01 and D03 cancel digit for digit;
\item the diagonal keys have finite part
\(+c(\alpha)\kappa^2P_L^2\) on \((1,1),(2,2)\) and
\(-c(\alpha)\kappa^2P_L^2\) on \((6,6),(8,8)\), with \(c\) given in
Table~\ref{tab:v10vec}. They are branch independent.
\end{enumerate}
\end{proposition}
\begin{table}[ht]\centering\small
\caption{Vector-key finite parts (declared \(k\)-\(\zeta\) scheme).}\label{tab:v10vec}
\begin{tabular}{@{}ll@{}}
\toprule
\(\alpha\) & \(c(\alpha)\) (\(k\)-\(\zeta\) finite part)\\
\midrule
\(3/10\) & \(0.0206627201453\)\\
\(41/100\) & \(0.0204800603712\)\\
\(3/5\) & \(0.0203890448352\)\\
\(7071/10000\) & \(0.020476362628\)\\
\bottomrule
\end{tabular}
\end{table}
The finite part depends on the scheme. Rescaling the regulator
\(|k|\to\lambda|k|\) shifts \(c\) by
\((\alpha/60\pi)\log\lambda\), and a sharp cutoff differs from it
by the \(\zeta(-1),\zeta(0)\) terms. It is a declared-scheme value
and not an apex coefficient.

%======================================================================
\section{v10: changes to v9 statements}
\label{sec:v10changes}
%======================================================================
\begin{itemize}
\item Table~\ref{tab:v9dJ} discussion: ``PSLQ was not attempted;
no elementary closed form'' is superseded by
Corollary~\ref{cor:v10dJ} (CR derived, F identified).
\item Theorem~\ref{thm:v9ghost}: ``two fitted normalizations'' is
now ``two normalizations derived from tree operators''
(Theorem~\ref{thm:v10lambda}).
\item Theorem~\ref{thm:v9generic}: ``F\((>1/2)\) between 29 and 30''
is now ``exactly 29 at every angle'', conditional on (D)
(Theorem~\ref{thm:v10chamber}).
\item Section~\ref{sec:v9D04}: obstruction (a), ``the slice does not
export jet densities'', is removed by the jet replay. Obstruction
(b) is made precise: the D04 jet term is purely mixed, and scheme Z
regularizes the per-mode power divergence. The mode sum is still
missing.
\item Proposition~\ref{prop:v9vector}: the finite parts are now
computed in a declared scheme.
\end{itemize}
No v9 number was found to be in error. The frozen-tree digests
before and after v10 are in Appendix~\ref{app:v10artifacts}.

%======================================================================
\section{What v11 adds, and what it does not}
\label{sec:whatv11}
%======================================================================

v11 keeps every v10 statement \cite{CheungConeV10}, except those
amended in Section~\ref{sec:v11changes}. The new items and their
status are:

\begin{center}
\small
\begin{longtable}{@{}L{.30\linewidth}L{.62\linewidth}@{}}
\caption{v11 verdicts (all executed; hashes in Appendix~\ref{app:v11artifacts}).}\label{tab:v11verdict}\\
\toprule
Item & Status\\
\midrule
Mixed \(r^2\) trace evaluator & \emph{Corrected.} The v10 evaluator used a Cauchy radius that is too large near the endpoints and for large \(h\). Every v10 \(\varphi\) value is withdrawn (Section~\ref{sec:v11bug}).\\
F mode \(I_{\rm F}\) & \emph{Derived} for \(0<\mu<2\) (Theorem~\ref{thm:v11IF}). The v10 PSLQ form is now a theorem there. As a by-product, \(\Delta J=-(\mu-1)B/W_1\) is a single incomplete-Beta integral.\\
Degree bound (D) & \emph{Machine-checked} by symbolic degree tracking on all 8 chambers and cross-checked against the exact engine (Proposition~\ref{prop:v11D}). Theorem~\ref{thm:v10chamber} is now unconditional in its scope.\\
\(M_{10}\) (D04 mode sum) & \emph{Executed} in two declared pairings (P0 same-\(k\); P1, the unit angular-momentum shift used as the frame intertwiner \(H_I\)) and two regulators, at four angles (Table~\ref{tab:v11M10}). The precision is set by the data, about 4--8 digits, not 30. The pairing and the regulator are \emph{declarations}, so D04 stays \PARTIAL.\\
Weight \(\pm1,\pm2\) and mixed jet traces & \OPEN{} (Section~\ref{sec:v11open}).\\
Jet-key \(\to m=2\) atom dictionary; 20 Ward combinations; 34-/18-dim maps; \(e_5,e_6\) & \OPEN. Not constructed and not evaluated; no residual is reported.\\
Transverse loop momentum & \OPEN; obstruction stated in Section~\ref{sec:v11open}.\\
Registry & Unchanged (258 ids); no new sheet directory written.\\
\bottomrule
\end{longtable}
\end{center}

%======================================================================
\section{v11: correction of the v10 mixed-trace evaluator}
\label{sec:v11bug}
%======================================================================

The v10 evaluator \loc{code/v10/t1_fp.py} extracts the Taylor
coefficients \(\phi_0,\phi_2,\phi_4\) of \(\Phi(t)\) by a Cauchy
integral on \(|t|=0.04\). Here \(\lambda=1-t^2\), and
\(T_1=\tfrac14(4x(1-x))^2\frac{\Gamma(h+2)\Gamma(h+1)}{\Gamma(\mu+1)\Gamma(\nu+1)}x^\mu(1-x)^\nu
\frac{h(h+1)\phi_0-2(h+1)\phi_2+2\phi_4}{h(h+1)}\). v11 finds two
defects.
\begin{enumerate}[label=(\roman*)]
\item The discriminant of the Bailey change of variables has a second
root at \(|t_2|=2\sqrt{x(1-x)}/|1-2x|\). For \(x\to0,1\), \(|t_2|\)
falls below \(0.04\), and the extracted coefficients are then wrong.
\item Even inside the disc, the integrand grows roughly like
\(e^{c h|t|}\). The extraction loses about 1.3 digits per unit of
\(h\): for \(\mu=\nu=7\), the exact diagonal value
\(\mu(\mu^2-1)/3\) is missed by \(2.6\times10^{-10}\).
\item Near the endpoints the natural variable is \(t/|t_2|\), and the
same growth sets in at the scale \(|t_2|/h\). With
\(r=|t_2|/4\) alone, \(T_1\) is aliased for \(h\gtrsim60\), and the
error does not depend on the working precision. For example, an
intermediate v11 run gave \(F(61,60)\) below both endpoint values.
\end{enumerate}
The v11 evaluator \loc{code/v11/t1v11.py} uses
\(r=\min(0.25,|t_2|/4,1/(1+h),|t_2|/(1+h))\) with 56 nodes, and
\loc{code/v11/m10_eval.py} adds two measures. The half \(x>\tfrac12\)
is evaluated through the symmetry
\(T_1(x;\mu,\nu)=T_1(1-x;\nu,\mu)\) (checked to 25 digits), and the
working precision is raised by 1.5 digits per decade near the
endpoint. Tanh--sinh is capped at degree 4, since degrees 3, 4 and 5
agree to \(3\times10^{-23}\) relative. On the diagonal the evaluator
reproduces \(\mu(\mu^2-1)/3\) to \(\le10^{-24}\) at \(\mu=20,60\)
and \(x=0.05,0.3\). Near the endpoints, halving the radius and
raising the node count to 80 change \(T_1\) by less than
\(10^{-14}\) relative at \(h=60.5\) and \(h\approx131\). Three
intermediate v11 runs with incomplete fixes are superseded and feed
no printed number (\loc{logs/v11/M10/cache_superseded*/}). As a
final sanity check, every cached \(F(\mu,\nu)\) lies between its two
endpoint values: true for 403 of 404 entries; the exception \(F(1,0)\) (\loc{logs/v11/M10/m10_sanity.json}) has both endpoint values zero, so its negative value is no contradiction.

Consequences: the corrected value is \(\varphi(3/2)=2.32745055013616982735\),
while v10 quoted \(2.32785224733468215021\), an error of
\(4.0\times10^{-4}\). The corrected value is
\(\varphi(0)=F(1,0)=-0.0578297132876329675\), while v10 quoted
\(-0.0578297132876329558555\); these agree to \(1.2\times10^{-17}\),
since the defect is mild at small \(h\). Because the evaluator was
defective, all 27 v10 \(\varphi\) values are withdrawn as a set,
together with the
v10 negative result ``no closed form in
\(\{n^j,n^j\psi_1(n),n^j\psi(n)\}\)'' are withdrawn. The v10
mid-range checks (\(x\in[0.3,0.8]\), \loc{logs/v10/jet/t1_fp_check.json})
and the diagonal identity are unaffected, and so are the
same-channel jets, scheme Z and the vector finite parts, none of
which use \(T_1\) off the diagonal.

%======================================================================
\section{v11: derivation of the F mode}
\label{sec:v11IF}
%======================================================================

\begin{theorem}[F mode; \loc{code/v11/fmode_derive.py}]
\label{thm:v11IF}
Let \(0<\mu<2\), \(\nu=2-\mu\), \(w(x)=x^{\mu-1}(1-x)^{1-\mu}\),
\(W(x)=\int_0^xw\) and \(W_1=W(1)=\Gamma(\mu)\Gamma(2-\mu)\). Then:
\begin{enumerate}[label=(\alph*)]
\item \(g(x;\mu,2-\mu)=\dfrac{W(x)\,(W_1-W(x))}{4W_1\,x(1-x)\,w(x)}\).
Both \({}_2F_1\) factors of Proposition~\ref{prop:v9g} are
incomplete Beta functions \cite{DLMFBeta}, by
\({}_2F_1(a,a-1;a+1;x)=a\,x^{-a}B_x(a,2-a)\).
\item \(I_{\rm F}=I_{\rm CR}+(\mu-1)B/W_1\) and
\(\Delta J=-(\mu-1)B/W_1\), where
\(B=\int_0^1(3-2x)W^2/w\,\mathrm dx\). The linear term
\(\int(3-2x)W/w=\mu^{-1}\sum_n(n+1)!K_n/(\mu+1)_n\) carries the same
\(K_n\) as the CR mode, and so equals \(-I_{\rm CR}/(\mu-1)\).
\item \((x(1-x)/w)'=(2-\mu-2x)/w\) gives
\(B=-2W_1E+(1+\mu)C\), with
\(E=\tfrac16-\tfrac{\mu(\mu+1)}{12}+\tfrac{\mu(\mu+1)(\mu+2)}{72}\)
and \(C=\int_0^1W^2/w\).
\item The series for \(W/w\), together with
\(\int_0^1Wx^m=W_1\bigl(1-(\mu)_{m+1}/(m+2)!\bigr)/(m+1)\), reduces
\(C/W_1\) to the v10 Gamma sums \(G_2,G_3\) plus two rational sums.
The result is
\[
\frac{C}{W_1}=\frac{(\mu-2)(5\mu-3)}3+(\mu-1)^2\Bigl[\psi_1(\mu-1)-\frac{\pi^2}6\Bigr].
\]
\item Hence
\(I_{\rm F}(\mu)=\tfrac{(\mu-1)(59\mu^3-57\mu^2-62\mu+48)}{36}-(\mu-1)^3(\mu+1)\tfrac{\pi^2}6\),
and the difference from the v10 PSLQ form simplifies to \(0\) in
SymPy.
\end{enumerate}
\end{theorem}
Checks: \(C/W_1\) agrees with direct quadrature to
\(\le1.2\times10^{-30}\) at \(\mu=7/4,4/3,10/7\)
(\loc{logs/v11/fmode_derive.json}). The reduction (a)--(b) agrees
with the v10 closed form to \(\le10^{-31}\) at \(\mu=7/4,5/4,3/2\)
(\loc{logs/v11/fmode_reduce.json}). The Euler constant and
\(\psi(\mu-1)\) cancel identically. Scope: \(0<\mu<2\), that is
\(\alpha>\tfrac12\), where \(B_x(\mu,2-\mu)\) converges. This is the
range in which the F mode enters \(\Delta J\).

%======================================================================
\section{v11: machine check of the degree bound (D)}
\label{sec:v11D}
%======================================================================

\begin{proposition}[\loc{code/v11/degree_check.py}, \loc{code/v11/degree_crosscheck.py}]
\label{prop:v11D}
In the frozen v8 engine, \(\alpha\) enters the ST operators only
through the derivative coefficients \((m\pm a)/2\), with
\(a=nt+s\), \(t=1/\alpha\) and \(m=\varepsilon|a|+2j\). The row
basis is a pivot \emph{selection} of candidates, with no
\(\alpha\)-dependent combination, and the candidate evaluation is
multilinear in the operator outputs and polynomial in \(q\).
Replace \(t\) by a SymPy symbol and apply \(K_2\), \(E(-q)\) and
\(C_2T\) to every unit atom up to \(\mathrm{DM}=6\) on each of the
8 chambers (CR and F, \(c_1\)--\(c_4\)). The maximal degrees in
\(t\) are then 1, 2 and 1, so (D) holds. The atom counts are
\(c\): 22--52, \(h\): 54--130, \(b\): 16--36 per chamber
(\loc{logs/v11/degree_check.json}). At \(t=1/\alpha_0\) the symbolic
outputs equal the exact numeric engine on 3744 coefficients (CR
\(c_4\), F \(c_3\), CR \(c_1\)), with 0 mismatches
(\loc{logs/v11/degree_crosscheck.json}).
\end{proposition}
The hypothesis of Theorem~\ref{thm:v10chamber} is therefore
discharged. The chamber ranks (26 on the six \(N=34\) chambers, 29
on F \((\tfrac12,\tfrac23)\) and \((\tfrac23,1)\), at every
\(\alpha\) in the chamber) now hold without condition, in the scope
of the declared ST rows of the frozen engine.

\paragraph{Certificate status.} (i) The v8 four-point exact rank
certificate is unaffected by v11 and remains
\CPIR{} in its v8 scope. (ii) The chamber-rank certificate is
\CPIR{} in the following precise scope: the rank of the declared ST
row matrix of the frozen v8 engine (\(\mathrm{DM}=6\), \(6\times6\)
\(q\)-grid) is 26 at every \(\alpha\) in the six \(N=34\) chambers,
and 29 at every \(\alpha\) in the two F chambers above
\(\tfrac12\). Every ingredient is executed: exact minors at the
sampled angles, the constancy and affinity arguments of
Theorem~\ref{thm:v10chamber}, and the machine-checked (D). The
certificate does not cover the identification with the interacting
basis or the walls between chambers.

%======================================================================
\section{v11: the D04 mode sum \texorpdfstring{\(M_{10}\)}{M10}}
\label{sec:v11M10}
%======================================================================

\paragraph{Definitions.} Take \(F(\mu,\nu)=\int_0^1T_1(x;\mu,\nu)\,\mathrm dx\)
with the corrected evaluator (tanh--sinh, 30 digits; endpoint limits
for \(x<10^{-22}\)). Two pairings of channel~1 with channel~0 are
declared:
\begin{itemize}
\item P0 (the v10 rule, same \(k\)):
\((\mu,\nu)=(\tfrac{k+\alpha}\alpha,\tfrac k\alpha)\) for
\(k\ge0\), and \((\tfrac{j-\alpha}\alpha,\tfrac j\alpha)\) for
\(j\ge1\).
\item P1 (the unit angular-momentum shift; v11's implementation of
the frame intertwiner \(H_I\), as a \emph{declared} pairing):
\((\tfrac{k+\alpha}\alpha,\tfrac{k+1}\alpha)\) for \(k\ge0\), and
\((\tfrac{|\alpha-k|}\alpha,\tfrac{|1-k|}\alpha)\) for \(k\ge1\).
\end{itemize}
Both involve only channels 0 and 1, so both are branch independent;
at \(\alpha=\tfrac12\), CR\(=\)F holds trivially. Two regulators are
used: \(R_k\), weight \(|k|^{-\epsilon}\) (the v10 declaration),
and \(R_h\), weight \(h^{-\epsilon}\) with \(h=(\mu+\nu)/2\). The
terms \(k<20\) are summed exactly. For \(k\ge20\) a QR least-squares
fit \(F\sim\sum_{i<12}c_iu^{3-i}\) on \(k=20,22,\dots,56\) is summed
through Hurwitz \(\zeta(i-3,u_{20})\). The \(u^{-1}\) term gives a
pole \(c_4/\epsilon\) (reported) and the finite part
\(-c_4\psi(u_{20})\), minus \(c_4\log(1/\alpha)\) for \(R_h\). The
error estimate is the spread over fit orders 10--13 and over fits on
half the points. A power expansion without logarithms is an
\emph{ansatz} for the large-\(k\) asymptotics. It is supported by the
fit residuals, not derived (\loc{code/v11/m10_sum.py},
\loc{logs/v11/M10/m10.json}).

\paragraph{Check at \(\alpha=1\).} P1 pairs identical modes, and
\(F(\mu,\mu)=\mu(\mu^2-1)/3\) exactly. The machinery returns
\(-19/180\) under \(R_k\) (to \(2.8\times10^{-26}\)) and \(11/180\) under \(R_h\)
(to \(2.2\times10^{-26}\)). The latter is the scheme-Z flat value
\(\Dcal^Z_0(1)=11/180\) of Proposition~\ref{prop:v10Z}. The difference
\(1/6\) is a pure regulator effect: the shift between the index
\(k\) and the mode label changes the zeta values. It shows that
\(M_{10}\) is \emph{not} regulator independent, and that only \(R_h\)
is consistent with scheme Z on the diagonal.

\begin{center}
\small
\begin{longtable}{@{}lllll@{}}
\caption{\(M_{10}\) finite parts (\(\kappa^2/4\pi\) times these is the D04 term under the stated declaration). Error = spread estimate; pole = coefficient of \(1/\epsilon\).}\label{tab:v11M10}\\
\toprule
\(\alpha\) & pairing, reg. & \(M_{10}\) (FP) & error & pole\\
\midrule
\(3/10\) & P0, \(R_h\) & \(0.1157438826\) & \(4.4\times10^{-9}\) & \(-0.0075\)\\
\(3/10\) & P0, \(R_k\) & \(0.0973390865\) & \(9.0\times10^{-10}\) & \(-0.0075\)\\
\(3/10\) & P1, \(R_h\) & \(0.0169054255\) & \(1.2\times10^{-8}\) & \(0.0196605\)\\
\(3/10\) & P1, \(R_k\) & \(1.9509542118\) & \(4.4\times10^{-7}\) & \(0.0196605\)\\
\(1/2\) & P0, \(R_h\) & \(-0.023918875\) & \(5.0\times10^{-8}\) & \(-0.0125\)\\
\(1/2\) & P0, \(R_k\) & \(-0.048208215\) & \(1.2\times10^{-8}\) & \(-0.0125\)\\
\(1/2\) & P1, \(R_h\) & \(-0.023918875\) & \(5.0\times10^{-8}\) & \(-0.0125\)\\
\(1/2\) & P1, \(R_k\) & \(0.2017917849\) & \(2.2\times10^{-7}\) & \(-0.0125\)\\
\(7071/10000\) & P0, \(R_h\) & \(-0.043222064\) & \(2.6\times10^{-7}\) & \(-0.0176775\)\\
\(7071/10000\) & P0, \(R_k\) & \(-0.071445663\) & \(6.5\times10^{-8}\) & \(-0.0176775\)\\
\(7071/10000\) & P1, \(R_h\) & \(0.0208317775\) & \(5.7\times10^{-8}\) & \(-0.00387076\)\\
\(7071/10000\) & P1, \(R_k\) & \(-0.044268878\) & \(1.7\times10^{-8}\) & \(-0.00387076\)\\
\(1\) & P0, \(R_h\) & \(-0.047837750\) & \(1.3\times10^{-6}\) & \(-0.025\)\\
\(1\) & P0, \(R_k\) & \(-0.079087750\) & \(3.5\times10^{-7}\) & \(-0.025\)\\
\(1\) & P1, \(R_h\) & \(0.0611111111\) & \(2.2\times10^{-26}\) & \(5.36872e-27\)\\
\(1\) & P1, \(R_k\) & \(-0.105555555\) & \(2.8\times10^{-26}\) & \(2.19251e-27\)\\
\bottomrule
\end{longtable}
\end{center}

\paragraph{Reading of Table~\ref{tab:v11M10}.} The leading fitted coefficient agrees with the prediction \(c_0=\alpha^{-3}/3\) (for \(u=k\) and for \(u=\alpha h\)) from \(F\sim h(h^2-1)/3\) to relative \(1.7\times10^{-15}\) in every family (\loc{logs/v11/M10/m10_leading.json}). This is the only analytic input to the asymptotics; the subleading coefficients are fitted, not derived.
The intertwiner choice is not a small effect: \(M_{10}^{\rm P1}-M_{10}^{\rm P0}\) under \(R_h\) is \(3/10\): \(-0.098838457\); \(7071/10000\): \(0.064053842\).
At \(\alpha=\tfrac12\) the P1 pairs \((2k+1,2k+2)\) are, after reindexing, the P0 pairs of the negative family and vice versa, so the two pairings must agree under the label-based regulator \(R_h\) (not under \(R_k\)). They agree to \(5.1\times10^{-12}\), a non-trivial check of the head/tail machinery on non-exact data.
The P0 pole coefficient equals \(-\alpha/40\) at all four angles and for both regulators, to \(9.7\times10^{-12}\) (identified numerically, not derived).
The pole coefficients are nonzero in general; a declared minimal subtraction of the \(1/\epsilon\) term is part of the finite-part definition.

\paragraph{Status of D04.} Under a declared pairing and regulator,
the D04 part of the second \(\phi\phi\) value is
\(\tfrac{\kappa^2}{4\pi}M_{10}\), with \(M_{10}\) as in
Table~\ref{tab:v11M10}. Neither the pairing nor the regulator is
derived. The difference between P0 and P1, and between \(R_k\) and
\(R_h\), measures the size of the undeclared choice. Status:
\PARTIAL{} (executed, scheme-dependent, precision about 4--8
digits). PSLQ in the basis
\(\{1,\alpha^{\pm1},\alpha^2,\alpha^{-2},\alpha^{-3},\zeta(3)/\alpha^2,\pi^2,\log\alpha\}\)
with coefficients up to 1000: relations among the basis elements alone (coefficient of \(M_{10}\) zero; e.g.\ \(1/\alpha^2=2/\alpha\) at \(\alpha=\tfrac12\)) for 7071/10000 P0 Rk, 7071/10000 P0 Rh, 1/2 P0 Rk, 1/2 P0 Rh, 1/2 P1 Rk, 1/2 P1 Rh; these are basis degeneracies; relations involving \(M_{10}\) for 3/10 P0 Rk, 3/10 P0 Rh, 3/10 P1 Rk, 3/10 P1 Rh, 7071/10000 P1 Rk, 7071/10000 P1 Rh. With ten basis elements and a tolerance of about \(10^{-8}\), small-coefficient relations exist generically, so these are \emph{not} credible and are not claimed. No closed form is claimed.
(\loc{logs/v11/M10/m10_pslq.json}).

%======================================================================
\section{v11: items left open, with precise obstructions}
\label{sec:v11open}
%======================================================================

\begin{enumerate}[label=(\alph*)]
\item \emph{Weight \(\pm1\) traces.} The \(k\to k\pm1\) traces with an
\(r\,e^{\pm\ii\theta}\) insertion were not built in v11. The task is
to split the scheme-Z pole at \(w=\tfrac12\) into its residue (the
genuine, scheme-independent log coefficient) and the finite part
(scheme dependent). Doing so needs the two-channel finite-part
evaluator at weight one: the analogue of \(T_1\) with
\({}_2F_1\) parameters shifted by \(\tfrac12\), and its own mode sum.
That evaluator was not written. Heuristically, in the same-channel
analogue the log arises from the \(z^{-1}\) term of
\(e^{-z}I_\mu(z)\sim(2\pi z)^{-1/2}[1-(4\mu^2-1)/(8z)+\dots]\), but
this is not a computation of the mixed coefficient, and no number is
claimed. Status: \OPEN.
\item \emph{Weight \(\pm2\) and the 84 mixed weight-0 entries.} The
corrected evaluator applies to the weight-0 mixed entries. Each
needs a mode sum like \(M_{10}\), with a declared pairing, and none
was run. Status: \OPEN.
\item \emph{Dictionary, Ward combinations, projections, \(e_5,e_6\).}
No map from jet keys \((s,t;w)\) to the engine's \(m=2\) atoms
\((C,n,j,i)\) was constructed. Without it, the 20 jet-level Ward
combinations cannot be formed, and neither can the remaining 12
coordinates of the 18-dim map, the 34-dim projection, or \(e_5,e_6\)
separately. No residual was evaluated, and none is reported as zero.
Status: \OPEN.
\item \emph{Transverse loop momentum.} The jet catalogue evaluates the
loop at transverse momentum on the slice. Going off the slice needs
the \(p_\perp\) dependence of every vertex in the catalogue, plus a
\(p_\perp\) integral in a declared scheme. The v8 atoms are
polynomial in the slice momenta \(q\) only, and the catalogue
exports no \(p_\perp\) jets. That export is the missing piece.
Status: \OPEN.
\end{enumerate}

%======================================================================
\section{Changes to v10 statements}
\label{sec:v11changes}
%======================================================================

\begin{enumerate}[label=(C\arabic*)]
\item Section~\ref{sec:v10conseq}, D04: the quoted values
\(\varphi(0)=-0.0578297132876329558555\) and
\(\varphi(3/2)=2.32785224733468215021\), the 27 cached values in
\loc{logs/v10/jet/phi/}, and the negative closed-form search are
\emph{withdrawn}, because the evaluator was defective
(Section~\ref{sec:v11bug}). Corrected values are in
Section~\ref{sec:v11bug}. ``\(M_{10}\) not evaluated'' is superseded
by Section~\ref{sec:v11M10}.
\item Section~\ref{sec:v10jet}: the claim that the closed
finite-part evaluator for the mixed trace is valid for all
\(x\in(0,1)\) is restricted. With the v10 radius it is valid only
where \(0.04<|t_2|\) and \(h\) is small. The v11 radius removes the
restriction.
\item Proposition~\ref{prop:v10IF}: ``identified by PSLQ, not
derived'' becomes \emph{derived} for \(0<\mu<2\)
(Theorem~\ref{thm:v11IF}).
\item Theorem~\ref{thm:v10chamber}: ``Assume (D) \dots\ not separately
machine-checked'' is replaced by Proposition~\ref{prop:v11D}, so the
theorem holds without condition. The v10 release items (xvii) and
(xiv) lose ``conditional on (D)''.
\item The v10 artifact manifest hashed \loc{CHECKPOINT_v10.md} one
append early; the corrected hash is in
Appendix~\ref{app:v11artifacts}.
\item Open table: the rows ``\(\Delta J\)'', ``Exact \(\Span\)
certificate'' and ``D04'' are updated (Table~\ref{tab:open}).
\end{enumerate}
No other v10 statement is changed.

%======================================================================
\section{What v12 adds, and what it does not}
\label{sec:whatv12}
%======================================================================

v12 keeps every v11 statement \cite{CheungConeV11} except those
amended in Section~\ref{sec:v12changes}. The new items and their
status are:

\begin{center}
\small
\begin{longtable}{@{}L{.30\linewidth}L{.62\linewidth}@{}}
\caption{v12 verdicts (all executed; hashes in Appendix~\ref{app:v12artifacts}).}\label{tab:v12verdict}\\
\toprule
Item & Status\\
\midrule
Closed form of the mixed \(r^2\) mode functional \(F(\mu,\nu)\) & \emph{Identified} and verified against all 414 quadrature values to \(4.8\times10^{-28}\) relative (Proposition~\ref{prop:v12F}). Not derived.\\
D04 pairing and frame intertwiner & \emph{Derived} from the channel labelling and the constant density: \(e^{\pm\ii\theta}\) shifts \(k\) by \(\pm\alpha\), which is the channel offset itself. The pairing is the same-\(k\) rule P0, and P1 is excluded (Proposition~\ref{prop:v12pair}).\\
Large-\(k\) asymptotics and the P0 pole & \emph{Exact} given the closed form: only odd powers of \(1/h\) appear, with \(a_4(d)=d^2(d^2-4)/240\), so the pole is \(-\alpha/40\) (Corollary~\ref{cor:v12pole}).\\
Regulator & \(R_k,R_h,R_\mu,R_\nu\) differ by an explicit local counterterm (Proposition~\ref{prop:v12ct}, verified to \(10^{-39}\)). The XVI-P prescription (DRED, \(\overline{\rm MS}\), normalization N4) induces no cone mode regulator, so the choice is a normalization. D04 is unique once it is declared; v12 declares \(R_h\).\\
\(M_{10}\) & 40 digits at four angles (Table~\ref{tab:v12M10}). Closed forms at \(\alpha=1,\tfrac12\) found by PSLQ; none at \(\tfrac3{10}\) or \(\tfrac{7071}{10000}\).\\
Chamber walls \(\alpha=\tfrac13,\tfrac12,\tfrac23\) & Exact \(\mathbb Q(\ii)\) certificates at the walls (Table~\ref{tab:v12walls}).\\
Weight \(\pm1,\pm2\) traces; 84 mixed weight-0 sums; dictionary; 20 Ward combinations; 34-/18-dim maps; \(e_5,e_6\) & \OPEN{} (Section~\ref{sec:v12open}). No residual is reported.\\
Transverse loop momentum & \OPEN; feasibility statement in Section~\ref{sec:v12open}.\\
Registry & Unchanged (258 ids, re-fetched 04:40 UTC); no new sheet directory.\\
\bottomrule
\end{longtable}
\end{center}

%======================================================================
\section{v12: closed form of the mixed \texorpdfstring{\(r^2\)}{r2} mode functional}
\label{sec:v12F}
%======================================================================

\begin{proposition}[Identified closed form; \loc{code/v12/F_closed.py}]
\label{prop:v12F}
Let \(h=(\mu+\nu)/2\) and \(d=\mu-\nu\). Then, to the precision of
every evaluation made,
\[
F(\mu,\nu)=\int_0^1T_1(x;\mu,\nu)\,\mathrm dx
=\frac{h^3}{3}+\Bigl(\frac{d^2}{6}-\frac13\Bigr)h+d^2(d^2-4)\,\Psi(h),
\]
\[
\Psi(h)=\frac{h^3}{48}+\frac{h^2}{96}-\frac{5h}{288}-\frac1{96}-\frac{h^2(h^2-1)}{48}\,\psi_1(h).
\]
\end{proposition}
\noindent\emph{How it was found.}
\begin{enumerate}[label=(\alph*)]
\item Fits to the v11 cache at fixed \(d\) give
\(F\sim\sum_ia_{2i}(d)h^{3-2i}\): the odd coefficients are
\(\le10^{-17}\) (\loc{code/v12/phi_asym_fit.py}).
\item At \(h=\tfrac32,\tfrac72,5\) and \(d=\tfrac12,1,3\), the ratio
\(\bigl(F-\tfrac{h^3}3-(\tfrac{d^2}6-\tfrac13)h\bigr)/(d^2(d^2-4))\)
is independent of \(d\) to 20 digits, and at \(d=2\) the remainder
vanishes to \(10^{-29}\) (\loc{code/v12/structure_test.py}).
\item PSLQ puts \(\Psi\) in \(\mathbb Q+\mathbb Q\pi^2\) at those
\(h\). For example, \(\Psi(5)=(987-100\pi^2)/48\) and
\(\Psi(\tfrac32)=(448-45\pi^2)/1536\).
\item Matching the asymptotic coefficients \(1/240\) and \(-1/840\)
fixes the coefficient of \(\psi_1\) as \(-h^2(h^2-1)/48\). This
predicts the next two coefficients, \(1/840\) and \(-1/440\); the
fits confirm both.
\end{enumerate}
\emph{Verification.} The formula agrees with all 414 cached
tanh--sinh values (v11 evaluator; \(h\) up to 190, \(d\) in
\(\{\tfrac12,0.414,1,2,\tfrac73,3,\dots\}\)). The maximum differences
are \(4.5\times10^{-22}\) absolute and \(4.8\times10^{-28}\) relative
(\loc{logs/v12/F_closed.json}). \emph{Status.} The formula is
identified and verified, not derived. A derivation would follow from
a closed evaluation of the Appell \(F_4\) finite part; v12 did not
attempt one.

\begin{corollary}[Asymptotics and the P0 pole]
\label{cor:v12pole}
Given Proposition~\ref{prop:v12F}, the expansion
\(\psi_1(h)\sim h^{-1}+\tfrac12h^{-2}+\sum_{n\ge1}B_{2n}h^{-2n-1}\)
gives \(F\sim\tfrac{h^3}3+(\tfrac{d^2}6-\tfrac13)h+d^2(d^2-4)
\bigl(\tfrac1{240h}-\tfrac1{840h^3}+\tfrac1{840h^5}-\tfrac1{440h^7}+\dots\bigr)\).
The expansion has no logarithms and no even powers. In a family
\(h_k=k/\alpha+\sigma\), the coefficient of \(k^{-1}\) is
\(\alpha a_4(d)\): the binomial re-expansion of \(h^{3-j}\), \(j\le3\),
produces no \(k^{-1}\) term. Summed over the two families of a
pairing, the pole is \(2\alpha a_4(d)=\alpha d^2(d^2-4)/120\). For
P0 (\(d=1\)) this is \(-\alpha/40\), the value seen numerically in
v11. For P1 (\(d=1/\alpha-1\)) it reproduces the v11 numbers, for
example \(637/32400\) at \(\alpha=\tfrac3{10}\).
\end{corollary}
The proof of the pole value is complete given
Proposition~\ref{prop:v12F}, and only to that extent.

%======================================================================
\section{v12: the D04 pairing and frame intertwiner}
\label{sec:v12pair}
%======================================================================

\begin{proposition}[Pairing]
\label{prop:v12pair}
Assume (i) the v8 channel labelling: the spin-weight-\(s\) component
with frame label \(k\) has Bessel order \(|k+s\alpha|/\alpha\); and
(ii) the v10 D04 density \(\rho_{(1,0)}=\kappa^2/4\pi\), which is
independent of \(\theta\) in the helicity frame. Then:
\begin{enumerate}[label=(\alph*)]
\item On \(C_\alpha\), angular modes are \(e^{\ii k\theta/\alpha}\),
\(\theta\in[0,2\pi\alpha)\). The frame rotation \(e^{\pm\ii\theta}\)
multiplies a mode by \(e^{\pm\ii\alpha\theta/\alpha}\), so it shifts
\(k/\alpha\) by \(\pm1\), that is, \(k\) by \(\pm\alpha\). This is
the channel offset in (i). The rotation is therefore already
absorbed in the channel orders, and the intertwiner \(H_I\) acts as
the identity on \(k\).
\item The helicity frame is the only frame defined globally on
\(C_\alpha\), because of the holonomy \(e^{2\pi\ii s\alpha}\). A
\(\theta\)-independent insertion in it commutes with
\(-\ii\partial_\theta\), so the \((1,0)\) trace is diagonal in \(k\).
This is pairing P0:
\((|k+\alpha|/\alpha,|k|/\alpha)\) for all \(k\in\mathbb Z\).
\item The v11 pairing P1 (\(k\to k+1\)) would need the insertion
\(e^{\pm\ii\theta/\alpha}\), a scalar angular harmonic that does
not occur in \(\rho_{(1,0)}\). P1 is excluded.
\end{enumerate}
\end{proposition}
Scope: this derivation uses (i) and (ii) as stated. It does not
re-derive the frozen vertex tensors.

%======================================================================
\section{v12: regulators, counterterm, and \texorpdfstring{\(M_{10}\)}{M10} to 40 digits}
\label{sec:v12reg}
%======================================================================

\begin{proposition}[Local counterterm; \loc{code/v12/counterterm.py}]
\label{prop:v12ct}
Let \(f(k)\sim\sum_ic_ik^{3-i}\) be a power expansion with no
logarithms (by Corollary~\ref{cor:v12pole}), and let
\(w_k=s(k+t)\). Then
\[
\mathrm{FP}\sum_kf(k)k^{-\epsilon}-\mathrm{FP}\sum_kf(k)w_k^{-\epsilon}
=c_4\log s+\sum_{m=1}^{4}\frac{(-1)^{m+1}t^m}{m}\,c_{4-m}.
\]
\end{proposition}
\begin{proof}
\(k^{-\epsilon}-w_k^{-\epsilon}=\epsilon k^{-\epsilon}[\log
s+\log(1+t/k)]+O(\epsilon^2)\). Both sums have only simple poles at
\(\epsilon=0\), so the difference of finite parts is the residue of
\(\sum_kf(k)k^{-\epsilon}[\log s+\log(1+t/k)]\), which is the
coefficient of \(k^{-1}\).
\end{proof}
For each pair of regulators, the formula matches the 40-digit
differences to \(\le1.3\times10^{-39}\). In particular,
\(M_{10}^{R_h}-M_{10}^{R_\mu}=M_{10}^{R_h}-M_{10}^{R_\nu}=\alpha/32\)
for P0. \emph{Which regulator?} XVI-P's declared prescription is
DRED with \(\overline{\rm MS}\) and the normalization condition N4,
under which the local ambiguity is fixed by normalization
(\loc{XVI-P/code/p5_entry/prescription_d11.py}, read only). It
defines no mode regulator on the cone. The regulator choice is
therefore exactly the choice of the local term in
Proposition~\ref{prop:v12ct}. v12 \emph{declares} \(R_h\): it is
symmetric in \(\mu\leftrightarrow\nu\), and it reproduces scheme Z on
the diagonal (v11: \(11/180\) at \(\alpha=1\)). Every other value
follows from the counterterm.

\paragraph{\(M_{10}\).} The terms \(k<K\) are evaluated exactly from
Proposition~\ref{prop:v12F}. The tail \(k\ge K\) uses the exact
asymptotic series in the regulator variable, summed through Hurwitz
\(\zeta\). Runs with \((K,J)=(300,36)\) and \((600,48)\) agree to
\(\le2\times10^{-34}\) (\loc{code/v12/m10_exact.py},
\loc{logs/v12/m10_exact_dps40.json}). All values lie inside the
v11 error bars.

\begin{center}
\small
\begin{longtable}{@{}lllll@{}}
\caption{\(M_{10}\) (finite part, 32 digits shown) and the \(1/\epsilon\) pole coefficient. P0 is the derived pairing (Proposition~\ref{prop:v12pair}); P1 is shown only for comparison with v11.}\label{tab:v12M10}\\
\toprule
\(\alpha\) & pairing & reg. & \(M_{10}\) & pole\\
\midrule
\(3/10\) & P0 & \(R_h\) & \(0.11574388259581850963848639109243\) & \(-0.0075\)\\
\(3/10\) & P0 & \(R_k\) & \(0.09733908656337398969381579445921\) & \(-0.0075\)\\
\(3/10\) & P0 & \(R_\mu\) & \(0.10636888259581850963848639109243\) & \(-0.0075\)\\
\(3/10\) & P0 & \(R_\nu\) & \(0.10636888259581850963848639109243\) & \(-0.0075\)\\
\(3/10\) & P1 & \(R_h\) & \(0.01690542553874612785881404187544\) & \(0.019660493827\ldots\)\\
\(3/10\) & P1 & \(R_k\) & \(1.95095421184601838796068716967526\) & \(0.019660493827\ldots\)\\
\(1/2\) & P0 & \(R_h\) & \(-0.02391887532526779349301220416248\) & \(-0.0125\)\\
\(1/2\) & P0 & \(R_k\) & \(-0.04820821508226710986072760568071\) & \(-0.0125\)\\
\(1/2\) & P0 & \(R_\mu\) & \(-0.03954387532526779349301220416248\) & \(-0.0125\)\\
\(1/2\) & P0 & \(R_\nu\) & \(-0.03954387532526779349301220416248\) & \(-0.0125\)\\
\(1/2\) & P1 & \(R_h\) & \(-0.02391887532526779349301220416248\) & \(-0.0125\)\\
\(1/2\) & P1 & \(R_k\) & \(0.20179178491773289013927239431928\) & \(-0.0125\)\\
\(7071/10000\) & P0 & \(R_h\) & \(-0.04322206410166229544037429665526\) & \(-0.0176775\)\\
\(7071/10000\) & P0 & \(R_k\) & \(-0.07144566327268729892016394092357\) & \(-0.0176775\)\\
\(7071/10000\) & P0 & \(R_\mu\) & \(-0.06531893910166229544037429665526\) & \(-0.0176775\)\\
\(7071/10000\) & P0 & \(R_\nu\) & \(-0.06531893910166229544037429665526\) & \(-0.0176775\)\\
\(7071/10000\) & P1 & \(R_h\) & \(0.02083177759556464749795711927522\) & \(-0.00387075577\ldots\)\\
\(7071/10000\) & P1 & \(R_k\) & \(-0.04426887899089171952349912499067\) & \(-0.00387075577\ldots\)\\
\(1\) & P0 & \(R_h\) & \(-0.04783775065053558698602440832496\) & \(-0.025\)\\
\(1\) & P0 & \(R_k\) & \(-0.07908775065053558698602440832496\) & \(-0.025\)\\
\(1\) & P0 & \(R_\mu\) & \(-0.07908775065053558698602440832496\) & \(-0.025\)\\
\(1\) & P0 & \(R_\nu\) & \(-0.07908775065053558698602440832496\) & \(-0.025\)\\
\(1\) & P1 & \(R_h\) & \(0.06111111111111111111111111111111\) & \(0.0\)\\
\(1\) & P1 & \(R_k\) & \(-0.10555555555555555555555555555555\) & \(0.0\)\\
\bottomrule
\end{longtable}
\end{center}

\paragraph{D04.} With the derived pairing P0 and the declared
regulator \(R_h\), the D04 part of the second \(\phi\phi\) value is
\(\tfrac{\kappa^2}{4\pi}M_{10}^{\rm P0,R_h}\), with a pole
\(-\tfrac{\kappa^2}{4\pi}\tfrac{\alpha}{40\epsilon}\). At the
certified angles,
\(M_{10}^{\rm P0,R_h}(\tfrac3{10})=0.11574388259581850963848639109243\)
and
\(M_{10}^{\rm P0,R_h}(\tfrac{7071}{10000})=-0.04322206410166229544037429665527\).
Status: \PARTIAL, for two reasons: the closed form is identified
rather than derived, and the regulator is a declared normalization.
Beyond that, the value is unique.

\paragraph{Closed forms (PSLQ, 40 digits, coefficients \(\le1600\);
\loc{code/v12/m10_pslq.py}).}
\[
M_{10}^{R_h}(1)=\frac{1-40\log2-20\gamma}{800},\qquad
M_{10}^{R_k}(1)=-\frac{6+10\log2+5\gamma}{200},
\]
\[
M_{10}^{R_h}(\tfrac12)=\frac{1-40\log2-20\gamma}{1600},\qquad
M_{10}^{R_k}(\tfrac12)=-\frac{6+15\log2+5\gamma}{400}.
\]
\(R_\mu=R_\nu=R_k\) at \(\alpha=1\), as Proposition~\ref{prop:v12ct}
requires. At \(\alpha=\tfrac3{10}\) and \(\tfrac{7071}{10000}\), no
relation was found in the bases \(\{1,\zeta(3),\zeta(3)/\pi^2,\pi^2,\log2\}\),
\(\{1,\zeta(3)/\pi^2,\log2,\zeta'(-1),\zeta'(-3),\gamma\}\) or
\(\{1,\zeta(3)/\pi^2,\pi^2,\log2,\log\alpha\}\). The large-coefficient
outputs are spurious. For rational \(\alpha=p/q\), the sum
\(\sum_kh_k^n\psi_1(h_k)\) is a Barnes-type double zeta value that
reduces to Hurwitz data at arguments \(j/(2p)\). That reduction was
not carried out.

%======================================================================
\section{v12: exact certificates at the chamber walls}
\label{sec:v12walls}
%======================================================================

\loc{code/v12/wall_certificate.py} is the v8 certificate program
with only the output path changed. It was run at
\(\alpha=\tfrac13,\tfrac12,\tfrac23\) on CR and F, in exact
\(\mathbb Q(\ii)\) arithmetic, with \(\mathrm{DM}=6\) and the
\(6\times6\) \(q\)-grid (\loc{logs/v12/walls/}).

\begin{center}
\small
\begin{longtable}{@{}llrrrrrr@{}}
\caption{Wall certificates: \(N\) = constrained basis size, rank = Ward-constrained count \(N-\dim Z\), image-column rank, \(\dim H^0\), seconds.}\label{tab:v12walls}\\
\toprule
\(\alpha\) & branch & \(N\) & rank & \(\dim Z\) & img. rank & \(\dim H^0\) & s\\
\midrule
\(1/2\) & CR & 42 & 33 & 9 & 40 & 2 & 687\\
\(1/2\) & F & 42 & 33 & 9 & 40 & 2 & 678\\
\(1/3\) & CR & 35 & 27 & 8 & 33 & 2 & 305\\
\(1/3\) & F & 35 & 27 & 8 & 33 & 2 & 305\\
\(2/3\) & CR & 35 & 27 & 8 & 33 & 2 & 304\\
\(2/3\) & F & 39 & 30 & 9 & 37 & 3 & 545\\
\bottomrule
\end{longtable}
\end{center}

\paragraph{Reading.} At each wall the constrained basis is larger
than in the adjacent chambers, because extra jets appear when Bessel
orders coincide. The rank rises with it: 27 at \(\tfrac13\) (CR, F)
and at \(\tfrac23\) (CR), against 26 in the chambers; 30 at
\(\tfrac23\) (F), against 29; and 33 at \(\tfrac12\) on both
branches. CR and F coincide exactly at \(\alpha=\tfrac12\), as they
must. The cocycle dimension \(\dim Z=N-\mathrm{rank}\) is 8 at
\(\tfrac13\) and at \(\tfrac23\) (CR), the \(N=34\) chamber value.
It is 9 at \(\tfrac23\) (F), the value for F above \(\tfrac12\), and
9 at \(\tfrac12\) on both branches. At every wall, \(B\subset Z\)
and \(B\) lies in the candidate span. The exceptional set in
\((\tfrac14,1)\) is therefore exactly the three walls, and the ranks
there are certified.

\paragraph{Certificate status.} The wall ranks are
\CPIR{} in this precise scope: the exact rank of the declared ST row
matrix of the frozen v8 engine (\(\mathrm{DM}=6\), \(6\times6\)
\(q\)-grid, exact \(\mathbb Q(\ii)\)) at \(\alpha=\tfrac13,\tfrac12,\tfrac23\)
on CR and F, as tabulated. The method is the v8 complete-evaluation
argument, unchanged.

%======================================================================
\section{v12: open items and precise obstructions}
\label{sec:v12open}
%======================================================================

\begin{enumerate}[label=(\alph*)]
\item \emph{Weight \(\pm1,\pm2\) traces and the 84 mixed weight-0
sums.} v12 wrote no weight-one evaluator. Proposition~\ref{prop:v12F}
supplies, for every channel pair \((s,t)\) under the derived same-\(k\)
pairing, the \emph{constant-density} part of a mixed weight-0 entry
in closed form. A low-cost path for the 84 entries therefore exists
for those with constant densities. Entries whose densities depend on
\(x\) need the moments \(\int x^nT_1\,\mathrm dx\), which are not
covered by the closed form. Neither the per-entry densities nor the
CR flip of channel~2 were tabulated, so no entry is evaluated. The
split of the weight-\(\pm1\) traces into a scheme-independent log
coefficient and a finite part is not computed. Status: \OPEN.
\item \emph{Dictionary, 20 Ward combinations, 34-/18-dim maps,
\(e_5,e_6\).} These depend on (a) and on the jet-key to \((C,n,j,i)\)
dictionary, which was not constructed. No residual was evaluated.
Status: \OPEN.
\item \emph{Transverse loop momentum (feasibility).} A minimal
construction needs three pieces. First, each catalogue vertex must
be exported as a polynomial in \(p_\perp\) up to the order needed:
degree 4 for the hh sector, by the N4 power counting. Second, the
\(p_\perp\) integrals in DRED (\(d-2\) transverse dimensions) reduce
these polynomials to \(\Gamma\)-function multiples of
\(s^{-(d-2)/2}\) times powers of the heat time. Third, the extra
powers of \(s\) shift the Mellin variable of scheme Z, so the
per-mode functionals change from \(F(\mu,\nu)\) to its Mellin-shifted
analogues. Nothing in this chain is beyond the present tools, but the
first step, a \(p_\perp\) export from the v8 slice engine, does not
exist. Status: \OPEN{} (feasible, not started).
\end{enumerate}

%======================================================================
\section{Changes to v11 statements}
\label{sec:v12changes}
%======================================================================

\begin{enumerate}[label=(C\arabic*)]
\item Section~\ref{sec:v11M10}: ``a power expansion without
logarithms is an ansatz'' is now a consequence of
Proposition~\ref{prop:v12F}, which is identified, not derived. The
fitted tail is replaced by the exact tail, and the quoted precision
of ``4--8 digits'' by 40 digits (Table~\ref{tab:v12M10}). The v11
values lie inside their error bars.
\item Section~\ref{sec:v11M10}: P1 was ``v11's implementation of the
frame intertwiner \(H_I\)''. This is withdrawn: \(H_I\) acts as the
identity on \(k\), and P1 is excluded (Proposition~\ref{prop:v12pair}).
\item Section~\ref{sec:v11M10}: ``the P0 pole coefficient equals
\(-\alpha/40\) (identified numerically, not derived)'' becomes
Corollary~\ref{cor:v12pole}, which is exact given the closed form.
\item Section~\ref{sec:v11M10}: ``\(M_{10}\) is not regulator
independent'' is sharpened. The regulators differ by the explicit
local counterterm of Proposition~\ref{prop:v12ct}.
\item Section~\ref{sec:v11M10}: ``PSLQ: no closed form'' is
superseded at \(\alpha=1,\tfrac12\) (Section~\ref{sec:v12reg}).
\item The certificate scope of Section~\ref{sec:v11D}, which excluded
the walls, is supplemented by Section~\ref{sec:v12walls}.
\end{enumerate}
No other v11 statement is changed.

\section{Open obligations and the exact release}
\label{sec:release}
%======================================================================

\begingroup\small
\begin{longtable}{@{}L{.22\linewidth}L{.40\linewidth}L{.31\linewidth}@{}}
\caption{Open obligations after v12. Each status is stated
exactly; a missing value is never zero.}
\label{tab:open}\\
\toprule
Obligation & v12 status (evidence) & What is still missing\\
\midrule
\endfirsthead
\toprule
Obligation & v12 status & Still missing\\
\midrule
\endhead
\bottomrule
\endfoot

v6 \(hh\) branch difference &
Restored (v9, Proposition~\ref{prop:v9TT}). &
Still conditional on Assumption~\ref{ass:reduction}.\\

\(\Delta J\) (TT branch difference) &
\textbf{Closed form, both modes derived} (CR v10; F v11 for
\(0<\mu<2\), Theorem~\ref{thm:v11IF}). &
A declared scheme for absolute TT finite parts.\\

Ghost normalizations &
\textbf{Derived (v10)} from tree operators
(Theorem~\ref{thm:v10lambda}). &
Independence from the v9 dictionary sign choice.\\

26 Ward combinations &
\PARTIAL. The 6 constant-profile combinations hold (v9), now
without fitted parameters. &
The 20 jet-level combinations: obstructions (a)--(c) of
Section~\ref{sec:v11open}. No residual is evaluated.\\

34-dimensional projection; 18-dim engine map &
\PARTIAL{} (12 of 34; 6 of 18). Same-channel weight-0 jet values are
now available (Proposition~\ref{prop:v10dressed}). &
Mixed and weight-\(\pm1,\pm2\) jet values; a jet-key to \(m=2\)
atom dictionary.\\

\(e_5,e_6\) separately &
\OPEN{} (excluded at constant profile, v9). &
Jet data along the kernel direction.\\

D04 part of second \(\phi\phi\) value &
\PARTIAL{} (v12). Pairing P0 derived (Proposition~\ref{prop:v12pair});
regulators differ by an explicit counterterm (Proposition~\ref{prop:v12ct});
\(M_{10}\) to 40 digits under the declared \(R_h\) (Table~\ref{tab:v12M10}). &
A derivation of the closed form of \(F\) (Proposition~\ref{prop:v12F}); a
derived (not declared) regulator normalization.\\

Transverse-jet catalogue &
\PARTIAL{} (v10, Section~\ref{sec:v10jet}); loop momentum on the
slice. &
Transverse loop momentum (no \(p_\perp\) jets exported); the 84 mixed
weight-0 mode sums; weight-\(\pm1\) residue/finite-part split
(\(w=\tfrac12\)); weight \(\pm2\). Mixed evaluator corrected in v11
(Section~\ref{sec:v11bug}); constant-density mixed weight-0 sums now have a
closed-form path (Proposition~\ref{prop:v12F}); transverse feasibility in
Section~\ref{sec:v12open}.\\

Exact \(\Span\) certificate &
\CPIR{} on the four points (v8). Chamber ranks 26 (six \(N=34\)
chambers) and 29 (F above \(\tfrac12\)) at every \(\alpha\):
\CPIR{} in the scope of Section~\ref{sec:v11D}; (D) machine-checked
(Proposition~\ref{prop:v11D}). Walls \(\tfrac13,\tfrac12,\tfrac23\): exact
ranks, \CPIR{} (Section~\ref{sec:v12walls}). &
Identification with the interacting basis; independent review.\\

Evanescent coefficients &
\CONDITIONAL{} on slice-DRED (v8); excluded from the frozen
\(\tilde g\) rule (v9). &
Transverse \(\tilde G\) components.\\

Vector-type keys &
Finite parts computed in a declared \(k\)-\(\zeta\) scheme
(Proposition~\ref{prop:v10vector}). &
A physical scheme; profile Ward partners.\\

Branch input &
CR and F remain inputs. &
A derived graviton operator.\\
\end{longtable}
\endgroup

In addition, graph identification still needs (R2), (R4), and
(R5), and the undisposed parts of (R3), for every relevant
graph. The D01 Mellin prescription obligation is explicit in
Corollary~\ref{cor:D01}.
The general joint-apex and \(\Hreg\) statements remain open.

\begin{theorem}[Scoped v12 release]
\label{thm:release}
Under the displayed definitions and source qualifications:
\begin{enumerate}[label=\textup{(\roman*)}]
\item The normalized heat coefficients \(\Ccal_b^\Fr\) and
\(\Dcal_b\), their specified singular-channel shifts, and the
orbifold identities are established.

\item A complete graph satisfying
Assumption~\ref{ass:reduction} has the apex residue
\eqref{eq:graphreduction}.

\item The D01 composite has the local obstruction
\eqref{eq:obstruction}; its specified Mellin continuation is
\eqref{eq:Mclosed} and is exactly zero at \(z=0\).
The Mellin value does not select the graph prescription.

\item The declared coefficient catalogue satisfies the ghost
identity and the constant-profile \(hh/Kc\) row.
With the stated tree input,
\(\lambda=-2\ii\) is fixed before loop fitting.
The complete lattice branch difference is
\eqref{eq:branchdifference}.

\item These rows do not determine the full interacting basis,
a physical tensor-core limit, or a complete Slavnov--Taylor
theorem on the cone.

\item The continuum/preparation/instrument objects in
\eqref{eq:pipeline} remain distinct.
The record-only recovery bound applies to a separately declared
finite interface, not to a continuum residue.
The spectral heat coefficients are not thermodynamic heat.

\item Image and echo-window functionals are not identified with
the apex residue; Example~\ref{ex:image} is an explicit negative
control.

\item None of the obligations in Table~\ref{tab:open} is set to
zero or discharged by the companion comparisons.

\item \emph{(v8, executed.)} Theorem~\ref{thm:v8cert} holds on its
stated finite scope; on that scope the status is \CPIR.

\item \emph{(v8, executed, conditional on
Assumption~\ref{ass:reduction}.)} The \(hh\) catalogue sum satisfies
the two \(hh\)-supported Ward combinations exactly
(Proposition~\ref{prop:v8hhward}); four of the six exact values are
fixed in the stated basis; six of the eighteen \(hh\) engine
coordinates are mapped.

\item \emph{(v8, amended in v9.)} The hidden D01 mixed-channel term of
Proposition~\ref{prop:v8hidden} is pure TT and vanishes on the lattice
(Proposition~\ref{prop:v9TT}); the v6 \(hh\) branch difference is
restored, subject only to Assumption~\ref{ass:reduction}.
The evanescent coefficients of Section~\ref{sec:v8evan} are
conditional on the declared prescription and are excluded from the
frozen \(\tilde g\) rule (Proposition~\ref{prop:v9evan}).

\item \emph{(v9, executed.)} The mixed trace has the closed form of
Proposition~\ref{prop:v9g}; the TT branch difference \(\Delta J\) is
finite and computed (Table~\ref{tab:v9dJ}); absolute TT values carry
the IR logarithm of Proposition~\ref{prop:v9log}.

\item \emph{(v9, executed, exact on four points.)} The CB/CK
dictionaries and all six constant-profile Ward functionals hold with
two fitted \(\alpha\)-independent normalizations; \(e_5,e_6\) are not
separately determined at constant profile
(Theorem~\ref{thm:v9ghost}).

\item \emph{(v9.)} The generic ST rank over \(\mathbb Q(\alpha)\) is as
in Theorem~\ref{thm:v9generic} on the chambers of
Proposition~\ref{prop:v9walls}; the exceptional set inside chambers is
open. \emph{(Amended in v10: empty, conditional on (D); item (xviii).)}

\item \emph{(v10, executed.)} \(\Delta J\) has the closed form of
Corollary~\ref{cor:v10dJ}: the CR mode is derived, and the F mode is
identified by PSLQ and verified, not derived.

\item \emph{(v10, executed.)} \(\lambda_{\rm CK}=-8\) and
\(\lambda_{\rm CB}=-4\ii\) follow from tree operator ratios and the
v9 dictionaries (Theorem~\ref{thm:v10lambda}).

\item \emph{(v10; unconditional since v11.)} The rank is constant inside
chambers: 26 on \(N=34\) chambers and 29 on F above \(\tfrac12\)
(Theorem~\ref{thm:v10chamber}; (D) is Proposition~\ref{prop:v11D}).

\item \emph{(v10, partial.)} The slice jet catalogue, the scheme-Z
\(r^2\) moments, and the vector-key finite parts in a declared
scheme are as in Sections~\ref{sec:v10jet} and
\ref{sec:v10vector}. The D04 part, the jet Ward combinations and
\(e_5,e_6\) remain open.

\item \emph{(v11, executed.)} \(I_{\rm F}\) is derived for \(0<\mu<2\)
(Theorem~\ref{thm:v11IF}), and \(\Delta J=-(\mu-1)B/W_1\).

\item \emph{(v11, executed.)} The degree bound (D) is machine-checked
(Proposition~\ref{prop:v11D}). The chamber-rank certificate is
\CPIR{} in the scope stated in Section~\ref{sec:v11D}, and the v8
four-point certificate remains \CPIR{} in its v8 scope.

\item \emph{(v11, partial.)} The mixed-trace evaluator is corrected
(Section~\ref{sec:v11bug}). \(M_{10}\) is evaluated in declared
pairings and regulators (Table~\ref{tab:v11M10}), so D04 is
scheme dependent. The weight \(\pm1,\pm2\) traces, the
dictionary, the 20 jet Ward combinations, the projections,
\(e_5,e_6\) and the transverse loop momentum remain open
(Section~\ref{sec:v11open}).

\item \emph{(v12, identified and verified.)} The closed form of
\(F(\mu,\nu)\) (Proposition~\ref{prop:v12F}) and its consequences,
Corollary~\ref{cor:v12pole} and the 40-digit \(M_{10}\)
(Table~\ref{tab:v12M10}).

\item \emph{(v12, derived under stated assumptions.)} The D04
pairing is P0 (Proposition~\ref{prop:v12pair}). The regulators differ
by the local counterterm of Proposition~\ref{prop:v12ct}. D04 is
\PARTIAL{} only because the closed form is identified and the
regulator \(R_h\) is a declared normalization.

\item \emph{(v12, executed.)} The wall ranks of
Table~\ref{tab:v12walls} are \CPIR{} in the scope of
Section~\ref{sec:v12walls}.
\end{enumerate}
The release status is
\[
 \boxed{
 \text{O.1 on the cone}=\NOTCLOSED,\qquad
 \texttt{bulk\_symmetry\_restored}=\texttt{false}.
 }
\]
\end{theorem}

\begin{proof}
Items (i)--(iii) are Theorems~\ref{thm:C},
\ref{thm:flip}, \ref{thm:D}, \ref{thm:orbifold},
\ref{thm:reduction}, \ref{thm:obstruction}, and
\ref{thm:Mellin}, with Proposition~\ref{prop:Dflip} and
Corollary~\ref{cor:D01}.
Item (iv) is Theorems~\ref{thm:ghost} and \ref{thm:ST},
Input~\ref{input:tree}, and
Proposition~\ref{prop:decomposition}.

Proposition~\ref{prop:noninjective}, the identity classification
in Section~\ref{sec:ST}, and the explicit missing data prove
the limitation in (v).
Proposition~\ref{prop:instrument},
Theorem~\ref{thm:classical}, and the separate hypotheses of
\eqref{eq:thermal} establish (vi).
Example~\ref{ex:image} and the distinct definitions establish
(vii). Table~\ref{tab:open} records the remaining obligations.

Items (xix)--(xxi) are Theorem~\ref{thm:v11IF},
Proposition~\ref{prop:v11D} and Section~\ref{sec:v11M10}.
Items (xv)--(xviii) are Corollary~\ref{cor:v10dJ},
Theorems~\ref{thm:v10lambda}, \ref{thm:v10chamber} and
Propositions~\ref{prop:v10Z}, \ref{prop:v10dressed},
\ref{prop:v10vector}.
Items (xii)--(xiv) are Propositions~\ref{prop:v9g},
\ref{prop:v9log}, \ref{prop:v9walls} and Theorems~\ref{thm:v9ghost},
\ref{thm:v9generic}.
Items (ix)--(xi) are Theorem~\ref{thm:v8cert},
Propositions~\ref{prop:v8hhward} and \ref{prop:v8hidden}, and the
executed records of Sections~\ref{sec:v8map} and \ref{sec:v8evan}.

The complete O.1 requirements are not consequences of these
restricted statements. The status is non-closure, not a
universal no-go theorem for a future completion.
\end{proof}

\scopebox{Submission boundary}{%
v8 adds executed evidence, and every executed statement is listed
with its hash in Appendix~\ref{app:v8artifacts}. The listed runs
are the replay, the exact certificates, the maps, the slice and
evanescent runs, and the Duhamel corroboration. v8 was compiled
with pdfLaTeX. It does not claim a proof-assistant check, an
independent referee review, physical acquisition, a new graph
evaluation beyond the declared catalogue, or a general-\(\alpha\)
theorem from finitely many exact points. Nothing was pushed or
posted.

The exact rank statement is \CPIR{} only on the scope of
Theorem~\ref{thm:v8cert}. Every other new item is \PARTIAL,
\CONDITIONAL, or \OPEN, as listed in Table~\ref{tab:open}. The
present conditional statements must not be relabelled as
unconditional graph theorems.

v9 adds the items of Section~\ref{sec:whatv9}, each hashed in
Appendix~\ref{app:v9artifacts}. It claims no closed form for
\(\Delta J\) or for the D04 part, no derivation of the two ghost
normalizations, no proof that \(\pi_KB\) is independent of \(\alpha\)
beyond the four points, and no identification of the exceptional set
inside chambers. No independent review was started; nothing was
pushed, posted or written to the site.

v10 adds the items of Section~\ref{sec:whatv10}, each hashed in
Appendix~\ref{app:v10artifacts}. It does not claim a derivation of
the F mode, a machine check of the degree bound (D), any jet-level
Ward residual, the D04 value, or a physical scheme for the vector
finite parts. No scope is labelled
\texttt{CLOSED\_PENDING\_INDEPENDENT\_REVIEW}. No independent
review was started; nothing was pushed, posted or written to the
site.

v11 adds the items of Section~\ref{sec:whatv11}, each hashed in
Appendix~\ref{app:v11artifacts}, and withdraws the v10 \(\varphi\)
values (Section~\ref{sec:v11changes}). Two scopes carry
\texttt{CLOSED\_PENDING\_INDEPENDENT\_REVIEW}: the v8 four-point
rank certificate and the chamber ranks of Section~\ref{sec:v11D}.
Nothing else does. v11 claims no derived pairing or regulator for
D04, no weight-\(\pm1\) log coefficient, no jet Ward residual, and
no value of \(e_5,e_6\). No independent review was started; nothing
was pushed, posted or written to the site.

v12 adds the items of Section~\ref{sec:whatv12}, each hashed in
Appendix~\ref{app:v12artifacts}. Three scopes carry
\texttt{CLOSED\_PENDING\_INDEPENDENT\_REVIEW}: the v8 four-point
rank certificate, the chamber ranks of Section~\ref{sec:v11D}, and
the wall ranks of Section~\ref{sec:v12walls}. Nothing else does.
v12 claims no derivation of the closed form of \(F\), no derived
regulator, no weight-\(\pm1\) log coefficient, no jet Ward residual,
and no value of \(e_5,e_6\). No independent review was started;
nothing was pushed, posted or written to the site.}

%======================================================================
\appendix
\section{Embedded verification and replay program}
\label{app:code}
%======================================================================

\subsection{Use and acceptance conditions}

Save the listing below as \loc{cone_v7.py}. It uses Python and
SymPy; numerical quadrature is not required for its exact checks.
The listing is ASCII Python. It is included directly with
\loc{listings}, so typesetting needs neither shell escape nor an
external code file.

Example commands:
\begin{verbatim}
python3 cone_v7.py selftest

# New logs with explicit coverage:
python3 cone_v7.py logs /path/to/logs

# Historical nonzero-only logs: comparison, not coverage proof:
python3 cone_v7.py logs /path/to/logs --legacy

# Optional replay; frozen engine is an external required input:
python3 cone_v7.py replay D01 CR \
  --engine /path/to/XVI-P/code/p5_entry --out /path/to/newlogs

python3 cone_v7.py tree \
  --engine /path/to/XVI-P/code/p5_entry

# Redirect outside the hashed tree:
python3 cone_v7.py manifest /path/to/archive > manifest.json
\end{verbatim}

\paragraph{What \loc{selftest} means.}
It checks the displayed finite algebraic identities. It does
not independently derive the catalogue from vertices.

\paragraph{What \loc{logs} means.}
Without \loc{--legacy}, every compared coordinate must have an
explicit evaluated-key status. With \loc{--legacy}, missing zero
entries are reported as coverage-uncertified. A successful legacy
comparison is not an execution or graph certificate.

\paragraph{What \loc{replay} means.}
It rebuilds the historical zero-transverse-momentum algebraic
reduction from a supplied frozen engine. It verifies the coarse
projector assignment, records all mixed pairs, and excludes an
entire affected physical-branch key.
For UNIT it retains all channel pairs.
The output is explicitly marked
\loc{ALGEBRAIC_SLICE_ONLY}; it does not establish (R2)--(R5).
The graph prefactors are declared source conventions, not a
new automatic authentication of Q02.

\paragraph{What is not embedded.}
The frozen \loc{amplitude_rules.py}, \loc{feynman.py},
\loc{dred.py}, vertex arrays, Q02 record, and historical JSON
logs are not reproduced here. Missing dependencies fail rather
than being replaced by zero.

% BEGIN EMBEDDED PYTHON cone_v7.py
\begin{lstlisting}[caption={Consolidated exact checker and coverage-aware optional replay.},
                  label={lst:program}]
#!/usr/bin/env python3
"""Scoped cone v7 checker.

Status:
  Source supplied for review; no execution is asserted by the paper.

selftest:
  Exact identities of the printed formulas, not a vertex rebuild.

logs:
  Compare explicit retained coordinates. New logs require coverage.
  --legacy permits polynomial comparison with uncertified zero coverage.

replay:
  Frozen-engine zero-transverse-momentum algebra only.
  Entire physical-branch keys with mixed poles are quarantined.
  UNIT retains all pairs. This is NOT an R2--R5 certificate.

tree:
  Check the declared tree normalization using a supplied engine.

manifest:
  Byte hashes with explicit path serialization. This is not the
  historical find/sort/xargs tree-hash algorithm.
"""
import argparse
import ast
import hashlib
import importlib
import json
import sys
import time
from pathlib import Path

import sympy as sp

I = sp.I
Q = sp.Rational
alpha = sp.Symbol("alpha", positive=True)
kappa, q0, q3 = sp.symbols("kappa q0 q3")
c0, c1, c2 = sp.symbols("c0 c1 c2")
P2 = q0**2 - q3**2
ETA = (1, -1, -1, -1)
HC = ((0, 0), (0, 1), (0, 2), (0, 3), (1, 1),
      (1, 2), (1, 3), (2, 2), (2, 3), (3, 3))
LAT = (0, 3, 9, 4, 7)
DIAGNOSTIC = (0, 3, 9, 4, 7, 5)
GRAPHS = ("D01", "D02", "D03", "D04", "D05a",
          "D06", "D11", "D12", "D13")


def zero(expr, label):
    v = sp.factor(sp.cancel(sp.expand(expr)))
    if v != 0:
        raise AssertionError(label + ": " + str(v))


def no_duplicate_pairs(pairs):
    out = {}
    for k, v in pairs:
        if k in out:
            raise ValueError("duplicate JSON key: " + str(k))
        out[k] = v
    return out


def read_json(path):
    with Path(path).open(encoding="utf-8") as f:
        return json.load(f, object_pairs_hook=no_duplicate_pairs)


def write_json(path, data):
    path = Path(path)
    path.parent.mkdir(parents=True, exist_ok=True)
    with path.open("w", encoding="utf-8") as f:
        json.dump(data, f, indent=2, sort_keys=True)
        f.write("\n")


def key_parse(text):
    obj = ast.literal_eval(text)
    if not (isinstance(obj, tuple) and len(obj) == 2
            and all(type(x) is int and 0 <= x < 10 for x in obj)):
        raise ValueError("invalid component key: " + text)
    return obj


# Restricted expression parser: no eval, attributes, function calls,
# comprehensions, or implicit imports. Serialized logs are arithmetic.
NAMES = {"alpha": alpha, "kappa": kappa, "q0": q0, "q3": q3,
         "I": I, "pi": sp.pi}


def expression(text):
    if len(text) > 100000:
        raise ValueError("expression exceeds declared input bound")

    def walk(n):
        if isinstance(n, ast.Expression):
            return walk(n.body)
        if isinstance(n, ast.Constant) and type(n.value) is int:
            return sp.Integer(n.value)
        if isinstance(n, ast.Name) and n.id in NAMES:
            return NAMES[n.id]
        if isinstance(n, ast.UnaryOp):
            if isinstance(n.op, ast.UAdd):
                return walk(n.operand)
            if isinstance(n.op, ast.USub):
                return -walk(n.operand)
        if isinstance(n, ast.BinOp):
            a, b = walk(n.left), walk(n.right)
            if isinstance(n.op, ast.Add):
                return a + b
            if isinstance(n.op, ast.Sub):
                return a - b
            if isinstance(n.op, ast.Mult):
                return a * b
            if isinstance(n.op, ast.Div):
                return a / b
            if isinstance(n.op, ast.Pow):
                if not (b.is_Integer and abs(int(b)) <= 100):
                    raise ValueError("undeclared power")
                return a**b
        raise ValueError("unsupported symbolic syntax")

    return sp.expand(walk(ast.parse(text, mode="eval")))


def C(a, b):
    return (1/a - a)/12 - b*(1-b)/(2*a)


def D(a, b):
    B2 = b*b - b + Q(1, 6)
    B4 = b**4 - 2*b**3 + b*b - Q(1, 30)
    return -11*a/180 + B2/(3*a) - B4/(6*a**3)


def physical_cs(branch):
    a = alpha
    z0 = (1-a*a)/(12*a)
    z1 = (5*a*a-6*a+1)/(12*a)
    if branch == "CR":
        z2 = (23*a*a-12*a+1)/(12*a)
    elif branch == "F":
        z2 = (23*a*a-36*a+13)/(12*a)
    elif branch == "UNIT":
        return (sp.Integer(1),)*3
    else:
        raise ValueError("unknown branch")
    return z0, z1, z2


def put(out, key, val):
    val = sp.expand(val)
    if key in out:
        zero(out[key]-val, "inconsistent duplicate " + str(key))
    out[key] = val


def symmetric(d):
    out = {}
    for (a, b), v in d.items():
        put(out, (a, b), v)
        put(out, (b, a), v)
    return out


def tensor_images(d):
    # Exchange transverse coordinates 1 and 2, then transpose slots.
    exch = {1: 2, 2: 1, 4: 7, 7: 4, 6: 8, 8: 6}
    out = dict(d)
    for (a, b), v in d.items():
        put(out, (exch.get(a, a), exch.get(b, b)), v)
    return symmetric(out)


def catalogue(cs):
    z0, z1, z2 = cs
    A, B = z0+2*z1+4*z2, -z0+2*z1
    R, T, V = 3*z0+4*z1, 4*z0+3*z1, z1-z0
    W = (z0+2*z1+z2)/2
    k16 = kappa**2/(16*sp.pi)
    out = {}
    out["D04"] = {(0, 0): -kappa**2*z0*P2/(4*sp.pi)}
    out["D05a"] = {(0, 0): -kappa**2*(z2-z0)*P2/(4*sp.pi)}
    out["D11"] = symmetric({
        (4, 4): z0*P2*k16/2,
        (4, 7): z0*P2*k16/2,
        (7, 7): z0*P2*k16/2})
    d03 = {
        (0, 0): -z0*P2, (0, 9): -z0*P2, (9, 9): -z0*P2,
        (3, 3): 4*z0*P2,
        (0, 4): z0*(q0**2-3*q3**2),
        (3, 4): 4*z0*q0*q3,
        (4, 9): -z0*(3*q0**2-q3**2),
        (4, 4): -(z0+4*z1)*P2,
        (4, 7): -z0*P2, (5, 5): -8*z1*P2}
    d01 = {
        (0, 0): z0*P2, (0, 9): z0*P2, (9, 9): z0*P2,
        (3, 3): -4*z0*P2,
        (4, 4): (z0-4*z1+z2)*P2,
        (4, 7): (z0+z2)*P2,
        (3, 4): -4*(z0+2*z2)*q0*q3,
        (4, 9): A*q0**2-B*q3**2,
        (0, 4): -B*q0**2+A*q3**2,
        (5, 5): -8*z1*P2}
    out["D03"] = {k: v*k16 for k, v in tensor_images(d03).items()}
    out["D01"] = {k: v*k16 for k, v in tensor_images(d01).items()}
    d02 = {
        (1, 1): q3**2, (2, 2): q3**2, (5, 5): P2,
        (6, 6): q0**2, (8, 8): q0**2,
        (1, 6): -q0*q3, (2, 8): -q0*q3,
        (0, 4): -q3**2/2, (0, 7): -q3**2/2,
        (3, 4): q0*q3, (3, 7): q0*q3,
        (4, 7): -P2/2,
        (4, 9): -q0**2/2, (7, 9): -q0**2/2}
    out["D02"] = {k: kappa**2*W*v/sp.pi
                  for k, v in symmetric(d02).items()}
    out["D12"] = {}
    out["D06"] = {
        (0, 0): I*kappa**2*R*P2/(8*sp.pi),
        (3, 3): I*kappa**2*R*P2/(8*sp.pi)}
    out["D13"] = {
        (0, 0): -2*T*q0, (0, 3): 2*V*q3,
        (3, 0): -R*q3, (3, 3): R*q0,
        (9, 0): -2*V*q0, (9, 3): 2*T*q3}
    out["D13"] = {k: I*kappa**2*v/(16*sp.pi)
                  for k, v in out["D13"].items()}
    return out


def invariant_data():
    x = sp.symbols("x00 x03 x33 xf")
    y = sp.symbols("y00 y03 y33 yf")

    def tr(v):
        return v[0]-v[2]

    def php(v):
        return q0**2*v[0]-2*q0*q3*v[1]+q3**2*v[2]

    dot = x[0]*y[0]-2*x[1]*y[1]+x[2]*y[2]
    S = (P2*tr(x)*tr(y), P2*dot,
         php(x)*tr(y)+tr(x)*php(y),
         P2*(x[3]*tr(y)+tr(x)*y[3]),
         php(x)*y[3]+x[3]*php(y), P2*x[3]*y[3])
    idx = {0: 0, 3: 1, 9: 2, 4: 3, 7: 3}
    return x, y, S, idx


def structure_table(cs):
    a, b, c = cs
    W = (a+2*b+c)/2
    return {
        "D01": (-a/16, a/8, 0, -a/8-b/4-c/2,
                a/4+c/2, a/4-b/2+c/4),
        "D02": (0, 0, 0, W, -W, -W),
        "D03": (a/16, -a/8, 0, 3*a/8, -a/4, -a/4-b/2),
        "D11": (0, 0, 0, 0, 0, a/8),
        "D12": (0, 0, 0, 0, 0, 0)}


def bilinear(d, x, y, idx):
    return sp.expand(sum(v*x[idx[a]]*y[idx[b]]
                         for (a, b), v in d.items()
                         if a in idx and b in idx))


def einstein_coordinates(hv):
    p = (q0, 0, 0, q3)
    pu = tuple(ETA[i]*p[i] for i in range(4))
    h = sp.zeros(4)
    for a, (m, n) in enumerate(HC):
        h[m, n] = h[n, m] = hv.get(a, 0)
    tr = sum(ETA[i]*h[i, i] for i in range(4))
    ph = [sum(pu[r]*h[r, n] for r in range(4)) for n in range(4)]
    pph = sum(pu[n]*ph[n] for n in range(4))
    E = sp.Matrix(4, 4, lambda m, n:
        P2*h[m, n]-p[m]*ph[n]-p[n]*ph[m]+p[m]*p[n]*tr
        +(ETA[m] if m == n else 0)*(pph-P2*tr))
    for n in range(4):
        zero(sum(pu[m]*E[m, n] for m in range(4)), "E transverse")
    return {a: (2-int(m == n))*ETA[m]*ETA[n]*E[m, n]
            for a, (m, n) in enumerate(HC)}


def check_catalogue_algebra():
    cat = catalogue((c0, c1, c2))
    x, y, S, idx = invariant_data()
    table = structure_table((c0, c1, c2))
    for g, coeff in table.items():
        lhs = bilinear(cat[g], x, y, idx)
        rhs = kappa**2/sp.pi*sum(a*b for a, b in zip(coeff, S))
        zero(lhs-rhs, "structure " + g)

    variables = x+y+(q0, q3)
    polys = [sp.Poly(sp.expand(s), *variables) for s in S]
    mons = sorted(set(m for p in polys for m in p.monoms()))
    matrix = sp.Matrix([[p.coeff_monomial(m) for p in polys]
                        for m in mons])
    if matrix.rank() != 6:
        raise AssertionError("six-structure independence failed")

    h00, h03, h33, f = sp.symbols("h00 h03 h33 f")
    g0, g3 = sp.symbols("ghost0 ghost3")
    hv = {0: h00, 3: h03, 9: h33, 4: f, 7: f}
    dh = {0: 2*q0*g0, 3: -q0*g3+q3*g0,
          9: -2*q3*g3, 4: 0, 7: 0}
    E = einstein_coordinates(hv)
    bterm = 0
    for g in ("D01", "D02", "D03", "D11", "D12"):
        bterm += sum(v*dh.get(a, 0)*hv.get(b, 0)
                     for (a, b), v in cat[g].items())
    gterm = sum(v*E[a]*({0: g0, 3: g3}[n])
                for (a, n), v in cat["D13"].items())
    core = kappa**2*(c1-c0)*P2*f*(q0*g0+q3*g3)/sp.pi
    zero(bterm+core/2, "hh gauge variation")
    zero(gterm-I*core/4, "antifield contraction")
    lam = 1/(I*Q(1, 2))  # Declared tree input N=1, A_GH=1/2.
    zero(lam+2*I, "a-priori lambda arithmetic")
    zero(bterm+lam*gterm, "ST row")

    # Ghost equation, separately for each upper ghost column.
    for n in (0, 3):
        Y00 = cat["D13"].get((0, n), 0)
        Y03 = cat["D13"].get((3, n), 0)
        Y33 = cat["D13"].get((9, n), 0)
        F0 = q0*(Y00+Y33)/2-q3*Y03
        F3 = q0*Y03-q3*(Y00+Y33)/2
        zero(cat["D06"].get((0, n), 0)+2*F0, "ghost row 0")
        zero(cat["D06"].get((3, n), 0)-2*F3, "ghost row 3")

    # Endpoint and branch checks are polynomial checks only.
    for branch in ("CR", "F"):
        vals = catalogue(physical_cs(branch))
        if branch == "F":
            for d in vals.values():
                for v in d.values():
                    zero(v.subs(alpha, 1), "F flat catalogue")
    cr = catalogue(physical_cs("CR"))
    fr = catalogue(physical_cs("F"))
    for g in GRAPHS:
        for key in set(cr[g]) | set(fr[g]):
            zero((cr[g].get(key, 0)-fr[g].get(key, 0))
                 .subs(alpha, Q(1, 2)), "CR=F at half")
    bhcr = sum((bilinear(cr[g], x, y, idx)
                for g in table), sp.Integer(0))
    bhf = sum((bilinear(fr[g], x, y, idx)
               for g in table), sp.Integer(0))
    nu = 2-1/alpha
    zero(bhcr-bhf+nu*kappa**2*P2*x[3]*y[3]/(4*sp.pi),
         "complete branch difference")

    if len(cat["D02"]) != 23:
        raise AssertionError("D02 entry count")
    zero(sum(table[g][0] for g in table), "longitudinal S1")
    zero(sum(table[g][1] for g in table), "longitudinal S2")
    return {"catalogue_algebra": "PASS",
            "tree_normalization": "DECLARED_INPUT_NOT_REBUILT"}


def check_heat_and_mellin():
    u = sp.Symbol("u", positive=True)
    b = sp.Symbol("b")
    T = sp.cosh((Q(1, 2)-b)*u/alpha) / (
        2*sp.sinh(u)*sp.sinh(u/(2*alpha)))
    bulk = alpha/(2*(sp.cosh(u)-1))
    cser = sp.series(T-bulk, u, 0, 2).removeO()
    zero(sp.expand(cser).coeff(u, 0)-C(alpha, b), "C coefficient")

    weighted = -2*sp.diff(T, u)/sp.sinh(u)
    wbulk = alpha/(sp.cosh(u)-1)**2
    wser = sp.expand(sp.series(weighted-wbulk, u, 0, 2).removeO())
    for k in range(1, 5):
        zero(wser.coeff(u, -k), "weighted pole cancellation")
    zero(wser.coeff(u, 0)-D(alpha, b), "D coefficient")
    nu = 2-1/alpha
    zero(D(alpha, 2*alpha)-D(alpha, 2*alpha-1)
         -2*nu*(1-nu**2)/3, "weighted flip")
    zero(D(alpha, 2*alpha).subs(alpha, 1), "weighted CR flat")

    # Exact finite checks in Q[zeta_N], not numerical nsimplify.
    x = sp.Symbol("x")
    for N in range(2, 9):
        cyclo = sp.cyclotomic_poly(N, x)
        for j in range(N):
            acc = 0
            for k in range(1, N):
                inv = sp.invert((x**k-1)**4, cyclo, x)
                phase = x**(((j+2)*k) % N) + x**(((2-j)*k) % N)
                acc = sp.rem(acc+Q(2, N)*phase*inv, cyclo, x)
            target = D(Q(1, N), Q(j, N))
            zero(sp.rem(acc-target, cyclo, x), "cyclotomic moment")

    # Finite rational algebra underlying the closed Mellin proof.
    t, n = sp.symbols("t nu")
    ratio = (t*t-n*n)/(2*t+1)
    bracket = (t+1)*ratio+n*(n+t)
    zero(bracket-t*(t+n)*(t+n+1)/(2*t+1),
         "Mellin recurrence reduction")
    # Gamma reflection gives M'_z(0)=-(nu+1)/(2s).
    return {"heat_moments": "PASS",
            "finite_cyclotomic_checks": "PASS",
            "Mellin_rational_reduction": "PASS",
            "Mellin_integral_and_continuation":
                "ANALYTICAL_PROOF_IN_PAPER"}


def retained_keys(g):
    if g in ("D04", "D05a"):
        return [(0, 0)]
    if g == "D06":
        return [(a, b) for a in (0, 3) for b in (0, 3)]
    if g == "D13":
        return [(a, b) for a in LAT for b in (0, 3)]
    if g in ("D01", "D03"):
        return [(a, b) for a in DIAGNOSTIC for b in DIAGNOSTIC]
    return [(a, b) for a in range(10) for b in range(10)]


def compare_logs(root, legacy):
    root = Path(root)
    report = {}
    for branch in ("CR", "F"):
        expected = catalogue(physical_cs(branch))
        for g in GRAPHS:
            path = root/("spec_" + g + "_" + branch + ".json")
            data = read_json(path)  # Missing files are errors.
            if data.get("graph") != g or data.get("branch") != branch:
                raise ValueError("graph/branch metadata mismatch: " + str(path))
            got = {key_parse(k): expression(v)
                   for k, v in data[
                       "Gamma_pole_coeff_of_1/(4-d)"].items()}
            status = data.get("key_status", {})
            unverified = []
            for key in retained_keys(g):
                skey = str(key)
                if not legacy:
                    st = status.get(skey)
                    if st not in ("COMPUTED_ZERO_SLICE",
                                  "COMPUTED_NONZERO_SLICE"):
                        raise ValueError("missing/excluded coverage "
                                         + str(path) + " " + skey)
                    if st == "COMPUTED_NONZERO_SLICE" and key not in got:
                        raise ValueError("nonzero status without value")
                    if st == "COMPUTED_ZERO_SLICE" and key in got:
                        zero(got[key], "zero status with nonzero value")
                elif key not in got:
                    unverified.append(skey)
                zero(got.get(key, 0)-expected[g].get(key, 0),
                     "catalogue mismatch " + g + " " + skey)
            report[g+"_"+branch] = {
                "polynomial_comparison": "PASS",
                "coverage": ("LEGACY_UNCERTIFIED" if legacy
                             else "EXPLICIT_SLICE_COVERAGE"),
                "absent_legacy_keys": unverified,
                "graph_identification": "R2_R5_NOT_CERTIFIED"}
    return report


def load_engine(folder):
    root = Path(folder).resolve(strict=True)
    for name in ("amplitude_rules.py", "feynman.py", "dred.py"):
        if not (root/name).is_file():
            raise FileNotFoundError(root/name)
    sys.path.insert(0, str(root))
    AR = importlib.import_module("amplitude_rules")
    if Path(AR.__file__).resolve().parent != root:
        raise RuntimeError("import resolved outside declared engine")
    return AR


def projector_check(AR):
    Pi = AR.pinv_sparse()
    M = sp.Matrix(10, 10, lambda a, b: sp.sympify(Pi.get((a, b), 0)))
    P0 = sp.zeros(10)
    for a in (0, 3, 9):
        P0[a, a] = 1
    for a in (4, 7):
        for b in (4, 7):
            P0[a, b] = Q(1, 2)
    P1 = sp.zeros(10)
    for a in (1, 2, 6, 8):
        P1[a, a] = 1
    P2m = sp.eye(10)-P0-P1
    projectors = (P0, P1, P2m)

    def ctype(a):
        m, n = HC[a]
        count = int(m in (0, 3))+int(n in (0, 3))
        return {2: 0, 1: 1, 0: 2}[count]

    def channel(a, b):
        ca, cb = ctype(a), ctype(b)
        if ca == cb == 1:
            return 1
        if ca == cb == 2:
            return 2
        return 0

    for s, P in enumerate(projectors):
        zero_matrix = P*P-P
        if any(sp.simplify(v) != 0 for v in zero_matrix):
            raise AssertionError("projector is not idempotent")
        PM = P*M
        for a in range(10):
            for b in range(10):
                want = M[a, b] if channel(a, b) == s else 0
                zero(PM[a, b]-want, "coarse helicity marker")
    return Pi, channel


def laurent_coefficient(expr, x, power):
    expr = sp.cancel(expr)
    num, den = sp.fraction(expr)
    dp = sp.Poly(sp.expand(den), x)
    if len(dp.terms()) != 1:
        raise ValueError("not a Laurent polynomial")
    (deg,), cf = dp.terms()[0]
    np = sp.Poly(sp.expand(num), x)
    return sp.expand(np.coeff_monomial(x**(deg+power))/cf) \
        if deg+power >= 0 else sp.Integer(0)


def replay(AR, graph, branch, outdir):
    """Historical flat momentum-slice algebra, with honest coverage."""
    start = time.time()
    F = AR.F
    if tuple(F.HC) != HC:
        raise ValueError("independent-coordinate order mismatch")
    Pi, channel = projector_check(AR)
    Gh = AR.ghost_num()
    V = F.vertices()
    kap = F.KAPPA
    LS, PS = sp.symbols("l0:4"), sp.symbols("p0:4")
    names = {**{r: LS[r] for r in range(4)},
             **{4+r: PS[r] for r in range(4)}}
    u, v = sp.symbols("u0:3"), sp.symbols("v0:3")
    slice_sub = {LS[1]: 0, LS[2]: 0, PS[1]: 0, PS[2]: 0,
                 PS[0]: q0, PS[3]: q3}

    def route(vt, routing):
        out = {}
        for key, poly in F.vt_route(vt, routing).items():
            e = poly.to_sympy(names) if hasattr(poly, "to_sympy") else poly
            # Polynomial substitution commutes with the finite contractions.
            out[key] = sp.expand(sp.sympify(e).subs(slice_sub))
        return out

    def pw(a, b, marker):
        return sp.sympify(Pi[(a, b)])*marker[channel(a, b)]

    def gw(a, b, marker):
        return sp.sympify(Gh[(a, b)])*marker[0 if a in (0, 3) else 1]

    N = {}
    def add(key, value):
        N[key] = N.get(key, 0)+value

    dens = "bubble"
    if graph == "D01":
        vt = V["V3_grav"]
        va = route(vt, [(0, 1), (-1, 0), (1, -1)])
        vb = route(vt, [(0, -1), (1, 0), (-1, 1)])
        for (a, c, e), pa in va.items():
            for (b, d, f), pb in vb.items():
                if (c, d) in Pi and (e, f) in Pi:
                    add((a, b), pa*pb*pw(c, d, u)*pw(e, f, v))
        pref = sp.sympify(AR.base_rules()["graphs"]["D01"]["S_f"]) \
            *(I*kap)**2*I**(2*vt.phase)*I**2

    elif graph == "D03":
        vt = V["V_cb_c_h"]
        va = route(vt, [(-1, 0), (1, -1), (0, 1)])
        vb = route(vt, [(-1, 1), (1, 0), (0, -1)])
        for (rho, mu, a), pa in va.items():
            for (sg, nu, b), pb in vb.items():
                if (mu, sg) in Gh and (nu, rho) in Gh:
                    add((a, b), pa*pb*gw(mu, sg, u)*gw(nu, rho, v))
        pref = -(I*kap)**2*I**(2*vt.phase)

    elif graph == "D11":
        vt = V["V_h_phi_phi"]
        va = route(vt, [(0, 1), (-1, 0), (1, -1)])
        vb = route(vt, [(0, -1), (1, 0), (-1, 1)])
        for (a, _, _), pa in va.items():
            for (b, _, _), pb in vb.items():
                add((a, b), pa*pb*u[0]*v[0])
        pref = Q(1, 2)*(I*kap)**2*I**(2*vt.phase)*I**2

    elif graph == "D04":
        vt = V["V_h_phi_phi"]
        va = route(vt, [(-1, 0), (0, 1), (1, -1)])
        vb = route(vt, [(1, 0), (0, -1), (-1, 1)])
        for (c, _, _), pa in va.items():
            for (d, _, _), pb in vb.items():
                if (c, d) in Pi:
                    add((0, 0), pa*pb*pw(c, d, u)*v[0])
        pref = (I*kap)**2*I**(2*vt.phase)*I**2

    elif graph == "D06":
        vt = V["V_cb_c_h"]
        va = route(vt, [(0, -1), (1, 0), (-1, 1)])
        vb = route(vt, [(-1, 0), (0, 1), (1, -1)])
        for (rho, mu, a), pa in va.items():
            for (sg, nu, b), pb in vb.items():
                if (mu, sg) in Gh and (a, b) in Pi:
                    add((rho, nu), pa*pb*gw(mu, sg, u)*pw(a, b, v))
        pref = (I*kap)**2*I**(2*vt.phase)*I

    elif graph == "D13":
        rules = AR.base_rules()
        vk, vg = AR.K_vertex(rules), V["V_cb_c_h"]
        va = route(vk, [(0, -1), (1, 0), (-1, 1)])
        vb = route(vg, [(-1, 0), (0, 1), (1, -1)])
        for (ap, mu, c), pa in va.items():
            for (rho, nu, d), pb in vb.items():
                if (mu, rho) in Gh and (c, d) in Pi:
                    add((ap, nu), pa*pb*gw(mu, rho, u)*pw(c, d, v))
        r = rules["graphs"]["D13"]
        pref = sp.sympify(r["ext_sign"]*r["S_f"]*r["loop_sign"]) \
            *(I*kap)**2*I**(vk.phase+vg.phase)*I

    elif graph == "D02":
        dens = "tadpole"
        vt = V["V4_grav"]
        va = route(vt, [(0, 1), (0, -1), (1, 0), (-1, 0)])
        for (a, b, c, d), pa in va.items():
            if (c, d) in Pi:
                add((a, b), pa*pw(c, d, u)*v[0])
        pref = Q(1, 2)*(I*kap**2)*I**vt.phase*I

    elif graph == "D12":
        dens = "tadpole"
        vt = V["V_hh_phi_phi"]
        va = route(vt, [(0, 1), (0, -1), (1, 0), (-1, 0)])
        for (a, b, _, _), pa in va.items():
            add((a, b), pa*u[0]*v[0])
        pref = Q(1, 2)*(I*kap**2)*I**vt.phase*I

    elif graph == "D05a":
        dens = "tadpole"
        vt = V["V_hh_phi_phi"]
        va = route(vt, [(1, 0), (-1, 0), (0, 1), (0, -1)])
        for (c, d, _, _), pa in va.items():
            if (c, d) in Pi:
                add((0, 0), pa*pw(c, d, u)*v[0])
        pref = Q(1, 2)*(I*kap**2)*I**vt.phase*I
    else:
        raise ValueError("unknown graph")

    # Exact circle average via a Laurent variable z=e^(i theta).
    z, w = sp.symbols("z w", nonzero=True)
    cosine = (z+1/z)/2
    sine = (z-1/z)/(2*I)
    lm0, lm3 = I*cosine/w, sine/w
    l2 = -1/w**2
    den = l2 if dens == "tadpole" else \
        l2*((lm0-q0)**2-(lm3-q3)**2)

    def pole(expr):
        e = sp.expand(expr.subs({LS[0]: lm0, LS[3]: lm3}))
        if e == 0:
            return sp.Integer(0)
        ser = sp.series(e/den, w, 0, 3).removeO()
        g2 = sp.expand(ser).coeff(w, 2)
        average = laurent_coefficient(g2, z, 0)
        # i/(4 pi^2) times 2 pi times the circle average.
        return sp.cancel(I*average/(2*sp.pi))

    if graph in ("D04", "D05a"):
        domain = [(0, 0)]
    elif graph == "D06":
        domain = [(a, b) for a in range(4) for b in range(4)]
    elif graph == "D13":
        domain = [(a, b) for a in range(10) for b in range(4)]
    else:
        domain = [(a, b) for a in range(10) for b in range(10)]

    cs = physical_cs(branch)
    values, states, mixed = {}, {}, {}
    for key in domain:
        e = sp.expand(N.get(key, 0))
        pair_poles = {}
        if e != 0:
            for mon, cf in sp.Poly(e, *u, *v).terms():
                iu = [i for i in range(3) if mon[i]]
                iv = [i for i in range(3) if mon[3+i]]
                if len(iu) != 1 or len(iv) != 1 \
                        or sum(mon[:3]) != 1 or sum(mon[3:]) != 1:
                    raise ValueError("unexpected channel-marker degree")
                s, t = iu[0], iv[0]
                if dens == "tadpole":
                    t = s
                pv = pole(cf)
                pair_poles[(s, t)] = pair_poles.get((s, t), 0)+pv
        pair_poles = {st: sp.factor(vv)
                      for st, vv in pair_poles.items()
                      if sp.simplify(vv) != 0}
        bad = {str(st): str(vv) for st, vv in pair_poles.items()
               if st[0] != st[1]}
        skey = str(key)
        if bad:
            mixed[skey] = bad
        if bad and branch != "UNIT":
            states[skey] = "EXCLUDED_MIXED_CHANNEL"
            continue
        total = 0
        for (s, t), pv in pair_poles.items():
            total += pv if branch == "UNIT" else pv*cs[s]
        value = sp.factor(sp.cancel(-I*pref*total))
        if value == 0:
            states[skey] = "COMPUTED_ZERO_SLICE"
        else:
            states[skey] = "COMPUTED_NONZERO_SLICE"
            values[skey] = str(value)

    result = {
        "graph": graph, "branch": branch,
        "status": "ALGEBRAIC_SLICE_ONLY",
        "prescription_status": "R2_R3_R4_R5_NOT_CERTIFIED",
        "prefactor_status": "DECLARED_FROZEN_SOURCE_CONVENTIONS",
        "pref": str(pref),
        "key_status": states,
        "mixed_channel_poles_before_prefactor": mixed,
        "Gamma_pole_coeff_of_1/(4-d)": values,
        "seconds": time.time()-start}
    path = Path(outdir)/("spec_"+graph+"_"+branch+".json")
    write_json(path, result)
    return {"written": str(path), "status": result["status"],
            "excluded_keys": sum(s == "EXCLUDED_MIXED_CHANNEL"
                                 for s in states.values())}


def tree_normalization(AR):
    F = AR.F
    if tuple(F.HC) != HC:
        raise ValueError("coordinate order mismatch")
    vt = F.grav_vertex(2)
    ks = [F.ksym(j) for j in range(2)]
    names = {4*j+r: ks[j][r] for j in range(2) for r in range(4)}
    p = (q0, 0, 0, q3)
    h00, h03, h33, f = sp.symbols("h00 h03 h33 f")
    # De Donder equations on h11=h22=f.
    F0 = q0*(h00+h33+2*f)/2-q3*h03
    F3 = q0*h03-q3*(h00+h33-2*f)/2
    solutions = sp.solve((F0, F3), (h00, h03), dict=True)
    if len(solutions) != 1:
        raise AssertionError("generic de Donder solve failed")
    sol = solutions[0]
    hv = {0: sol[h00], 3: sol[h03], 9: h33, 4: f, 7: f}
    Eu = einstein_coordinates(hv)
    ratios = []
    for a in range(10):
        val = 0
        for b in range(10):
            poly = vt.T.get((a, b))
            if poly is None:
                continue
            e = I**vt.phase*poly.to_sympy(names)
            sub = {ks[0][r]: p[r] for r in range(4)}
            sub.update({ks[1][r]: -p[r] for r in range(4)})
            val += e.subs(sub)*hv.get(b, 0)
        val = sp.cancel(val)
        den = sp.cancel(Eu[a])
        if den == 0:
            zero(val, "tree zero-coordinate comparison")
        else:
            ratios.append(sp.cancel(val/den))
    if not ratios:
        raise AssertionError("tree check had no nonzero probe")
    N = ratios[0]
    for r in ratios[1:]:
        zero(r-N, "tree normalization not constant")
    zero(N-1, "declared N")
    Agh = sp.sympify(F.A_GH)
    zero(Agh-Q(1, 2), "declared A_GH")
    zero(N/(I*Agh)+2*I, "tree-predicted lambda")
    return {"N": str(N), "A_GH": str(Agh),
            "lambda_pred": str(N/(I*Agh)),
            "scope": "TREE_NORMALIZATION_ON_DECLARED_DD_CLASS"}


def manifest(root):
    root = Path(root).resolve(strict=True)
    if not root.is_dir():
        raise ValueError("manifest root is not a directory")
    files = []
    for p in sorted(root.rglob("*"),
                    key=lambda p: p.relative_to(root).as_posix()):
        if p.is_symlink():
            raise ValueError("symlink requires a separate policy: " + str(p))
        rel = p.relative_to(root)
        if ".pytest_cache" in rel.parts or p.suffix == ".pyc":
            continue
        if p.is_file():
            digest = hashlib.sha256(p.read_bytes()).hexdigest()
            files.append((rel.as_posix(), digest))
    tree = hashlib.sha256()
    for rel, digest in files:
        name = rel.encode("utf-8")
        tree.update(len(name).to_bytes(8, "big"))
        tree.update(name)
        tree.update(bytes.fromhex(digest))
    return {"serialization":
                "sorted UTF8 relative names; 8-byte length; name; raw SHA256",
            "exclusions": [".pytest_cache path component", "*.pyc"],
            "symlinks": "REJECTED",
            "files": dict(files), "tree_sha256": tree.hexdigest()}


def main():
    p = argparse.ArgumentParser()
    sub = p.add_subparsers(dest="command", required=True)
    sub.add_parser("selftest")
    lp = sub.add_parser("logs")
    lp.add_argument("root")
    lp.add_argument("--legacy", action="store_true")
    rp = sub.add_parser("replay")
    rp.add_argument("graph", choices=GRAPHS)
    rp.add_argument("branch", choices=("CR", "F", "UNIT"))
    rp.add_argument("--engine", required=True)
    rp.add_argument("--out", required=True)
    tp = sub.add_parser("tree")
    tp.add_argument("--engine", required=True)
    mp = sub.add_parser("manifest")
    mp.add_argument("root")
    a = p.parse_args()
    if a.command == "selftest":
        result = {**check_heat_and_mellin(), **check_catalogue_algebra(),
                  "execution_claim":
                      "THIS_OUTPUT_EXISTS_ONLY_IF_THE_USER_RUNS_THE_PROGRAM",
                  "O1": "NOT_CLOSED",
                  "bulk_symmetry_restored": False}
    elif a.command == "logs":
        result = compare_logs(a.root, a.legacy)
    elif a.command == "replay":
        result = replay(load_engine(a.engine), a.graph, a.branch, a.out)
    elif a.command == "tree":
        result = tree_normalization(load_engine(a.engine))
    else:
        result = manifest(a.root)
    print(json.dumps(result, indent=2, sort_keys=True))


if __name__ == "__main__":
    main()
\end{lstlisting}
% END EMBEDDED PYTHON cone_v7.py

\subsection{Reproducibility obligations not replaced by this listing}

A full computational submission should additionally provide:
\begin{enumerate}[label=\textup{(\arabic*)}]
\item immutable frozen input files and their actual byte manifests;
\item the graph inventory and the complete coordinate/domain
coverage ledger;
\item exact contracted spectral expressions before evaluation at
zero transverse spectral parameter;
\item the regulated trace reduction, with boundary defects;
\item the UV/IR separation and meromorphic interchange certificate;
\item source normalization and the complete relevant ST functional;
\item the full basis projection and exact rank/span certificates;
\item software versions, commands, process return codes, raw
outputs, failed runs, and replay-versus-rebuild declarations.
\end{enumerate}
A successful run of the embedded program would not, by itself,
provide items 3--5 or 7.
v8 ran it (Section~\ref{sec:v8response}) and separately supplies
item 7 on a finite scope (Section~\ref{sec:v8cert}); items 3--5 are
not supplied.

%======================================================================
\section{v8 artifacts, hashes and rerun commands}
\label{app:v8artifacts}
%======================================================================

Paths are relative to \loc{/workspace/XVI-O1-CONE/} on the audit box.
Long jobs were run detached with \loc{setsid} and logged. The
environment is Python~3.13.5, SymPy~1.14.0, with \loc{LC_ALL=C.UTF-8}.
The frozen trees were hashed before the first v8 run and after the last
one; the before/after values are identical: XVI-P: yes, XVI-Q: yes, XVI-R: yes.
A full manifest is \loc{logs/v8/artifact_manifest_v8.sha256}; the
\TeX{} source and PDF are hashed in \loc{SHA256SUMS_v8}, outside this
document.

\paragraph{Rerun commands.}
\begin{verbatim}
cd /workspace/XVI-O1-CONE
python3 code/v8/cone_v7.py selftest
python3 code/v8/cone_v7.py logs logs/v4 --legacy
bash code/v8/run_replay.sh   # 27 runs
python3 code/v8/cone_v7.py logs logs/v8/replay
python3 code/v8/ward_scan_radial.py 0.85
python3 code/v8/apex_exact_run.py 7071/10000 CR
python3 code/v8/apex_certificate.py 3/10 CR      # also 3/10 F, 7071/10000 CR|F
python3 code/v8/k_support.py 3_10_CR             # per tag
python3 code/v8/hh_map.py 3/10 CR                # per (alpha, branch)
python3 code/v8/exact_values.py 3_10_CR 3_10_F 7071_10000_CR 7071_10000_F
bash code/v8/run_slice.sh && bash code/v8/run_slice2.sh
python3 code/v8/evanescent_combine.py
python3 code/v8/duhamel_mode_check2.py
python3 code/artifacts_v8.py && python3 code/build_v8.py
pdflatex -halt-on-error XVI-O1-CONE-v8.tex   # twice
\end{verbatim}
The certificate runs take 5--10 minutes each and stay below 2~GB of
memory. The box limit is about 7~GB, so they were run at most three
at a time.

{\small
\begin{longtable}{@{}p{.50\linewidth}p{.46\linewidth}@{}}
\caption{SHA-256 of v8 artifacts.}\\
\toprule
File & SHA-256\\
\midrule
\endfirsthead
\toprule
File & SHA-256\\
\midrule
\endhead
\bottomrule
\endfoot
\multicolumn{2}{@{}l}{\textit{Base source (read-only)} (1 files)}\\
{\scriptsize\loc{XVI-O1-CONE-v7.tex}} & \texttt{\tiny d0ba7b9cdbf535d1296a74146011c1cb1269ff53148d3a96add64574873d5132}\\
\multicolumn{2}{@{}l}{\textit{v8 programs} (23 files)}\\
{\scriptsize\loc{code/v8/apex_certificate.py}} & \texttt{\tiny 14a49b42dcb7c1fb53a69825a75a9a6dcfa36ad4f88f7f8ca80e95f47818b1ca}\\
{\scriptsize\loc{code/v8/apex_certificate_la.py}} & \texttt{\tiny 803c363e5aa3fd32824286bf5f4da891067c2a6e91c80705acf14178787cfb21}\\
{\scriptsize\loc{code/v8/apex_exact.py}} & \texttt{\tiny ba6b0946e84e9f2fb3a60a24d38a3e4343e3887ae6c79c385b88e5c0db4ec46b}\\
{\scriptsize\loc{code/v8/apex_exact_run.py}} & \texttt{\tiny db8f950dd31e012044776bddc3bd594aa2ededead73794459b71301321e23d48}\\
{\scriptsize\loc{code/v8/cone_v7.py}} & \texttt{\tiny 55b197552eb4c400c526cd39e4964a4cc8ec1fb89cf45dd41b1336dd42ba9b07}\\
{\scriptsize\loc{code/v8/duhamel_mode_check.py}} & \texttt{\tiny b7a1183a11c4f0219b700a7d6ca0056f63aa8fc9efd3012d89a156d8310794f5}\\
{\scriptsize\loc{code/v8/duhamel_mode_check2.py}} & \texttt{\tiny b010c66de54a5b7fde6bfc8efd42a4d618533f65390337ee0cdb8d250b698aaa}\\
{\scriptsize\loc{code/v8/evanescent_combine.py}} & \texttt{\tiny e967018847f31c26448a832381c558fb01dde75dfb36bf60402e8873c8e58a0c}\\
{\scriptsize\loc{code/v8/exact_values.py}} & \texttt{\tiny 4b4ad1682afb5629f24df91b97a9a905a980ba1fa034fa060c087dfe9f7bfae3}\\
{\scriptsize\loc{code/v8/gq.py}} & \texttt{\tiny a48693860dde5d3d0b244ca0b13c00cea1988849a45a3b16fe39d1cc7e1eae3b}\\
{\scriptsize\loc{code/v8/hh_map.py}} & \texttt{\tiny 372b3f01f084e1bb8091c933d3d979a445b445febc6610694048a7253c850cee}\\
{\scriptsize\loc{code/v8/k_support.py}} & \texttt{\tiny 471f02fa05515d76d4fd94b57b38811bdcafac450d0cd900e9466dfa076aa94f}\\
{\scriptsize\loc{code/v8/make_apex_exact.py}} & \texttt{\tiny 2f4e74c6c3ce1c0d92470f7babbb27f0b9d53515cd390270d352303cdc6da6fa}\\
{\scriptsize\loc{code/v8/make_apex_exact_run.py}} & \texttt{\tiny c8e9abbb5402d35f1c97dafce62215c94d0bfae44e3ba6b266bedc23902fcb18}\\
{\scriptsize\loc{code/v8/slice_engine.py}} & \texttt{\tiny 296bfb8541ded14d27f1050cf8602693fb1a6f4c7c99cdd378b57eaed40069ed}\\
{\scriptsize\loc{code/v8/slice_run.py}} & \texttt{\tiny 2b46ef8d8c6dc8c9085725eccdf0d89aad609098ea913e57ba55f9a92b2e0688}\\
{\scriptsize\loc{code/v8/test_slice.py}} & \texttt{\tiny 75acc5a9c6deea06be6623038e745ddbff481bb12ae7b2a1dc23d65f4279cf70}\\
{\scriptsize\loc{code/v8/ward_scan_radial.py}} & \texttt{\tiny 453a8f3dfd959465b5722a710485fea655b4265af5cca914e7d09499049fafca}\\
{\scriptsize\loc{code/v8/run_replay.sh}} & \texttt{\tiny c459fd5759b7cdec921d12fdd27ad3241a118c5c12361541e23b3916e07ef968}\\
{\scriptsize\loc{code/v8/run_slice.sh}} & \texttt{\tiny 2ff93bab2df9006a0054b71f39b4f055f475eb2c230fc9d740bec708b64e8d7b}\\
{\scriptsize\loc{code/v8/run_slice2.sh}} & \texttt{\tiny bd47fcb81bdfa021554c3d667088249b3973ac8dde448239d696955778e90251}\\
{\scriptsize\loc{code/build_v8.py}} & \texttt{\tiny c83d29a050937ab65347255ece236c47fc3ef8f78063f50710391532f2bf1263}\\
{\scriptsize\loc{code/artifacts_v8.py}} & \texttt{\tiny adb9abcb3d879cdb7eb32e1956c16ae21ada684b5bb3aabcb0b4c8ab0f7f7766}\\
\multicolumn{2}{@{}l}{\textit{v8 source fragments} (9 files)}\\
{\scriptsize\loc{v8src/abstract.tex}} & \texttt{\tiny f5d9e2b3fd61a71d616d41562c68ba2d1bec5537f5c069b29dcaa0abb5c7d130}\\
{\scriptsize\loc{v8src/boundary.tex}} & \texttt{\tiny 33dd134192a6b257f9065161aba3ebf2df4238cb834ebb9dd066465c57e40fa3}\\
{\scriptsize\loc{v8src/opentable.tex}} & \texttt{\tiny 03a4908ff49b8f49d0ad402839aafcd70ad891e57a8a20f01b54322758fbe4bb}\\
{\scriptsize\loc{v8src/response.tex}} & \texttt{\tiny 7dd77765cbdc314c760cebc8eaa66839f0a4c329cb706075098a0399802d9336}\\
{\scriptsize\loc{v8src/sec_evan.tex}} & \texttt{\tiny 47e4623fa8bc6200026cf430642b705674154a95d40e472bc37ad95a17673422}\\
{\scriptsize\loc{v8src/sec_map.tex}} & \texttt{\tiny 83e8eb0f2e6b6b24ee812bd5924a0da5344c98bd34a7d091fdafed4ef2d0d1f8}\\
{\scriptsize\loc{v8src/sec_mixed.tex}} & \texttt{\tiny 23807ea310beb231d885e08bed37aebbe69014c601b91becbc7d08b0f57fb1ed}\\
{\scriptsize\loc{v8src/whatv8.tex}} & \texttt{\tiny 7d654d0008959917d3dee81e41930b6e68ad7bc058f9897b3fa6c621becc021b}\\
{\scriptsize\loc{v8src/bibmarks.json}} & \texttt{\tiny 12e6130eabfd85e6fc1841c039149ad8688cd32a66f934b606825675f3cc4cd3}\\
\multicolumn{2}{@{}l}{\textit{v7-program runs} (4 files)}\\
{\scriptsize\loc{logs/v8/v7_logs_legacy.json}} & \texttt{\tiny 0674a32ac36a26c0a485c5591fb006bc4739e4159c6019912fb2860c37b1ca4f}\\
{\scriptsize\loc{logs/v8/v7_logs_replay.json}} & \texttt{\tiny 856162e1d1ba8365213ee7af80e1a661a02535b0512d011d99aa5a353d8871de}\\
{\scriptsize\loc{logs/v8/v7_selftest.json}} & \texttt{\tiny 5d0668fc426d17a17cb0b058f78ce919532f7003c0d8d548f0440a4bb10277ef}\\
{\scriptsize\loc{logs/v8/v7_tree.json}} & \texttt{\tiny acf9699d11beaca179da32ffab88ee57366f1cbe1da8ea48aa8e7876538d5391}\\
\multicolumn{2}{@{}l}{\textit{Replay} (27 files)}\\
{\scriptsize\loc{logs/v8/replay/spec_D01_CR.json}} & \texttt{\tiny acce791a88e9204de5691ef10b411cc9809625ea5fa1cfe2afb0eb21c54b6003}\\
{\scriptsize\loc{logs/v8/replay/spec_D01_F.json}} & \texttt{\tiny 79a84353204e8e9e50fc91d924b4c4da91bab831d5a10e98ffe74dc118f13f03}\\
{\scriptsize\loc{logs/v8/replay/spec_D01_UNIT.json}} & \texttt{\tiny 5649acd95b4f21ca4a564945e3eb35cc968a051c0b6e5f29934e5384b5cbea8c}\\
{\scriptsize\loc{logs/v8/replay/spec_D02_CR.json}} & \texttt{\tiny 7f4a1d36f5f1723abe33324b654b07cb5e51bde054f82c7e9235c3abe5decd50}\\
{\scriptsize\loc{logs/v8/replay/spec_D02_F.json}} & \texttt{\tiny 99bb7ba60647bc744d2da60077179727729109379e857bbf3f9f8ff466035062}\\
{\scriptsize\loc{logs/v8/replay/spec_D02_UNIT.json}} & \texttt{\tiny fd11e5c3f5272aa46177f9adcd2cf56fd30b31b1f28ff2f6118f6afdc5959dbb}\\
{\scriptsize\loc{logs/v8/replay/spec_D03_CR.json}} & \texttt{\tiny 2bec5983c7c649fe57e29452e4819847e581d9a4ff059c7df4f50127034f14c6}\\
{\scriptsize\loc{logs/v8/replay/spec_D03_F.json}} & \texttt{\tiny d3f142193c21fb2bcb45bb44692cba7cd09a59131cfc2557ee084cbf7aa23abb}\\
{\scriptsize\loc{logs/v8/replay/spec_D03_UNIT.json}} & \texttt{\tiny 3e1e4792fb8ec91035b5ef3c9bb6343255c46d1f099c28750a56ef3044e6f24c}\\
{\scriptsize\loc{logs/v8/replay/spec_D04_CR.json}} & \texttt{\tiny a2151dc64d8159bd2183dd4fb84b9fa69b2500ccb8e959282279fae9ba2d3d71}\\
{\scriptsize\loc{logs/v8/replay/spec_D04_F.json}} & \texttt{\tiny 542b00df420c2d5bb8cc884b8103cc9fa705d3e60a06458bcbd0ef8f5da2ae19}\\
{\scriptsize\loc{logs/v8/replay/spec_D04_UNIT.json}} & \texttt{\tiny 512cdf20a2a5df3a0ef7df19acd867b5c70e65a53ca3c57b34adf14b36cb7446}\\
{\scriptsize\loc{logs/v8/replay/spec_D05a_CR.json}} & \texttt{\tiny 179e86ad7825cfebf3a48ab4e670b639518fbab389513f9fd220bc458e734317}\\
{\scriptsize\loc{logs/v8/replay/spec_D05a_F.json}} & \texttt{\tiny 4f5cbe0ae2b9a9965780a290bc3652e960c356383716794b659f2549f94e12ea}\\
{\scriptsize\loc{logs/v8/replay/spec_D05a_UNIT.json}} & \texttt{\tiny ca56b30462aa4fca5a73bac77870b9f55e377f43fe9cd7afe17d1ce0ecdd105f}\\
{\scriptsize\loc{logs/v8/replay/spec_D06_CR.json}} & \texttt{\tiny 62ef7ca5c49060ff207c284497ff1eaf60e7214227e0164e5e8fb13a6fb0d1f0}\\
{\scriptsize\loc{logs/v8/replay/spec_D06_F.json}} & \texttt{\tiny 43575824758f99a429e9980f88bb9441687a209a0e7f6f4b0172a039e554ed2b}\\
{\scriptsize\loc{logs/v8/replay/spec_D06_UNIT.json}} & \texttt{\tiny e0967ee852a1f35946ea020ca4afc39b83f961bcae556b0ff6f890493dd97916}\\
{\scriptsize\loc{logs/v8/replay/spec_D11_CR.json}} & \texttt{\tiny 05f947251f8a3d1cee8c721e998a5e431c71fe87d75abb76881aac32d99698d5}\\
{\scriptsize\loc{logs/v8/replay/spec_D11_F.json}} & \texttt{\tiny a350cacd428979ff494fdc2b7d0069a12a895520420344259fee3252bac0b0a5}\\
{\scriptsize\loc{logs/v8/replay/spec_D11_UNIT.json}} & \texttt{\tiny b9c70d687aee15bd7c9f53697fb04264da32cba89e4824c0eb5d1fe1400fe30a}\\
{\scriptsize\loc{logs/v8/replay/spec_D12_CR.json}} & \texttt{\tiny 6675f26bf142c5e549d2d4e4ed78640f30de99bbc2c8c3d57c6b5de4a19059fa}\\
{\scriptsize\loc{logs/v8/replay/spec_D12_F.json}} & \texttt{\tiny 356db7ea109fe8dfce4e29c5c67bf7c3a6c525bcb0fb7cb4c3e5f66b12674ee6}\\
{\scriptsize\loc{logs/v8/replay/spec_D12_UNIT.json}} & \texttt{\tiny 8c3d2c8aa6018bac44281d5d21aeb77cd35c53cb3663fabe1ea920da5196f6c0}\\
{\scriptsize\loc{logs/v8/replay/spec_D13_CR.json}} & \texttt{\tiny 794463632e6f4f6303047f8c94b12172874b60f1dce1c5d35bdfec781c6eb249}\\
{\scriptsize\loc{logs/v8/replay/spec_D13_F.json}} & \texttt{\tiny ff97041eff19ac02e6009fd36f30748715f8856eaaa5dc6fef165d82ba78342f}\\
{\scriptsize\loc{logs/v8/replay/spec_D13_UNIT.json}} & \texttt{\tiny d4041679d0a9c8b27237d1746a9c72dc661bba20b8c76591f8ab7a37570b7d9b}\\
\multicolumn{2}{@{}l}{\textit{Ward scan} (2 files)}\\
{\scriptsize\loc{logs/v8/ward/wardscan_radial_0.85_fried.json}} & \texttt{\tiny c802124092dc962b8fa7ccac419650eeae2a906b78b5ab5efbdd32bfa6675258}\\
{\scriptsize\loc{logs/v8/ward/wardscan_radial_0.85_geom.json}} & \texttt{\tiny 007709fdd124593107898379c9c318e188c6a926d67905ff2c3908d496ef57a3}\\
\multicolumn{2}{@{}l}{\textit{Exact certificates and maps} (17 files)}\\
{\scriptsize\loc{logs/v8/cert/apex_cert_3_10_CR.json}} & \texttt{\tiny 001aa482da32b879365eaad7cd97e584b587895074354840f14b0ad94df53a96}\\
{\scriptsize\loc{logs/v8/cert/apex_cert_3_10_CR_vectors.json}} & \texttt{\tiny a719330531cb070bea918f2f8a0b755e48e0100974756bc2aa932e779cff9319}\\
{\scriptsize\loc{logs/v8/cert/apex_cert_3_10_F.json}} & \texttt{\tiny fc76e1a9d2fb95eb2256c5ee4a24353d8efce1c9e38942a67b65ca64d979f93b}\\
{\scriptsize\loc{logs/v8/cert/apex_cert_3_10_F_vectors.json}} & \texttt{\tiny a719330531cb070bea918f2f8a0b755e48e0100974756bc2aa932e779cff9319}\\
{\scriptsize\loc{logs/v8/cert/apex_cert_7071_10000_CR.json}} & \texttt{\tiny 204c46e1b8ac59609072a060018373f67ac19e93717fbdc56b0547b863330a0b}\\
{\scriptsize\loc{logs/v8/cert/apex_cert_7071_10000_CR_vectors.json}} & \texttt{\tiny a719330531cb070bea918f2f8a0b755e48e0100974756bc2aa932e779cff9319}\\
{\scriptsize\loc{logs/v8/cert/apex_cert_7071_10000_F.json}} & \texttt{\tiny f3fd2f7010379d22b7236f92247f19c8a82a1eaa9947b7c7be27e9cdcff02af0}\\
{\scriptsize\loc{logs/v8/cert/apex_cert_7071_10000_F_vectors.json}} & \texttt{\tiny ef66990df199ceca552c9986688ebfe82f3c1e050bedfe76b16a415ed40dc728}\\
{\scriptsize\loc{logs/v8/cert/exact_values.json}} & \texttt{\tiny 21aefbf917ccda4a9a70bc0bc8a1405ea2e556713bd6bd4aa104aeed757926e7}\\
{\scriptsize\loc{logs/v8/cert/hh_map_3_10_CR.json}} & \texttt{\tiny ee9d20a4b630496ba5340c74e05ae16d2613d21a1b6cb05433b21f5cc1939894}\\
{\scriptsize\loc{logs/v8/cert/hh_map_3_10_F.json}} & \texttt{\tiny de28306ff18ec729efd27eca9269fa6d75a2fef8839950c4f11978ae41c1c40e}\\
{\scriptsize\loc{logs/v8/cert/hh_map_7071_10000_CR.json}} & \texttt{\tiny d6c9abe28830ccd985c69466ee1eca53344b69b27216ea3e8a8ef68100aebad5}\\
{\scriptsize\loc{logs/v8/cert/hh_map_7071_10000_F.json}} & \texttt{\tiny e6f492d4cf6f19b6c9b75f3be96b13ac97a0382c23f8902f92261beb9f30742e}\\
{\scriptsize\loc{logs/v8/cert/k_support_3_10_CR.json}} & \texttt{\tiny 0e442beaef6ed47ef2e53bfc7f7f9034a83536a4cacda009f88e0fb5fbe22a45}\\
{\scriptsize\loc{logs/v8/cert/k_support_3_10_F.json}} & \texttt{\tiny 0e442beaef6ed47ef2e53bfc7f7f9034a83536a4cacda009f88e0fb5fbe22a45}\\
{\scriptsize\loc{logs/v8/cert/k_support_7071_10000_CR.json}} & \texttt{\tiny 0e442beaef6ed47ef2e53bfc7f7f9034a83536a4cacda009f88e0fb5fbe22a45}\\
{\scriptsize\loc{logs/v8/cert/k_support_7071_10000_F.json}} & \texttt{\tiny 0e442beaef6ed47ef2e53bfc7f7f9034a83536a4cacda009f88e0fb5fbe22a45}\\
\multicolumn{2}{@{}l}{\textit{Exact sampled port} (1 files)}\\
{\scriptsize\loc{logs/v8/apex_exact/apex_7071_10000_CR_seed0.json}} & \texttt{\tiny 9b14e0b3f879676cdcd97d162f5507ff9fb9ae64026473b76c00da0d2ad8f4d1}\\
\multicolumn{2}{@{}l}{\textit{Slice / evanescent / hidden pairs} (11 files)}\\
{\scriptsize\loc{logs/v8/slice/evanescent_combined.json}} & \texttt{\tiny de6d4bf42a5477bb9c82a331d0a5d77447dee43f8b950ee1140aaf5832ef1054}\\
{\scriptsize\loc{logs/v8/slice/hidden_mixed.json}} & \texttt{\tiny 4f4c0cb1d3e8a0b53adcd0c94c39515f8d7b06b64f00c043de47d8a2b8731077}\\
{\scriptsize\loc{logs/v8/slice/slice_D01.json}} & \texttt{\tiny 0a55a654154890d9c8c50e70638d47f59340c05b8afa87842c6e9be2b492f6fb}\\
{\scriptsize\loc{logs/v8/slice/slice_D02.json}} & \texttt{\tiny 06e19f1c9e8f2acd0f4546c03e3b3e8a4273456c87b5304eab5551f23ef92e74}\\
{\scriptsize\loc{logs/v8/slice/slice_D03.json}} & \texttt{\tiny c7eb6a6efcfc389f9d2439f2472b62b47ef0ef6201deff6f4ef92d968acf418d}\\
{\scriptsize\loc{logs/v8/slice/slice_D04.json}} & \texttt{\tiny cc23ad207c516682715544fbafe2434dce83e84785d19e8383adeab0c1a833b7}\\
{\scriptsize\loc{logs/v8/slice/slice_D05a.json}} & \texttt{\tiny 5a6b5c35ae58f9197a75128b7e002ebee0a8688682b0d4e8f8414687861fa897}\\
{\scriptsize\loc{logs/v8/slice/slice_D06.json}} & \texttt{\tiny cf0bcdad0d20b76076fcf327fafbbe5c5f2a6a568537e98fd7d4383e64d09786}\\
{\scriptsize\loc{logs/v8/slice/slice_D11.json}} & \texttt{\tiny f4afed2f0745d95d685294001dab6e0099720733a6e624df8a5971336fcd5656}\\
{\scriptsize\loc{logs/v8/slice/slice_D12.json}} & \texttt{\tiny acfa367bac579ab8f18c2bd731b3b3550a83aa1862722a4da2250774cdef88be}\\
{\scriptsize\loc{logs/v8/slice/slice_D13.json}} & \texttt{\tiny eeb384f033363738b9bcd5d0160ee3d69299c8426c2c2fb35ee9b62d628226c6}\\
\multicolumn{2}{@{}l}{\textit{Duhamel corroboration} (2 files)}\\
{\scriptsize\loc{logs/v8/duhamel_mode_check.json}} & \texttt{\tiny 2931da69cbc9873704391c64ac13299b41666f1ad59eeaeb67ac4eab63eb86db}\\
{\scriptsize\loc{logs/v8/duhamel_mode_check2.json}} & \texttt{\tiny 74bc8b04bc98559da3c5e3521d6b96d46e0056d5d09bd411c025b02c56b5fae7}\\
\multicolumn{2}{@{}l}{\textit{Manifests, environment} (5 files)}\\
{\scriptsize\loc{logs/v8/manifest_XVIP.json}} & \texttt{\tiny 8dcd372a5977065d7f30e9698917c9b371faf5caa111a84bd9d27a2e4c2714a6}\\
{\scriptsize\loc{logs/v8/manifest_XVIQ.json}} & \texttt{\tiny c33ef1abe8d6f63a28885a7d41bc70d27455cf45eba06a94697270d5cd0ab337}\\
{\scriptsize\loc{logs/v8/manifest_XVIR.json}} & \texttt{\tiny 44f17586766b4c13c088080143a99864b397dbc1893a7b39422855e88876fa17}\\
{\scriptsize\loc{logs/v8/tree_hashes_v8.json}} & \texttt{\tiny a6b3246b70cf27eed7eb7d2f7f475903fbb8d1cd7199283ed1b90b5f42bceb8f}\\
{\scriptsize\loc{logs/v8/env.json}} & \texttt{\tiny 4a40f99a724e190dad2fa6e65c49b318415204377fa0723aa71128884aed714e}\\
\end{longtable}}


%======================================================================
\section{v9 artifacts, hashes and rerun commands}
\label{app:v9artifacts}
%======================================================================

Paths are relative to the project directory \loc{XVI-O1-CONE/} under \loc{/workspace}; \loc{intake/} abbreviates the companion intake directory \loc{XVI-R-companion-intake/}. Long jobs
were run detached with \loc{setsid}, logged, and made resumable
with per-item JSON caches, because the box runs out of memory at
about 7~GB and is restored from time to time. Environment:
Python~3.13.5, SymPy~1.14.0, mpmath~1.3.0, \loc{LC_ALL=C.UTF-8}.
Frozen trees before = after: XVI-P: yes, XVI-Q: yes, XVI-R: yes.
The v8 source and PDF are unchanged: yes.
The modular scan comprises 192 result files, which are hashed individually in
\loc{logs/v9/scan/SHA256SUMS}. The full manifest is
\loc{logs/v9/artifact_manifest_v9.sha256}, and the \TeX{} source and
PDF are hashed in \loc{SHA256SUMS_v9}.
\loc{code/v9/mixed02.py} is kept only because it produced the
asymptotic coefficients. Its \({}_3F_2\) shortcut is unreliable at
large parameter and is not used for any printed value.

\paragraph{Rerun commands.}
\begin{verbatim}
cd /workspace/XVI-O1-CONE
python3 code/v9/g_formula_check.py
python3 code/v9/mixed02.py   # -> logs/v9/asym02.json (~250 s)
setsid nohup python3 code/v9/hidden02_run.py \
   > logs/v9/hidden/run.log 2>&1 &
python3 code/v9/cbck_map.py 3/10 CR   # and 3/10 F, 7071/10000 CR|F
python3 code/v9/ck_free.py 3/10 CR    # and the other three points
python3 code/v9/e_solve.py
python3 code/v9/chamber_scan.py       # resumable cache
setsid nohup bash code/v9/scan_driver.sh \
   > logs/v9/scan/driver.log 2>&1 &
setsid nohup bash code/v9/scan_driver_rev.sh \
   > logs/v9/scan/driver_rev.log 2>&1 &
python3 code/v9/evan_match.py
python3 code/v9/vector_logs.py
python3 code/artifacts_v9.py && python3 code/build_v9.py
pdflatex -halt-on-error XVI-O1-CONE-v9.tex   # twice
\end{verbatim}

{\small
\begin{longtable}{@{}p{.50\linewidth}p{.46\linewidth}@{}}
\caption{SHA-256 of v9 artifacts.}\\
\toprule
File & SHA-256\\
\midrule
\endfirsthead
\toprule
File & SHA-256\\
\midrule
\endhead
\bottomrule
\endfoot
\multicolumn{2}{@{}l}{\textit{Base source (read-only)} (1 files)}\\
{\scriptsize\loc{XVI-O1-CONE-v8.tex}} & \texttt{\tiny 153e4e9dbcf116a138d0f3919607bed1f3646b1a8437f8e5cdbcc5dd686e85f8}\\
\multicolumn{2}{@{}l}{\textit{v9 programs} (15 files)}\\
{\scriptsize\loc{code/v9/apex_scan.py}} & \texttt{\tiny 9e03be0cce4c825d42b995384f40846fea9fbac73677acb5e2dfb11edef36529}\\
{\scriptsize\loc{code/v9/cbck_map.py}} & \texttt{\tiny 68539cf07ee63910f5ddab8e3a309ce9a305070b69d9448cb3bacbe43123600e}\\
{\scriptsize\loc{code/v9/chamber_scan.py}} & \texttt{\tiny 499d77465702a3d471325b0abf81324a4545dd57a0f76bd775b7504aa917fe24}\\
{\scriptsize\loc{code/v9/ck_free.py}} & \texttt{\tiny 8cffb01060ba5899e3f1a4055dd0ecf4c1d8eefee915033f656135fe63ea79bf}\\
{\scriptsize\loc{code/v9/e_solve.py}} & \texttt{\tiny c29ad9480e827430f572ef433425ed1123cd86fd98ce93d654d68142287be119}\\
{\scriptsize\loc{code/v9/evan_match.py}} & \texttt{\tiny 260a30e2fa25ba065f9ec141bb104e1b6f3ef018a381f74ebc1db55cf0e28005}\\
{\scriptsize\loc{code/v9/g_formula_check.py}} & \texttt{\tiny 1661f1bc5ec9c3aa4e66e5415630761390e3790e2d92ccb36f4534a284c29def}\\
{\scriptsize\loc{code/v9/hidden02_run.py}} & \texttt{\tiny 9bcafe3ad4254c68bfdaa95d9d11f0e6db8fc235bdbda39829c5c1668f130487}\\
{\scriptsize\loc{code/v9/mixed02.py}} & \texttt{\tiny 808a312ab7bfacb834e58e3d1b8988178e220845df162fd2014c19bd483d2530}\\
{\scriptsize\loc{code/v9/modefun.py}} & \texttt{\tiny 451c91de56316ab65955d0d439cf8f0ed40d63c1431d7e5dcea12149a050689c}\\
{\scriptsize\loc{code/v9/vector_logs.py}} & \texttt{\tiny cd8a433208136e21180d8f14ed7a7e67e3cf31c5be9718192fca17c0c0973064}\\
{\scriptsize\loc{code/v9/scan_driver.sh}} & \texttt{\tiny fbc160089288c70e1a201a9c63481cd74bc13c35b20fbbf27876c785c526851f}\\
{\scriptsize\loc{code/v9/scan_driver_rev.sh}} & \texttt{\tiny e062603c4cdcd258400c4f932ff4583ce2bc1024bd3a75f6c81b2fbe839280dd}\\
{\scriptsize\loc{code/build_v9.py}} & \texttt{\tiny 6449baa74d39ee947f6c537f7adb3aa0d14565b422ee81362f818b7ac430af71}\\
{\scriptsize\loc{code/artifacts_v9.py}} & \texttt{\tiny c5e51e8aa1ab4acaa8918e940872c9440abf81e9134a778867820cf40fd480e3}\\
\multicolumn{2}{@{}l}{\textit{v9 source fragments} (4 files)}\\
{\scriptsize\loc{v9src/abstract.tex}} & \texttt{\tiny 1367810f9f3f59426a4c29b553724f4955113b48b99c71b140aab7e7e5142ea5}\\
{\scriptsize\loc{v9src/bib_v9.tex}} & \texttt{\tiny 0fc845a1e879e4863832c91b1c0e41d42068ba1507cdaaa0d32fb73d3734733b}\\
{\scriptsize\loc{v9src/opentable.tex}} & \texttt{\tiny 5cf973cdb2281f71777c962af3b6056e948ee9dfdbea28b9f2bc7bf2a4c6e477}\\
{\scriptsize\loc{v9src/sec_v9.tex}} & \texttt{\tiny 362ce962df3d6d5f29c9a547ee3d5da7da4c8a3a22b7d37981750806b1396cdb}\\
\multicolumn{2}{@{}l}{\textit{Mixed trace and hidden pairs} (16 files)}\\
{\scriptsize\loc{logs/v9/g_formula_check.json}} & \texttt{\tiny 00a65b245b5217c9c77c05598404dc37334b218e5613aa2088ea74ad44f280aa}\\
{\scriptsize\loc{logs/v9/asym02.json}} & \texttt{\tiny a3cc63868cdb683638e305f62dfca309e33de9e1046e65f70b6e49613a909abe}\\
{\scriptsize\loc{logs/v9/hidden/dJ_11_20.json}} & \texttt{\tiny 6c3d9cea71400827e43587ddd37bec984243ddbb35648662b0f5c5f4e7d12fd8}\\
{\scriptsize\loc{logs/v9/hidden/dJ_17_20.json}} & \texttt{\tiny 512aa07e728aeec0f47b10be0f4860a00d25d9021e0200ce67fd846ef9b89bf3}\\
{\scriptsize\loc{logs/v9/hidden/dJ_19_20.json}} & \texttt{\tiny e7abcfd32ff1e0e0894484b05e13277f8340c2d0c5b1ef88b4799e8e71c1dce6}\\
{\scriptsize\loc{logs/v9/hidden/dJ_1_2.json}} & \texttt{\tiny 6e2ba8202b076ce5ec92c54e284c3ce7cfdbf2210ceff6c4a772dc4aac78bc25}\\
{\scriptsize\loc{logs/v9/hidden/dJ_2_3.json}} & \texttt{\tiny 76194023826555748db3130ed6f04ae20c33000936c38ec69c84b1336ea31ed9}\\
{\scriptsize\loc{logs/v9/hidden/dJ_2_5.json}} & \texttt{\tiny 601caae331eaf8283387fed0336e7fcdc0308dbb8211ccf2fa6d723c437eb573}\\
{\scriptsize\loc{logs/v9/hidden/dJ_3_4.json}} & \texttt{\tiny 3a0eb389c6a9337fb0b455aec786aae697edb21e760d87378ddfd808228b67e8}\\
{\scriptsize\loc{logs/v9/hidden/dJ_3_5.json}} & \texttt{\tiny ece5ba7a75d11370982406ce2f0c3385d754b21308857c409a0e85dec28117e1}\\
{\scriptsize\loc{logs/v9/hidden/dJ_4_5.json}} & \texttt{\tiny b9693cb7179baba41b508894675c27736167ee0ba694b70b975fbc9d999fcad8}\\
{\scriptsize\loc{logs/v9/hidden/dJ_7071_10000.json}} & \texttt{\tiny 29acf5b74f782ea5708a880eb9081d12a5ad9831da748428bdc749e6e1fac245}\\
{\scriptsize\loc{logs/v9/hidden/dJ_7_20.json}} & \texttt{\tiny 5b99ee97bac9315bb96959dd3de3c3ce3b34c734fddf169bb30550fd96b01fc7}\\
{\scriptsize\loc{logs/v9/hidden/dJ_99_100.json}} & \texttt{\tiny fe6dcf8f07d79fb333b66b44f9898fd2335fd2820b36405f8c4ddc7bae5fabd9}\\
{\scriptsize\loc{logs/v9/hidden/dJ_9_10.json}} & \texttt{\tiny 33f5d6897ea5a3dd20f58437155b0b7bb62734038648fcf95ef17b11dabaa621}\\
{\scriptsize\loc{logs/v9/hidden/dJ_9_20.json}} & \texttt{\tiny 7ae4b36c77cbdae96e3d5d22d89f48ca882552947b263fb515d7459a6c330668}\\
\multicolumn{2}{@{}l}{\textit{Ghost dictionary and exact values} (9 files)}\\
{\scriptsize\loc{logs/v9/cbck_map_3_10_CR.json}} & \texttt{\tiny 265f5928ebc1c396f0c7a55e7d5b115631966c95b9f12fef38bf61f5e2cddff9}\\
{\scriptsize\loc{logs/v9/cbck_map_3_10_F.json}} & \texttt{\tiny 0f5e005e2284889ca69c165ca2e225656c4e7c0cf47b41f2449ad129dab30e4d}\\
{\scriptsize\loc{logs/v9/cbck_map_7071_10000_CR.json}} & \texttt{\tiny 8c8f3fb53f10d9160c091206de68e8a1a7dca7544ae87a4a433e6e6b805fdf5b}\\
{\scriptsize\loc{logs/v9/cbck_map_7071_10000_F.json}} & \texttt{\tiny c81a43ba76ad1387775b36a1d317e23bb6c4b030ded84ced361f754ac3cc5831}\\
{\scriptsize\loc{logs/v9/ck_free_3_10_CR.json}} & \texttt{\tiny 5fceb30dbe2e706dff67e8b717d44a8367b5114da9ed7e6d17907de73c55002b}\\
{\scriptsize\loc{logs/v9/ck_free_3_10_F.json}} & \texttt{\tiny 689499331b250796e9f48a458a928f0f379ccf2b9840ebf114a1d395599d8e98}\\
{\scriptsize\loc{logs/v9/ck_free_7071_10000_CR.json}} & \texttt{\tiny 4a2dbdb45f12039736b80afd150f92cb952148c6bf3ae367756e89121a4091b3}\\
{\scriptsize\loc{logs/v9/ck_free_7071_10000_F.json}} & \texttt{\tiny ba3b72e7cd1718401711bca6ce8893b3fe67710e0b2d86e187c2230eb152c404}\\
{\scriptsize\loc{logs/v9/e_solve.json}} & \texttt{\tiny 99c4b03a828c4ec0560ee2178b11500293e2c9a283ec516d59e50ca974b8d8ee}\\
\multicolumn{2}{@{}l}{\textit{Chambers and modular scan} (3 files)}\\
{\scriptsize\loc{logs/v9/chamber_cache2.json}} & \texttt{\tiny 8058faead811c4c2fd3d5cba97de65070dcb27c08447d22a2f62f2963e24e393}\\
{\scriptsize\loc{logs/v9/scan/SHA256SUMS}} & \texttt{\tiny 2f17bf26b9a0f3f989267752da85a21ddc8d23ae406f81943417d5ad034da4e5}\\
{\scriptsize\loc{logs/v9/scan/alphas.txt}} & \texttt{\tiny fbdd209b93ebf0cd45dc65800bba9b24f7be32bc38ce5b1cfa0c0190e12a3239}\\
\multicolumn{2}{@{}l}{\textit{Obstruction checks} (2 files)}\\
{\scriptsize\loc{logs/v9/evan_match.json}} & \texttt{\tiny 1a273d335cbf50dc9a264ba4bed7dcc4edd6f6c279e799f810ae3a579a6dca81}\\
{\scriptsize\loc{logs/v9/vector_logs.json}} & \texttt{\tiny 7531cd19399da49e3952ffcaaa70e76a96b70eb4bfa775b0395f4603ec22a6da}\\
\multicolumn{2}{@{}l}{\textit{Registry} (3 files)}\\
{\scriptsize\loc{logs/v9/sheets_extract.txt}} & \texttt{\tiny de32592311345e01e242f29888b51b1f4e55e2f78c9b0714391cdfc3fbe6d1b6}\\
{\scriptsize\loc{logs/v9/sheet_ICEx2_extract.txt}} & \texttt{\tiny b2598ce64a6f6847598fe4896004ca6ad15a9a4263b6c7cc0e390e068e91635d}\\
{\scriptsize\loc{intake/new_oct18g/SHA256SUMS}} & \texttt{\tiny 4911a83afa44259453723f4a6b81da4b20fad1462642a0bcf8b3ccfbf9249267}\\
\multicolumn{2}{@{}l}{\textit{Integrity, environment} (7 files)}\\
{\scriptsize\loc{logs/frozen_v9_before.sha}} & \texttt{\tiny cb9c2a11921e376aa078f2081cea58b7c0a8a0363a1fd6976bf7272df89e7b5b}\\
{\scriptsize\loc{logs/v9/tree_hashes_v9.json}} & \texttt{\tiny a2358223bf1146b2f118651cb82aad0c6721553eb0a64f91a3b5a80110a02ea7}\\
{\scriptsize\loc{logs/v9/v8_hashes_before.sha}} & \texttt{\tiny 3dd87be0833964d8d5d9302a2cecca67f0df333d207b829eb57c154cbfe8e5ba}\\
{\scriptsize\loc{logs/v9/v8_hashes_after.sha}} & \texttt{\tiny 3dd87be0833964d8d5d9302a2cecca67f0df333d207b829eb57c154cbfe8e5ba}\\
{\scriptsize\loc{logs/v9/earlier_versions_sha.json}} & \texttt{\tiny 7ea1fb72a1f61af6ec47b14a3e0441d2d13c65eefde6c48494ce519115431399}\\
{\scriptsize\loc{logs/v9/env.json}} & \texttt{\tiny 5633c0d1b2c59fd04977cd8dbe7742a0417a910c10290567be167e4de639c425}\\
{\scriptsize\loc{CHECKPOINT_v9.md}} & \texttt{\tiny cf3b80bfdbb3340176ac0d147d1da6eb43a95238f5dbb550eedd890003e40df9}\\
\end{longtable}}


%======================================================================
\section{v10 artifacts, hashes and rerun commands}
\label{app:v10artifacts}
%======================================================================

Paths are relative to \loc{/workspace/XVI-O1-CONE/}. Long jobs ran
detached with \loc{setsid}, were logged, and were resumable through
per-item JSON caches. Environment: Python~3.13.5, SymPy~1.14.0, mpmath~1.3.0.
Frozen trees before \(=\) after: XVI-P: yes, XVI-Q: yes, XVI-R: yes.
Digests: XVI-P \texttt{ec965c15ae25cd3b\ldots}, XVI-Q \texttt{0df9a0d10409d31a\ldots}, XVI-R \texttt{57187bc6825f5c6a\ldots}.
The v9 source and PDF are unchanged: yes. The v1--v8 sources are unchanged: yes.
The cached \(\varphi\) values (20 files) and the minor evaluations (106 files) are hashed in
\loc{logs/v10/jet/phi/SHA256SUMS} and \loc{logs/v10/minor/SHA256SUMS}. The full manifest is
\loc{logs/v10/artifact_manifest_v10.sha256}; the \TeX{} source and PDF are hashed in \loc{SHA256SUMS_v10}.
The large-\(n\) \(\varphi\) runs were stopped before completion; they feed no printed value.

\paragraph{Rerun commands.}
\begin{verbatim}
cd /workspace/XVI-O1-CONE
python3 code/v10/cr_mode_closed.py
python3 code/v10/f_mode_pslq.py
python3 code/v10/dJ_closed_verify.py
python3 code/v10/lambda_derive.py
setsid nohup bash code/v10/minor_driver.sh \
   > logs/v10/minor/driver.log 2>&1 &
python3 code/v10/minor_entries.py
python3 code/v10/cocycle9.py F_c4 7071/10000 7/10 3/4 \
   4/5 5/6 7/8 9/10 19/20   # F_c3 likewise
setsid nohup bash code/v10/jet_driver.sh > /dev/null 2>&1 &
python3 code/v10/jet_analyse.py; python3 code/v10/r2_scheme.py
python3 code/v10/t1_fp.py; python3 code/v10/t1_fp_check.py
python3 code/v10/jet_density_D04.py
setsid nohup bash code/v10/phi_driver.sh > /dev/null 2>&1 &
setsid nohup bash code/v10/vector_driver.sh > /dev/null 2>&1 &
python3 code/v10/vector_combine.py
python3 code/artifacts_v10.py && python3 code/build_v10.py
pdflatex -halt-on-error XVI-O1-CONE-v10.tex   # twice
\end{verbatim}

{\small
\begin{longtable}{@{}p{.50\linewidth}p{.46\linewidth}@{}}
\caption{SHA-256 of v10 artifacts.}\\
\toprule
File & SHA-256\\
\midrule
\endfirsthead
\toprule
File & SHA-256\\
\midrule
\endhead
\bottomrule
\endfoot
\multicolumn{2}{@{}l}{\textit{Base source (read-only)} (1 files)}\\
{\scriptsize\loc{XVI-O1-CONE-v9.tex}} & \texttt{\tiny 3a8dd43489cf29b96d953a544e5a6758feac16c307449f738ea235ce41c6f1b0}\\
\multicolumn{2}{@{}l}{\textit{v10 programs} (26 files)}\\
{\scriptsize\loc{code/v10/cocycle9.py}} & \texttt{\tiny c6d73ccf092b90a9afda8ba3b65b26ca1a387d746ad57dafc44c8e650c7f7543}\\
{\scriptsize\loc{code/v10/cr_mode_closed.py}} & \texttt{\tiny f9339e76d370390523fc7213f21430cb5b2d4a7ba10f96841069836650ca78df}\\
{\scriptsize\loc{code/v10/dJ_closed_verify.py}} & \texttt{\tiny 26d96b59e31da187cc61c29919aab675fc4885b0364954a9e7a2c5bea51a0690}\\
{\scriptsize\loc{code/v10/f_mode_pslq.py}} & \texttt{\tiny b99327b9686e07428c15037508931f3776acb4b50853eb17f6c0d3a39a8bfe77}\\
{\scriptsize\loc{code/v10/jet_analyse.py}} & \texttt{\tiny f7618580da4559d7fa92e4aeecb5489dd910347cbf937248f538aa5a3528bfe7}\\
{\scriptsize\loc{code/v10/jet_density_D04.py}} & \texttt{\tiny d25ea7a43495f6dfad1ce3b03542ee671d720ada7617ba472ffdcfbd40462bfd}\\
{\scriptsize\loc{code/v10/jet_replay.py}} & \texttt{\tiny 9b93ab0600ed67a6296b8e3374cad4cebf7e9bcf9c328f9cb14664ed454fcfde}\\
{\scriptsize\loc{code/v10/lambda_derive.py}} & \texttt{\tiny 94af1831e985bfdfc288cbed8e27e0d11c3e16b8be7e26e5ffb5e5b485908d98}\\
{\scriptsize\loc{code/v10/minor_entries.py}} & \texttt{\tiny 85655f6a8fd58dd70aa5cfe3ee736e51a5d3932c25804be9612d1f6900ee6043}\\
{\scriptsize\loc{code/v10/minor_poly.py}} & \texttt{\tiny 06456aa92632b8d75ab673850d83baccb856b38c994b632ecf1d49ef2fcbf640}\\
{\scriptsize\loc{code/v10/mixed_r2.py}} & \texttt{\tiny a2776419626cc47100505424cccabd3dfa966c75d4174ac52a461bc699fcb60b}\\
{\scriptsize\loc{code/v10/mixed_r2_check.py}} & \texttt{\tiny 67a2856a2c8e8f2830f558dc6325d99aba81e0c5f50c802fa39924d498fac97d}\\
{\scriptsize\loc{code/v10/phi_fit.py}} & \texttt{\tiny 24ce224a5e9f4079f022f518e31af7061c5657eab361660fef100a80b835094f}\\
{\scriptsize\loc{code/v10/phi_mixed.py}} & \texttt{\tiny b51174ecd2ae1216caf550b0d9e02a31e1c1d7cd634bdeca62227f8ab81f88ef}\\
{\scriptsize\loc{code/v10/r2_scheme.py}} & \texttt{\tiny cd5360f41f048b3ad5c964f15d2fda9aca84e37ff8c2f08df1783a51b257a213}\\
{\scriptsize\loc{code/v10/t1_fp.py}} & \texttt{\tiny 8db8da98a37588fc30854ccfe04a6ed3c59b663108e157b9432c4435f3648a06}\\
{\scriptsize\loc{code/v10/t1_fp_check.py}} & \texttt{\tiny 4fcd4587a4d8a2447cedc2de8a88743232e852f591efed7485d8fd10419fed0d}\\
{\scriptsize\loc{code/v10/vector_combine.py}} & \texttt{\tiny 7c52024edbc55c226720c62e362d4028065fae39035803aff8d9a6771168e280}\\
{\scriptsize\loc{code/v10/vector_fp.py}} & \texttt{\tiny 1748e39a63b7376d0277a19d58717b7886046e7bf75dfe597ae50201f790cd08}\\
{\scriptsize\loc{code/v10/jet_driver.sh}} & \texttt{\tiny 4711578a57f8540d893b633d86decc55f2106523d06fe466881149cba8cd6b88}\\
{\scriptsize\loc{code/v10/minor_driver.sh}} & \texttt{\tiny 5489cf621510e8d23c22e9ae3a31ecc1b35f509f3e6db3664322a8bfdd1621e3}\\
{\scriptsize\loc{code/v10/phi_driver.sh}} & \texttt{\tiny 685ef1dcd00c971fa85ccf04769995ed13d629c0ca5e941c1c68ad8efd0abfda}\\
{\scriptsize\loc{code/v10/phi_driver2.sh}} & \texttt{\tiny d06e442cc222ef8c43cc8d750956c062e28d8d16140e805e3bcbda24e665b163}\\
{\scriptsize\loc{code/v10/vector_driver.sh}} & \texttt{\tiny dee6090a4b399f9193563d35e942830728c95b992a42a22a8bfb2175c955ef86}\\
{\scriptsize\loc{code/build_v10.py}} & \texttt{\tiny 75cb95c20748413a366029c40b9f0280cc90a425981fc304b80a7b7847e6c7f4}\\
{\scriptsize\loc{code/artifacts_v10.py}} & \texttt{\tiny 8de0cad6e1d9ad59cee1ee0b640ffdca7dcc9385e0c0660f305a8dcf0429932b}\\
\multicolumn{2}{@{}l}{\textit{v10 source fragments} (4 files)}\\
{\scriptsize\loc{v10src/abstract.tex}} & \texttt{\tiny 6d4a42db29b130adfc18702497cf6eef0e8d654cc821c08d1549ea5cf7be0db3}\\
{\scriptsize\loc{v10src/bib_v10.tex}} & \texttt{\tiny c648846e7462ed90a7a0907bc31bc748f7e95dd0358373e8342bf906372ebb7a}\\
{\scriptsize\loc{v10src/opentable.tex}} & \texttt{\tiny 060a3ade76c775984acbf220242f76498caca62a66829780e87951d6b7c316a6}\\
{\scriptsize\loc{v10src/sec_v10.tex}} & \texttt{\tiny 770e85953af1f47e47a387ffcf22156c4e2c737075530616b68a76555c63af4a}\\
\multicolumn{2}{@{}l}{\textit{Delta J closed form} (11 files)}\\
{\scriptsize\loc{logs/v10/cr_mode_closed.json}} & \texttt{\tiny 939528c5aaf52273d12f32a570dc6ceed917a0ceb21988848f59c3fc5e265ad8}\\
{\scriptsize\loc{logs/v10/dJ_closed_verify.json}} & \texttt{\tiny 5159c9163551c83a17161d1814e7dff061d9d3448cd65a0f88e8712d036a6a59}\\
{\scriptsize\loc{logs/v10/fmode/3_2.json}} & \texttt{\tiny 25daa838dce8ce663f92ab61cdc38b49b5f2b6f61355977416734d1647222181}\\
{\scriptsize\loc{logs/v10/fmode/4_3.json}} & \texttt{\tiny aca72774baf4e99f3855d1e5c2ac85118b7c50c7274f5f9c16e3e13f7dad9017}\\
{\scriptsize\loc{logs/v10/fmode/5_3.json}} & \texttt{\tiny 9dbd7ad598b5444481965b76ea81cf638c726376a074bc7ab57ea3a9ae2679cf}\\
{\scriptsize\loc{logs/v10/fmode/5_4.json}} & \texttt{\tiny ab984d9a4a13b211ece98e64fc6be0b03b89cdefac36cd365ee9e0eb8ae540f6}\\
{\scriptsize\loc{logs/v10/fmode/6_5.json}} & \texttt{\tiny ce0cc18323063200588745de47548e0aaa952dbee89ad04b0f1eb808475190f2}\\
{\scriptsize\loc{logs/v10/fmode/7_4.json}} & \texttt{\tiny 37e213fb00f05380f2d11d1c28d9711f75249261cdf9b86c76cd2c7f6becd092}\\
{\scriptsize\loc{logs/v10/fmode/7_5.json}} & \texttt{\tiny 445b722bd9d3f3729ccee37269bbe561386cff57c4be70da19217c3dae2d4945}\\
{\scriptsize\loc{logs/v10/fmode/8_5.json}} & \texttt{\tiny 4b464ff79c117543f91134fab5b185d5700a1fed5f10e941fad5397689bd5360}\\
{\scriptsize\loc{logs/v10/fmode/9_5.json}} & \texttt{\tiny cb11bbd8860cbf9a186a701d2d870900f7a706ed7eb0eb71e106d50788ba2aef}\\
\multicolumn{2}{@{}l}{\textit{Ghost normalizations} (1 files)}\\
{\scriptsize\loc{logs/v10/lambda_derive.json}} & \texttt{\tiny b15ccf4c87d1d0dc0a27fea25af189ee441684f0ad06030f46ce24a7832decaf}\\
\multicolumn{2}{@{}l}{\textit{Chamber minors} (19 files)}\\
{\scriptsize\loc{logs/v10/minor/CR_c1/pick.json}} & \texttt{\tiny ef8ed9fdf3a5b5bfb1d57321bf2e431b6d0a058c62046811fbc29c2f0e86dc13}\\
{\scriptsize\loc{logs/v10/minor/CR_c2/pick.json}} & \texttt{\tiny d6b21d7a8209f9e3ce179b7dbe1fbc20a3d8c4f374c9ff951ae27ecad084596f}\\
{\scriptsize\loc{logs/v10/minor/CR_c3/pick.json}} & \texttt{\tiny a9fbb9bb784a70b5818a6423122f18877cb6719c4c0f6427c6d18e6dd73c8448}\\
{\scriptsize\loc{logs/v10/minor/CR_c4/pick.json}} & \texttt{\tiny ff43cbf02a89e2081ce558e7c2b2c9d3f403c6621383926827a4ca0ae51772c6}\\
{\scriptsize\loc{logs/v10/minor/F_c1/pick.json}} & \texttt{\tiny bd008573b6d427d752ae6136c5ad44f721ccd35f6315c21e85b44c597c003f67}\\
{\scriptsize\loc{logs/v10/minor/F_c2/pick.json}} & \texttt{\tiny ade2b067e5fe3da91fc1ececd3ede26d86c6857f47ca814eb41c74d8a894b162}\\
{\scriptsize\loc{logs/v10/minor/F_c3/pick.json}} & \texttt{\tiny cb90901af66560cbdd497ceb8b0077ca7af89e4abd4e9cff54476d39e4851a2a}\\
{\scriptsize\loc{logs/v10/minor/F_c4/pick.json}} & \texttt{\tiny a4af66a3b268f73ac168bb4613f584ba818ccf94087f448a015bcb4740d22c4b}\\
{\scriptsize\loc{logs/v10/minor/CR_c1/entries_check.json}} & \texttt{\tiny 761c25903f00a6612b18e786ee39f7e674e68abfeea8ab9162fcef35211e7c8f}\\
{\scriptsize\loc{logs/v10/minor/CR_c2/entries_check.json}} & \texttt{\tiny 348bc38ecc754a39a22dab2434f9dfa8f9508a4d967659b6e5e101fe4ae61385}\\
{\scriptsize\loc{logs/v10/minor/CR_c3/entries_check.json}} & \texttt{\tiny 1c4cae65bd188f4015c057a43ee109481089a79173cc270743a9997275f1b516}\\
{\scriptsize\loc{logs/v10/minor/CR_c4/entries_check.json}} & \texttt{\tiny df135b94bd6e67cff2391c7cc666d80e138925aee88341abd72edbe25f7e606b}\\
{\scriptsize\loc{logs/v10/minor/F_c1/entries_check.json}} & \texttt{\tiny aa2e6d90dd5e394a2b71bdf02aa82b8aafb361c9de26dd9188f7d250b35c24eb}\\
{\scriptsize\loc{logs/v10/minor/F_c2/entries_check.json}} & \texttt{\tiny 807c5572509ec6966b278ba622e105b56cc0777b7cbd912cb03badaaf227e9db}\\
{\scriptsize\loc{logs/v10/minor/F_c3/entries_check.json}} & \texttt{\tiny f8fb788a64060d68b0a3c0b5123f5b26ea03cab8781715ed6ff6fbd3ff2d9f73}\\
{\scriptsize\loc{logs/v10/minor/F_c4/entries_check.json}} & \texttt{\tiny a310782e9ef118d07c13f07d7dcc2b36142ab5e33726f2fed036fd51464dfeeb}\\
{\scriptsize\loc{logs/v10/minor/F_c3/cocycle9.json}} & \texttt{\tiny 377c62498a429585ea9f1b6629697fb0d9c9114f3344c211758df947165644d9}\\
{\scriptsize\loc{logs/v10/minor/F_c4/cocycle9.json}} & \texttt{\tiny fbf314449c9bc32240477f4427087165efbcc7df931c4f29f1da60f099750a8d}\\
{\scriptsize\loc{logs/v10/minor/SHA256SUMS}} & \texttt{\tiny ed4653f9ecb3ac7c2443db90827bb1f73af34e6ed720b8e0caad22182311ef6d}\\
\multicolumn{2}{@{}l}{\textit{Transverse jets} (17 files)}\\
{\scriptsize\loc{logs/v10/jet/jet_D01.json}} & \texttt{\tiny ffed1138c1f2fb5b554c865593a5e48da10d19ee6900922f17d21538a21457df}\\
{\scriptsize\loc{logs/v10/jet/jet_D02.json}} & \texttt{\tiny 35643ee8d7aa9459d629a8d19d8ab38318d6ce68a39ddc643ed8f9fdd874e11f}\\
{\scriptsize\loc{logs/v10/jet/jet_D03.json}} & \texttt{\tiny 531ba4cfa9741a5d999c19da05c1b09c453f6c61c8d0f7cceb0c91ca88caff00}\\
{\scriptsize\loc{logs/v10/jet/jet_D04.json}} & \texttt{\tiny 4c3d5a5460806bb341648bf9b9f367fa2faeb70837294238a80d8f38145a8d89}\\
{\scriptsize\loc{logs/v10/jet/jet_D05a.json}} & \texttt{\tiny c695b8be7cdf20cb2d4d71ccfb524bffc5bf5a531056de7e959041736ab738bc}\\
{\scriptsize\loc{logs/v10/jet/jet_D06.json}} & \texttt{\tiny 04cebbf0bdc3eb40a9e254a8ab430026ff3b3b88773bbec70efcd00224adda6a}\\
{\scriptsize\loc{logs/v10/jet/jet_D11.json}} & \texttt{\tiny 7b68328a09b92eb169da152c5e9557be7e09f923798c05f50efd8ef9524ff7b6}\\
{\scriptsize\loc{logs/v10/jet/jet_D12.json}} & \texttt{\tiny 934e721f7e9ff41029c74b299aeba4e83960d42e02097a29e1bfeb64881d44ef}\\
{\scriptsize\loc{logs/v10/jet/jet_D13.json}} & \texttt{\tiny f856bcfff7707c4ca5eccbf26aac233c79e7597aead51b1acec8d7dc7249cc88}\\
{\scriptsize\loc{logs/v10/jet/jet_analysis.json}} & \texttt{\tiny 5daf13a48ff23cf576f8f4b8f1d02d7524ea52163b14dc3d08615103bf0fad08}\\
{\scriptsize\loc{logs/v10/jet/r2_scheme.json}} & \texttt{\tiny 7f6492a3ee94d407ddcaa10fffa39b6d52462aa699274decab531a24b4112131}\\
{\scriptsize\loc{logs/v10/jet/t1_fp_check.json}} & \texttt{\tiny 5beac73b57965a11d7ebe68397bb32d1cbbeb84a5866955f53c24fb764b4d168}\\
{\scriptsize\loc{logs/v10/jet/t1_fp_same.json}} & \texttt{\tiny 9b9385cdae53f49af5e845e24e1003cfa3d0e9e3e4c9f1655d6e0dcd4f6bed8b}\\
{\scriptsize\loc{logs/v10/jet/mixed_r2.json}} & \texttt{\tiny b17bd5d3f8788e573c7c4e9cd5b8c283746a129fb18a4d200cbcc7844bd03675}\\
{\scriptsize\loc{logs/v10/jet/mixed_r2_check.json}} & \texttt{\tiny 3416ad37869d05fcbc93748e3cb1b8037bc6f72356900dcdbd66b99288e4997b}\\
{\scriptsize\loc{logs/v10/jet/density_D04.json}} & \texttt{\tiny f8e9f35864582aec6338493bd81b1c6e66308cac09d0f76a77f03f9e8b14336f}\\
{\scriptsize\loc{logs/v10/jet/phi/SHA256SUMS}} & \texttt{\tiny a65d81990623d6593aa653c3e4d7219596865e45245fed47ccc2df874a7b8c41}\\
\multicolumn{2}{@{}l}{\textit{Vector finite parts} (5 files)}\\
{\scriptsize\loc{logs/v10/vector_fp/3_10.json}} & \texttt{\tiny 911f52f839f4be0cd2f80bb7b2c2df76b31810accf3727e3927f5087883ebcc4}\\
{\scriptsize\loc{logs/v10/vector_fp/3_5.json}} & \texttt{\tiny c138cb911e3e051a2b1013eb4c5273f3af486875c28277f50538621fc6be60bb}\\
{\scriptsize\loc{logs/v10/vector_fp/41_100.json}} & \texttt{\tiny 43619f5aa07f5362d8d9531b2c6405d16b25c7fc7887259a07711e13ed21ce17}\\
{\scriptsize\loc{logs/v10/vector_fp/7071_10000.json}} & \texttt{\tiny 2fb16311db70b19ce6b1cba7646c9263ec444348340c256e4294e345de7c791f}\\
{\scriptsize\loc{logs/v10/vector_fp/combined.json}} & \texttt{\tiny 920a24127826fda08054402b1255615c7f7bc31c0a63a7fba7c7f020fafc24c1}\\
\multicolumn{2}{@{}l}{\textit{Integrity, environment} (6 files)}\\
{\scriptsize\loc{logs/frozen_v10_before.sha}} & \texttt{\tiny cb9c2a11921e376aa078f2081cea58b7c0a8a0363a1fd6976bf7272df89e7b5b}\\
{\scriptsize\loc{logs/v10/tree_hashes_v10.json}} & \texttt{\tiny ddf510d823e0af9de28430bdf7be42309eaab04b1d6b9a6c87b39b8e9735626d}\\
{\scriptsize\loc{logs/v9_hashes_before_v10.sha}} & \texttt{\tiny 97840e3bd1107e61deb848ac53de97a8aa68275b16517f2b17579b465df2dc74}\\
{\scriptsize\loc{logs/v10/earlier_versions_sha.json}} & \texttt{\tiny c9d46ce4bbae98a1c22a08a7f91ba24f3ac242f054ee673fbe6d0a7286bbc43c}\\
{\scriptsize\loc{logs/v10/env.json}} & \texttt{\tiny f3c3b679dd87e7e87a4a5670b7e7dede9a62e737459a7a2a21d8fc562df914fa}\\
{\scriptsize\loc{CHECKPOINT_v10.md}} & \texttt{\tiny 4c01e1e654d8f96797988e7208517f82e7259ce2d5bf774b07b863924c3ee04c}\\
\end{longtable}}


%======================================================================
\section{v11 artifacts, hashes and rerun commands}
\label{app:v11artifacts}
%======================================================================

Paths are relative to \loc{/workspace/XVI-O1-CONE/}. Long jobs ran
detached with \loc{setsid}, were logged, and were resumable through
per-item JSON caches. Environment: Python~3.13.5, SymPy~1.14.0, mpmath~1.3.0.
Frozen trees before \(=\) after: XVI-P: yes, XVI-Q: yes, XVI-R: yes.
Digests: XVI-P \texttt{ec965c15ae25cd3b\ldots}, XVI-Q \texttt{0df9a0d10409d31a\ldots}, XVI-R \texttt{57187bc6825f5c6a\ldots}.
The v9 and v10 sources and PDFs are unchanged: yes. The v1--v8 sources are unchanged: yes.
The cached \(F(\mu,\nu)\) values (404 files) are hashed in \loc{logs/v11/M10/cache_SHA256SUMS}; the
superseded first run (Section~\ref{sec:v11bug}) is kept in \loc{logs/v11/M10/cache_superseded/} and feeds no printed value.
The full manifest is \loc{logs/v11/artifact_manifest_v11.sha256}; the \TeX{} source and PDF are hashed in \loc{SHA256SUMS_v11}.

\emph{Manifest fix.} The v10 manifest hashed \loc{CHECKPOINT_v10.md} one append before its final state
(\texttt{4c01e1e654d8f967\ldots}); its current hash is
\texttt{8bb24346c6f7d991\ldots} (\loc{logs/v11/checkpoint_v10_fix.json}). In v11 the last
append to \loc{CHECKPOINT_v11.md} precedes this hashing, and nothing is appended afterwards.

\paragraph{Rerun commands.}
\begin{verbatim}
cd /workspace/XVI-O1-CONE
python3 code/v11/degree_check.py; python3 code/v11/degree_crosscheck.py
python3 code/v11/fmode_reduce.py; python3 code/v11/fmode_derive.py
python3 code/v11/m10_jobs.py > logs/v11/M10/jobs.txt
setsid nohup bash code/v11/m10_driver.sh > /dev/null 2>&1 &
python3 code/v11/m10_sum.py
python3 code/artifacts_v11.py && python3 code/build_v11.py
pdflatex -halt-on-error XVI-O1-CONE-v11.tex   # until stable
\end{verbatim}

{\small
\begin{longtable}{@{}p{.50\linewidth}p{.46\linewidth}@{}}
\caption{SHA-256 of v11 artifacts.}\\
\toprule
File & SHA-256\\
\midrule
\endfirsthead
\toprule
File & SHA-256\\
\midrule
\endhead
\bottomrule
\endfoot
\multicolumn{2}{@{}l}{\textit{Base source (read-only)} (1 files)}\\
{\scriptsize\loc{XVI-O1-CONE-v10.tex}} & \texttt{\tiny d21c8d84b12b53e50a78227d86dbff204293748d7712ced02e11b08b983b8d25}\\
\multicolumn{2}{@{}l}{\textit{v11 programs} (12 files)}\\
{\scriptsize\loc{code/v11/degree_check.py}} & \texttt{\tiny c5c0c48ef4e5f2750235f8da59daa62e157d177250c556b275b6dfbfa93d7362}\\
{\scriptsize\loc{code/v11/degree_crosscheck.py}} & \texttt{\tiny 09139a504eaf3af98adeacfa3a8ae344455cbf6c133b58b4604d2e73e47d505e}\\
{\scriptsize\loc{code/v11/fmode_derive.py}} & \texttt{\tiny d16e4f844f54237fde609d74360fb7d81a1169e98b71cfbe96cf378423866a5f}\\
{\scriptsize\loc{code/v11/fmode_reduce.py}} & \texttt{\tiny 1bcf052a6baf0894b8b14cdd765d88a4933314cd8f595f26c86212c0f45b839c}\\
{\scriptsize\loc{code/v11/m10_eval.py}} & \texttt{\tiny d01c93f6689f8844efcf5acfca47f7ccaef33325b749b916e96ecb3f60ef904c}\\
{\scriptsize\loc{code/v11/m10_jobs.py}} & \texttt{\tiny 43f70dfbc390768faaf2573f5c8f4575b3373ea83ac2cca06e4431466e083b25}\\
{\scriptsize\loc{code/v11/m10_post.py}} & \texttt{\tiny 5b8a6aab7f5d46d4e913b70ae6d93702911326c99c6a31416c290b987189b7d9}\\
{\scriptsize\loc{code/v11/m10_sum.py}} & \texttt{\tiny 1595ff5aa5ed8e3fe300e71ffaf4514c680cdd3ea44e00ba0ece13c33d5abce3}\\
{\scriptsize\loc{code/v11/t1v11.py}} & \texttt{\tiny 4660ac0960ab9dfe2e92a01ce2c7900aabb9c295834e60306965f4abdb167dd1}\\
{\scriptsize\loc{code/v11/m10_driver.sh}} & \texttt{\tiny 627f6ff0aa57b4e6aed21f248af15e5e5dec784eea05f9105377da2558756596}\\
{\scriptsize\loc{code/build_v11.py}} & \texttt{\tiny a71bde039d87789829333866ab7a3bf31db9d1fa966dfd1500936e629c7bd3fc}\\
{\scriptsize\loc{code/artifacts_v11.py}} & \texttt{\tiny 8645d9bbfd275047c2a149914f9efc5f94511284f1147916d79a91a3e15385ae}\\
\multicolumn{2}{@{}l}{\textit{v11 source fragments} (4 files)}\\
{\scriptsize\loc{v11src/abstract.tex}} & \texttt{\tiny 82c1899a7e301fab8022a089f012d160b59576044a6294d86f0918b588453435}\\
{\scriptsize\loc{v11src/bib_v11.tex}} & \texttt{\tiny 60374aa56074952db5b9c95fb5965153933e35644ab481052613996402eaa4f9}\\
{\scriptsize\loc{v11src/opentable.tex}} & \texttt{\tiny 7451352004052c4ef79a1736ea0872691ea47162ac07312215aefec1e0655205}\\
{\scriptsize\loc{v11src/sec_v11.tex}} & \texttt{\tiny 7431c14dfd2eb064a6738b10fbe11f65cd7a9190a25bed923f0414663add9685}\\
\multicolumn{2}{@{}l}{\textit{F mode derivation} (2 files)}\\
{\scriptsize\loc{logs/v11/fmode_reduce.json}} & \texttt{\tiny 6289543dec1700c1b7944344756742677dc6a7e21cd417ed9f03de1e76a46dbc}\\
{\scriptsize\loc{logs/v11/fmode_derive.json}} & \texttt{\tiny 9b897c1cdd37263e49544181583dc633f2c9c3e365df8da26b0065a3c26402f6}\\
\multicolumn{2}{@{}l}{\textit{Degree bound (D)} (2 files)}\\
{\scriptsize\loc{logs/v11/degree_check.json}} & \texttt{\tiny 229f265c1dd566a0b02ccf5173434762b70ae03f162a162680bf45dbe299d35a}\\
{\scriptsize\loc{logs/v11/degree_crosscheck.json}} & \texttt{\tiny 2901c5d65f674a6a5124a8f600061bd63e0a1efb17ebfad5d943fdef9c985641}\\
\multicolumn{2}{@{}l}{\textit{Mode sum M10} (9 files)}\\
{\scriptsize\loc{logs/v11/M10/jobs.txt}} & \texttt{\tiny 5378b8e118373db999c3f5611d7054ef425cc2c5e6a875a7b1178b8034010829}\\
{\scriptsize\loc{logs/v11/M10/m10.json}} & \texttt{\tiny f58b460f79bc0fbfbbf7196be0e923585586bdaa47ad67058684c0651071980b}\\
{\scriptsize\loc{logs/v11/M10/cache_SHA256SUMS}} & \texttt{\tiny 4e63c1d913104cd121a044955d1dd4670371ffb88c5ae4f64ee0ec10e1818c71}\\
{\scriptsize\loc{logs/v11/M10/phi_corrected.json}} & \texttt{\tiny cd4cc6653e3ef9b4385e4f8d8323efa578ebfe71055c70e1984167424ef06558}\\
{\scriptsize\loc{logs/v11/M10/m10_text.tex}} & \texttt{\tiny a1a666abe117a79edc09c2ed15c5625a8e3e356b52e89345dfafbc9e2f75c3ee}\\
{\scriptsize\loc{logs/v11/M10/jobs_dedup.txt}} & \texttt{\tiny 72b5560949a87642c8a44ebb56044881a7c31961cbfc6bdff99f0fc88560c561}\\
{\scriptsize\loc{logs/v11/M10/m10_leading.json}} & \texttt{\tiny 173abaf669a892c4d963bb16080299f95c9f7a55a46ce96c878b04ff85b9e7ad}\\
{\scriptsize\loc{logs/v11/M10/m10_pslq.json}} & \texttt{\tiny 2e35be6bc699295dc53dda73fb79bec1efcb4dabf51b2d8b49333b3f362627e4}\\
{\scriptsize\loc{logs/v11/M10/m10_sanity.json}} & \texttt{\tiny 081018c7db3a61b81ed5ac4e47ece9f3e490a1cc49e926dd70872cc4550e7e49}\\
\multicolumn{2}{@{}l}{\textit{Registry} (1 files)}\\
{\scriptsize\loc{logs/v11/registry_oct19.sha}} & \texttt{\tiny ba94230ae35fbb1c4456568f0c3787d01310ef19195d5373d5821c5309d91538}\\
\multicolumn{2}{@{}l}{\textit{Integrity, environment} (7 files)}\\
{\scriptsize\loc{logs/frozen_v11_before.sha}} & \texttt{\tiny cb9c2a11921e376aa078f2081cea58b7c0a8a0363a1fd6976bf7272df89e7b5b}\\
{\scriptsize\loc{logs/v11/tree_hashes_v11.json}} & \texttt{\tiny ddf510d823e0af9de28430bdf7be42309eaab04b1d6b9a6c87b39b8e9735626d}\\
{\scriptsize\loc{logs/v10_hashes_before_v11.sha}} & \texttt{\tiny 1ca9d9369867da38fee1e2458b4e2d94fe277505ab39c3f84392c7b7985ee63a}\\
{\scriptsize\loc{logs/v11/earlier_versions_sha.json}} & \texttt{\tiny 111da13ad1339049d8fcaeaecf12373f70553db4f0cbe749e707a77aef0e7523}\\
{\scriptsize\loc{logs/v11/checkpoint_v10_fix.json}} & \texttt{\tiny cd7ab2d3abbebc904e1d32fbdfa78f0e4418ab42884ceba36fda28abb86c6860}\\
{\scriptsize\loc{logs/v11/env.json}} & \texttt{\tiny f3c3b679dd87e7e87a4a5670b7e7dede9a62e737459a7a2a21d8fc562df914fa}\\
{\scriptsize\loc{CHECKPOINT_v11.md}} & \texttt{\tiny 0bdb0d051f8b5ad8acb8effd0f8a5c357d86b2357a6fd3738eb1044cb530447c}\\
\end{longtable}}


%======================================================================
\section{v12 artifacts, hashes and rerun commands}
\label{app:v12artifacts}
%======================================================================

Paths are relative to \loc{/workspace/XVI-O1-CONE/}. Long jobs ran
detached with \loc{setsid} and were logged. Environment: Python~3.13.5, SymPy~1.14.0, mpmath~1.3.0.
Frozen trees before \(=\) after: XVI-P: yes, XVI-Q: yes, XVI-R: yes.
Digests: XVI-P \texttt{ec965c15ae25cd3b\ldots}, XVI-Q \texttt{0df9a0d10409d31a\ldots}, XVI-R \texttt{57187bc6825f5c6a\ldots}.
The v9, v10 and v11 sources and PDFs are unchanged: yes. The v1--v8 sources are unchanged: yes.
The full manifest is \loc{logs/v12/artifact_manifest_v12.sha256}; the \TeX{} source and PDF are hashed in \loc{SHA256SUMS_v12}.
The last append to \loc{CHECKPOINT_v12.md} precedes this hashing.

\paragraph{Rerun commands.}
\begin{verbatim}
cd /workspace/XVI-O1-CONE
python3 code/v12/phi_asym_fit.py; python3 code/v12/structure_test.py
python3 code/v12/F_closed.py
python3 code/v12/m10_exact.py 40; python3 code/v12/counterterm.py
python3 code/v12/m10_pslq.py
setsid nohup bash code/v12/walls_driver.sh > /dev/null 2>&1 &
python3 code/v12/make_tables.py
python3 code/artifacts_v12.py && python3 code/build_v12.py
pdflatex -halt-on-error XVI-O1-CONE-v12.tex   # until stable
\end{verbatim}

{\small
\begin{longtable}{@{}p{.50\linewidth}p{.46\linewidth}@{}}
\caption{SHA-256 of v12 artifacts.}\\
\toprule
File & SHA-256\\
\midrule
\endfirsthead
\toprule
File & SHA-256\\
\midrule
\endhead
\bottomrule
\endfoot
\multicolumn{2}{@{}l}{\textit{Base source (read-only)} (1 files)}\\
{\scriptsize\loc{XVI-O1-CONE-v11.tex}} & \texttt{\tiny 863bb08efc04416464d5628041d98ec7478015315b4c804cacc6b260181f203e}\\
\multicolumn{2}{@{}l}{\textit{v12 programs} (11 files)}\\
{\scriptsize\loc{code/v12/F_closed.py}} & \texttt{\tiny 0456cb703049723257452addde7e3783799d113deca58085b661b7a0c419ce00}\\
{\scriptsize\loc{code/v12/counterterm.py}} & \texttt{\tiny 7342f223780372c078030725bc1d11960f119b4615c340fb72c2e24ab3ce22bc}\\
{\scriptsize\loc{code/v12/m10_exact.py}} & \texttt{\tiny 202e65ba9d96c8b93b91a305061e845cc0a03e006e1a9525bcb242076876d773}\\
{\scriptsize\loc{code/v12/m10_pslq.py}} & \texttt{\tiny 74df977c35cc874c14b09ac4413ea56a4c22f745daea393afc140a02797de7e4}\\
{\scriptsize\loc{code/v12/make_tables.py}} & \texttt{\tiny c34407b1fdce5641ce96b36bb878e0534c0cc82ad4371e8cbdb79c4fdce16e53}\\
{\scriptsize\loc{code/v12/phi_asym_fit.py}} & \texttt{\tiny 7fc2ba6ecddf002e4c7b5e5536014a8b7aad9c082a61ebe53ec7e86ca4923037}\\
{\scriptsize\loc{code/v12/structure_test.py}} & \texttt{\tiny 93ed4137164b4514f64c991e3d254cb53c6c5757b689f1b5312bacfd7ab1aa8c}\\
{\scriptsize\loc{code/v12/wall_certificate.py}} & \texttt{\tiny 3f6e849b3abc9e9a838511693944ba7cb74914730d6c658189b10ae9bec5609d}\\
{\scriptsize\loc{code/v12/walls_driver.sh}} & \texttt{\tiny 29de260c64a06ea0d926f0f74904eb90f8ea2dd0f057ea2656ed5e09a7dd5c7b}\\
{\scriptsize\loc{code/build_v12.py}} & \texttt{\tiny aaeddb313b8530955ccb27be7eb4d869a45969ff44cf19a5647bf2b42906fb4d}\\
{\scriptsize\loc{code/artifacts_v12.py}} & \texttt{\tiny 1a14fb8922da23c875b74b57db326beae92a6200d466284ee1dfe3309ea5c134}\\
\multicolumn{2}{@{}l}{\textit{v12 source fragments} (10 files)}\\
{\scriptsize\loc{v12src/abstract.tex}} & \texttt{\tiny 35a9eaa1bee58b57325ddd2ced0a1766b71034a8fe7b65caad11b314961ccf3e}\\
{\scriptsize\loc{v12src/bib_v12.tex}} & \texttt{\tiny b50296bc16b71c6f30c0e46371fd6e0b51342d3899c4e0447c18b92c3dcc1d01}\\
{\scriptsize\loc{v12src/m10_table.tex}} & \texttt{\tiny b53cd6e0767926f7d17c0dcca0cfc4b675a78746c0c3af1467c57fd134f9940d}\\
{\scriptsize\loc{v12src/opentable.tex}} & \texttt{\tiny d8c870ee5a528d43f1b9a8f38394b718246d8849396343bc70df294857eb164f}\\
{\scriptsize\loc{v12src/release_items.tex}} & \texttt{\tiny 02a0ceb5cb2863352d4851f12f522f8904f9a8fb14345c566f25ee264c6fb3ad}\\
{\scriptsize\loc{v12src/scope.tex}} & \texttt{\tiny 6cc18e30fd344499325571b3e2ef2cbbdc217c0d113747a5f7eff80c5dd16062}\\
{\scriptsize\loc{v12src/sec_v12.tex}} & \texttt{\tiny b037e2e87c0f55ffb24335c0c7a4a8684e2ebfaa5c04f2699829b3b84cda8bf5}\\
{\scriptsize\loc{v12src/subtitle.tex}} & \texttt{\tiny af9a7488e1b84118cab4b33c11275717ef6c147f75fa36e630cc95e24bfbb284}\\
{\scriptsize\loc{v12src/walls_table.tex}} & \texttt{\tiny 69a47164e45e335d8fd835271391c8a0ab58dcd12ee4070d7999b25a70abcd21}\\
{\scriptsize\loc{v12src/walls_text.tex}} & \texttt{\tiny 949dbad49645b19e8788aa09057247367a7915c3e402c371fb6c1aa5ae17c65b}\\
\multicolumn{2}{@{}l}{\textit{Closed form and M10} (5 files)}\\
{\scriptsize\loc{logs/v12/F_closed.json}} & \texttt{\tiny 52555dbfa415868a37e5d028282ddefd4bba54b887256c0f7a830dcf36ae31de}\\
{\scriptsize\loc{logs/v12/structure_test.json}} & \texttt{\tiny 1d547e416e764d165eb735eccb5108acd4ce1ef290a677b61a53b9682ae94d6d}\\
{\scriptsize\loc{logs/v12/m10_exact_dps40.json}} & \texttt{\tiny e3289ebaebf2f4489950862d17a57820aea4bfae3adaa77b9fb5fac0d42f2e27}\\
{\scriptsize\loc{logs/v12/counterterm.json}} & \texttt{\tiny edcb0e4f321a645b8d66ae7570305a3322e27bcedb13959f15537572841b4af0}\\
{\scriptsize\loc{logs/v12/m10_pslq.json}} & \texttt{\tiny 45975a515bc74bb28632ac648909dd8af68129880056a2c083790aa6e7f9f308}\\
\multicolumn{2}{@{}l}{\textit{Wall certificates} (19 files)}\\
{\scriptsize\loc{logs/v12/walls/apex_cert_1_2_CR.json}} & \texttt{\tiny 1bd0bc1bb1b437d1f76b90328b672f79a3bcda95843cf4fc9cda565e058b4e57}\\
{\scriptsize\loc{logs/v12/walls/apex_cert_1_2_CR_vectors.json}} & \texttt{\tiny 801d520ac941843d8d728449035e904740c5b200dbd715898544a16af8a6ef14}\\
{\scriptsize\loc{logs/v12/walls/apex_cert_1_2_F.json}} & \texttt{\tiny c64cf4905b4234b3d083c4b0b8a97509950d34ae49b87348fc2be9ff6dbb57a8}\\
{\scriptsize\loc{logs/v12/walls/apex_cert_1_2_F_vectors.json}} & \texttt{\tiny 801d520ac941843d8d728449035e904740c5b200dbd715898544a16af8a6ef14}\\
{\scriptsize\loc{logs/v12/walls/apex_cert_1_3_CR.json}} & \texttt{\tiny fc1ac350a8b2b6cff9bf59f6bc549a8df90c6683373b0627c518b051b10c62e3}\\
{\scriptsize\loc{logs/v12/walls/apex_cert_1_3_CR_vectors.json}} & \texttt{\tiny 5a822527a80a1825993c7e1b506820040e94a73cdce924a3958281e35318f304}\\
{\scriptsize\loc{logs/v12/walls/apex_cert_1_3_F.json}} & \texttt{\tiny a1936f1bfc569d6884c5d4924d7735c9f6e3eded4473f1f5cb602e8aa4a1412a}\\
{\scriptsize\loc{logs/v12/walls/apex_cert_1_3_F_vectors.json}} & \texttt{\tiny 5a822527a80a1825993c7e1b506820040e94a73cdce924a3958281e35318f304}\\
{\scriptsize\loc{logs/v12/walls/apex_cert_2_3_CR.json}} & \texttt{\tiny a8ef79e8533f9e701dbdce05c852bb55febdeb5b1793e3ffe7c44c1c7b2a87d4}\\
{\scriptsize\loc{logs/v12/walls/apex_cert_2_3_CR_vectors.json}} & \texttt{\tiny 5a822527a80a1825993c7e1b506820040e94a73cdce924a3958281e35318f304}\\
{\scriptsize\loc{logs/v12/walls/apex_cert_2_3_F.json}} & \texttt{\tiny 9c91707d46277ebc95d7cdd15f80361929ce522cd2e9ff73d7c68f0126e6f983}\\
{\scriptsize\loc{logs/v12/walls/apex_cert_2_3_F_vectors.json}} & \texttt{\tiny 9a60468c4977ac366c0a0154237887c9ec10a5f03a5802dcbd1d1e6493cfb590}\\
{\scriptsize\loc{logs/v12/walls/driver.log}} & \texttt{\tiny 5c91bc5fb3c4c9ab3cc8ceee39df4d1ff7bb892e5eb34ed8f809873d8f1dd969}\\
{\scriptsize\loc{logs/v12/walls/run_1_2_CR.log}} & \texttt{\tiny 6743f05ef33f0e5533c6658d58c844bd0aa482e82945900266ebd70e9808ceb2}\\
{\scriptsize\loc{logs/v12/walls/run_1_2_F.log}} & \texttt{\tiny e6ccb77c4ae29bb0bb3bafd306c0558b5945b8e76b45184cd7f71106f3b53cd4}\\
{\scriptsize\loc{logs/v12/walls/run_1_3_CR.log}} & \texttt{\tiny eacceb0421c58b556a319afbb873531c2f58299023d02520cd98c5965c14ea0e}\\
{\scriptsize\loc{logs/v12/walls/run_1_3_F.log}} & \texttt{\tiny 1a262adc22c9a96f7030764301f02860a7e2cbf88f46d1ef49cce04436035cdc}\\
{\scriptsize\loc{logs/v12/walls/run_2_3_CR.log}} & \texttt{\tiny 62622c79c276713f0b69f8c498df5ef80592b7897678f2a311e1ec1a5e2ca1ea}\\
{\scriptsize\loc{logs/v12/walls/run_2_3_F.log}} & \texttt{\tiny e0b50e18a7cde49e044e7e09403a2bf7e97e82aca92cd696ea1160befafee5cf}\\
\multicolumn{2}{@{}l}{\textit{Registry} (1 files)}\\
{\scriptsize\loc{logs/v12/registry_oct19b.sha}} & \texttt{\tiny 828d9afa530fa1dd0628ee751730b88bc2645a01f41adee08f0561e8d5e777e3}\\
\multicolumn{2}{@{}l}{\textit{Integrity, environment} (6 files)}\\
{\scriptsize\loc{logs/frozen_v12_before.sha}} & \texttt{\tiny cb9c2a11921e376aa078f2081cea58b7c0a8a0363a1fd6976bf7272df89e7b5b}\\
{\scriptsize\loc{logs/v12/tree_hashes_v12.json}} & \texttt{\tiny ddf510d823e0af9de28430bdf7be42309eaab04b1d6b9a6c87b39b8e9735626d}\\
{\scriptsize\loc{logs/v11_hashes_before_v12.sha}} & \texttt{\tiny b8d08e5c00ece0b8988b6df73a0294246700b6ee14241b2614e364e0af92fa9c}\\
{\scriptsize\loc{logs/v12/earlier_versions_sha.json}} & \texttt{\tiny 93b6c79020f3ad890c9878cc20307aa56a913d0a456afd262b5b1beb67b22a0f}\\
{\scriptsize\loc{logs/v12/env.json}} & \texttt{\tiny f3c3b679dd87e7e87a4a5670b7e7dede9a62e737459a7a2a21d8fc562df914fa}\\
{\scriptsize\loc{CHECKPOINT_v12.md}} & \texttt{\tiny 64621e13a2f2765c05724debe7f5d481c93d51499e77ea3f5eff361aaeaf375a}\\
\end{longtable}}

\section{Source, version, and compilation declarations}
\label{app:declarations}

\paragraph{Supplied material.}
The v6 manuscript and the Python-listing manuscript were supplied
in full as text. Statements about their historical executions
are attributed to those sources.
In v8 the frozen XVI-P/XVI-Q engine files were read and replayed
from read-only copies; nothing in those trees was edited
(Section~\ref{sec:v8hash}). v7 was read in full.

\paragraph{Companion source scope.}
The I-B and I-B-G public records support the host and test-domain
comparisons cited above. The other Cheung companion descriptions
are cited at the version and locator level supplied with the
revision material. Their full manuscripts and proofs are not
represented as independently audited here.
In particular, the source-qualified tensor-core and
XVI-O1-ST comparisons are not silently upgraded to local proofs.

\paragraph{Primary analytical sources.}
The Bessel representation, standard special-function identities,
and self-adjoint functional calculus are identified inputs.
The normalizations, subtractions, Mellin reduction, and finite
catalogue algebra used in the conclusions are written out above.

\paragraph{Code status.}
v7: the embedded program is newly consolidated source.
v8: it was extracted unchanged and executed, together with the
v8 programs in \loc{code/v8/}; return codes, logs and SHA-256 values
are in Appendix~\ref{app:v8artifacts}. Software: Python~3 with SymPy
(versions in \loc{logs/v8/env.json}). No proof-assistant check is
asserted.
v9: programs in \loc{code/v9/}, logs in \loc{logs/v9/}, hashes in
Appendix~\ref{app:v9artifacts}; mpmath is additionally used for the
30-digit mode functionals.
v10: programs in \loc{code/v10/}, logs in \loc{logs/v10/}, hashes in
Appendix~\ref{app:v10artifacts}.
v11: programs in \loc{code/v11/}, logs in \loc{logs/v11/}, hashes in
Appendix~\ref{app:v11artifacts}.
v12: programs in \loc{code/v12/}, logs in \loc{logs/v12/}, hashes in
Appendix~\ref{app:v12artifacts}.

\paragraph{Typesetting.}
This is a single-file pdfLaTeX source with an internal
bibliography and inline listings. It requires no external figure,
bibliography database, code file, or shell escape for typesetting.
v12 was compiled with
\begin{verbatim}
pdflatex -halt-on-error XVI-O1-CONE-v12.tex   # repeated until stable
\end{verbatim}
and the final log was checked for unresolved references and
citations.

\paragraph{Authorship and approval.}
The author line follows the supplied manuscript.
This revision does not independently authenticate authorship,
archive provenance, or author approval.
No independent review of v8--v12 has taken place; nothing was pushed
or posted. Final submission requires the author's approval of the claims,
scope, versions, and computational artifacts.

\paragraph{Standing non-claims.}
The cone, deficit, holonomy, extension, and CR branch remain
prescribed inputs. No Einstein evolution, physical tensor
Hessian, spacetime emergence, gravitational evaporation channel,
or inherited characteristic quantum coefficient is constructed.
XVI-N is not discharged; XV-B is not opened; H2/I3/J2 are not
completed by this paper.

%======================================================================
\begin{thebibliography}{99}
\addcontentsline{toc}{section}{References}
%======================================================================
\item[] \textit{Marks (v12):} * read in full; \ddag{} confirmed online
(registry fetches of 18 October 2026, re-checked twice on 19 October 2026, \url{https://enter.grok.me/browse};
for the v9 sheets: source fetched, abstract and every theorem statement read);
\dag{} from memory or from a supplied description only.

\bibitem{CheungConeV9}
* N.~Cheung,
\emph{Scoped O.1 on the Prescribed Cone, Version 9},
18 October 2026. Source \loc{XVI-O1-CONE-v9.tex}
(SHA-256 \texttt{3a8dd434\ldots}); the base of this revision.

\bibitem{BaileyF4}
\dag{} W.~N.~Bailey, \emph{Generalized Hypergeometric Series},
Cambridge Tracts~32 (1935), reduction formulae for Appell's \(F_4\).

\bibitem{PSLQ}
\dag{} H.~R.~P.~Ferguson, D.~H.~Bailey and S.~Arno,
Analysis of PSLQ, an integer relation finding algorithm,
\emph{Math.\ Comp.}\ \textbf{68} (1999) 351--369.

\bibitem{CheungConeV11}
* N.~Cheung,
\emph{Scoped O.1 on the Prescribed Cone, Version 11},
19 October 2026. Source \loc{XVI-O1-CONE-v11.tex}
(SHA-256 \texttt{863bb08e\ldots}); the base of this revision.

\bibitem{CheungConeV10}
* N.~Cheung,
\emph{Scoped O.1 on the Prescribed Cone, Version 10},
18 October 2026. Source \loc{XVI-O1-CONE-v10.tex}
(SHA-256 \texttt{d21c8d84\ldots}); the base of this revision.

\bibitem{DLMFBeta}
\dag{} NIST Digital Library of Mathematical Functions, \S8.17
(incomplete Beta function) and \S15.4 (special values of \({}_2F_1\)).

\bibitem{CheungConeV8}
* N.~Cheung,
\emph{Scoped O.1 on the Prescribed Cone, Version 8},
18 October 2026. Source \loc{XVI-O1-CONE-v8.tex}
(SHA-256 \texttt{153e4e9d\ldots86e85f8}); the base of this revision.

\bibitem{CheungIBGSBv9}
\ddag N.~Cheung,
\emph{I-B-G-S-B: Positive tensor-core selection and validated
conical-kernel error bounds},
Enter Calculus Companion; registry source \loc{new_oct18f/all/I-B-G-S-B.tex}
(SHA-256 \texttt{a3ef5d15\ldots29864}).
Locators: \loc{thm:selection}, \loc{prop:warddefect},
\loc{prop:curvaturetable}. Cited input only.

\bibitem{CheungEchoV2}
\ddag N.~Cheung,
\emph{I-B-ECHO-v2: Diffracted echoes and finite supplied-report
instruments}, Enter Calculus Companion; registry source
(SHA-256 \texttt{9c3bfb55\ldots7a745}). Cited input only.

\bibitem{CheungXJST2}
\ddag N.~Cheung,
\emph{X-J-ST-2: Charged scalar reports and a Ward-compatible free
photon reference}, Enter Calculus Companion; registry source
(SHA-256 \texttt{3d0bae7f\ldots 08e81}). Cited input only.

\bibitem{CheungHawk05ST}
\ddag N.~Cheung,
\emph{Hawk05-ST: Finite Enter/Page Proxies on a Prescribed Conical
Spacetime}, registry source (SHA-256 \texttt{2cbd6c9f\ldots b7f8e}).
Not used for O.1.

\bibitem{CheungICv2}
\ddag N.~Cheung,
\emph{I-C-v2: The Shredder Interface}, registry source
(SHA-256 \texttt{6684c952\ldots 5d187}). Not used for O.1.

\bibitem{CheungICExv2}
\ddag N.~Cheung,
\emph{I-C-Ex-v2: Sharp Access Bounds, Residual Capacity}, registry
source (SHA-256 \texttt{8051b3e1\ldots 15a6f6}). Not used for O.1.

\bibitem{CheungICEx2}
\ddag N.~Cheung,
\emph{I-C-Ex2: Eight Further Extensions of the Shredder Cut},
registry identifier new on 18 October 2026, stored in
\loc{new_oct18g/} (SHA-256 \texttt{69833777\ldots 2647d}).
Not used for O.1.

\bibitem{CheungConeV7}
* N.~Cheung,
\emph{Scoped O.1 on the Prescribed Cone, Version 7},
18 October 2026. Source \loc{XVI-O1-CONE-v7.tex}
(SHA-256 \texttt{d0ba7b9c\ldots5d132}); the base of this revision.


\bibitem{CheungConeV6}
* N.~Cheung,
\emph{Scoped O.1 on the Cone, v6:
Executed Apex Records, a Closed Constant-Profile
Slavnov--Taylor Row, Weighted Apex Coefficients,
and Ward-Identity Boundaries},
18 October 2026.
Supplied complete source \loc{XVI-O1-CONE-v6.tex}.
Historical execution and manifest claims are attributed to this
source, not newly authenticated here.

\bibitem{CheungConeScripts}
* N.~Cheung,
\emph{XVI-O1-CONE: Python scripts (v1--v6)},
18 October 2026.
Supplied complete listing source
\loc{XVI-O1-CONE-scripts.tex}.
The listings state that frozen XVI-P/Q files and run logs are
external dependencies.

\bibitem{CheungP}
\ddag N.~Cheung,
\emph{XVI-P},
Enter Calculus Companion, unpublished manuscript and computational
records.
Frozen p22 contract as identified in the supplied cone audit;
later p23 excerpt is a comparison, not a silent contract change.
Named engine files:
\loc{code/p5_entry/amplitude_rules.py},
\loc{feynman.py}, and \loc{dred.py}.
Relevant supplied locators include
\loc{rem:o1}, \loc{rem:p5}, and \loc{rem:qds}.

\bibitem{CheungQ}
\ddag N.~Cheung,
\emph{XVI-Q},
Enter Calculus Companion, unpublished records.
Supplied Q02 graph-normalization description and separately
scoped Q11 release.
Named comparison artifacts include
\loc{Q02_bulk_graphs.json} and \loc{P18_D11.json}.
v8 inspected and replayed named engine files from a read-only copy.

\bibitem{CheungR}
\dag N.~Cheung,
\emph{XVI-R},
frozen comparison record identified in the supplied cone audit.
No theorem from its unavailable contents is imported here.

\bibitem{CheungIB}
\ddag N.~Cheung,
\emph{I-B: A Finite Enter Spacetime:
Calibrated Destructive Time Readout on a Prescribed Conical Product},
version 1, Enter Calculus Companion.
Relevant locators:
\loc{thm:instrument}, \loc{prop:not-source},
\loc{prop:no-return}, and \loc{prop:record-residual}.
Public manuscript record:
\url{https://enter.grok.me/browse/I-B}.

\bibitem{CheungIBG}
\ddag N.~Cheung,
\emph{I-B-G: Intrinsic and Holonomy-Twisted Green Operators on a
Prescribed Conical Spacetime},
version 2, Enter Calculus Companion.
Relevant locators:
\loc{prop:nonGH}, \loc{thm:green}, \loc{thm:speed},
and \loc{rem:compact-exact}.
Public manuscript record:
\url{https://enter.grok.me/browse/I-B-G}.

\bibitem{CheungIBGS}
\ddag N.~Cheung,
\emph{I-B-G-S: Intrinsic Spin-1 and Spin-2 Green Operators on the
Prescribed Conical Spacetime},
version 2, Enter Calculus Companion.
Supplied version-qualified excerpts, especially
\loc{thm:spin1}, \loc{thm:spin2},
\loc{def:geo}, and \loc{thm:selection}.
Free source-test compatibility is not an interacting loop theorem.

\bibitem{CheungTensor}
\ddag N.~Cheung,
\emph{I-B-G-S-B},
Enter Calculus Companion.
Supplied source descriptions of positive tensor-core selection
and exceptional source-test defects.
Relevant supplied locators:
\loc{cor:exclude}, \loc{thm:contrast}, and
\loc{prop:warddefect}.
A complete independent tensor-core proof is not asserted here.

\bibitem{CheungVP}
\ddag N.~Cheung,
\emph{I-B-VP: Vacuum Polarization of the Friedrichs Ground State
on a Prescribed Conical Spacetime},
version 3 and supplied predecessor passages,
Enter Calculus Companion.
Relevant supplied locators:
\loc{lem:infrared} and \loc{thm:friedrichs}.
One-point/compact-test comparison only.

\bibitem{CheungVP1}
\ddag N.~Cheung,
\emph{I-B-VP-1: Prescribed Apex Finite Parts on a Conical Enter
Spacetime: Wick Square, Stress Tensor, and Charged Current},
version 3, October 2026.
Supplied source-qualified descriptions of one-point extension
freedom and Ward conditions.

\bibitem{CheungVP4}
\ddag N.~Cheung,
\emph{I-B-VP-4},
Enter Calculus Companion.
Supplied frozen-source Maxwell-response comparison.
The current and its Maxwell response are distinct objects;
no interacting cone Ward theorem is imported.

\bibitem{CheungEx1B}
\ddag N.~Cheung,
\emph{I-B-VP-5-Ex1-B},
version 1, 18 October 2026.
Supplied excerpts on heat normalization, local stress running,
continuity of complete finite parts, and non-universality.
Relevant locators:
\loc{thm:running}, \loc{thm:continuity},
\loc{thm:equalradius}, and \loc{thm:trace}.

\bibitem{CheungEx2}
\ddag N.~Cheung,
\emph{I-B-VP-5-Ex2: Smooth Flux-Core Selection on a Prescribed Cone},
version 1, 18 October 2026.
Supplied excerpts on twisted scalar Friedrichs selection and
critical-window Dirac selection.

\bibitem{CheungEx2B}
\ddag N.~Cheung,
\emph{I-B-VP-5-Ex2-B},
version 1, 18 October 2026.
Supplied excerpts on exhausted-annulus Mosco arguments,
endpoint continuity, and finite-core qualifications.

\bibitem{CheungEcho}
\ddag N.~Cheung,
\emph{I-B-ECHO-v2: Diffracted Echoes and Finite Supplied-Report
Instruments},
version 2, October 2026.
Supplied source-qualified comparisons:
\loc{thm:readout}, \loc{thm:core},
and \loc{cor:nonreconstruction}.
No continuum acquisition or loop-residue identity is imported.

\bibitem{CheungIC}
\ddag N.~Cheung,
\emph{I-C-v2: The Shredder Interface:
Destructive Readout, Black-Hole Analogues, and Quantum Computers},
version 2, October 2026.
Supplied excerpts and source descriptions.
Theorems A--F separate referenced recovery, classical histories,
conditional thermal accounting, complete classical access,
endpoint underdetermination, and persistent classical output.

\bibitem{CheungEn09}
\ddag N.~Cheung,
\emph{En09-A: Realization of Relational Time by Destructive Enter
on a Clock Register},
version 1, Enter Calculus Companion.
Supplied excerpts, especially \loc{thm:main}.
Complete reference-preserving instrument and consumed-clock
conditioning; no continuum acquisition is inferred.

\bibitem{CheungEn07}
\ddag N.~Cheung,
\emph{En07-A}, \emph{En07-B}, and \emph{En07-C},
version 1 source records, Enter Calculus Companion.
Supplied complete-complement and classical-output comparisons.
Model-specific channels and physical endpoint assumptions are
not imported.

\bibitem{CheungXJ2}
\ddag N.~Cheung,
\emph{X-J-ST-2},
Enter Calculus Companion.
Identifier and \loc{prop:noacquisition} cited from the supplied
revision memorandum. No complete manuscript or additional
theorem is asserted to have been inspected.

\bibitem{CheungOai290}
\ddag N.~Cheung,
\emph{Oai290-A},
version 1, Enter Calculus Companion.
Supplied excerpts distinguish whole-code recovery, named
observable targets, and declared geometric models.
The companion identifier does not make an external catalogue
family a premise of this apex calculation.

\bibitem{CheungO1ST}
\ddag N.~Cheung,
\emph{XVI-O1-ST-v5},
Enter Calculus Companion.
Source-qualified sign-integrity witness described in the
supplied revision memorandum:
D01 sixth coordinate \(-4/3\) versus \(+4/3\), with
difference \((8/3)e_6\).
Full underlying artifacts were not supplied for this revision.

\bibitem{CheungHawk05}
\ddag \emph{Hawk05-A}, \emph{Hawk05-B}, and \emph{Hawk05-C},
version 1 supplied manuscript records,
attributed to N.~Cheung in the supplied source list.
Finite operational recovery, Page--Sen proxy, and conditional
extrapolation comparisons only.
Their authorship and full artifacts are not independently
authenticated here; no gravitational Page theorem is imported.

\bibitem{CheungN}
\ddag N.~Cheung,
\emph{XVI-N},
Enter Calculus Companion.
Supplied version-qualified boundary, including the identifier
\loc{thm:massless-obstruction}.
No absent theorem is reconstructed or discharged here.

\bibitem{CheungCharacteristic}
\ddag N.~Cheung,
\emph{XIV-H}, \emph{XIV-I}, and \emph{XIV-J},
version 1 source records, Enter Calculus Companion.
Supplied historical H2 ending, narrower I2 construction, and
unresolved full I3/J2 characteristic-map obligations.

\bibitem{CheungXVB}
\ddag N.~Cheung,
\emph{XV-B},
version 2, Enter Calculus Companion.
Supplied technical-interface descriptions.
The technical manuscript is distinguished from the unopened
downstream quantum-admission exit.

\bibitem{Fursaev1994}
\dag D.~V.~Fursaev,
``The heat kernel expansion on a cone and quantum fields near
cosmic strings,''
\emph{Classical and Quantum Gravity} \textbf{11} (1994), 1431.
DOI: \nolinkurl{10.1088/0264-9381/11/6/008}.
arXiv:\nolinkurl{hep-th/9309050}.
The full-opening-angle convention and equations (2.11)--(2.15)
fix the comparison in Remark~\ref{rem:Fursaevnorm}.

\bibitem{FursaevMiele}
\dag D.~V.~Fursaev and G.~Miele,
``Cones, spins and heat kernels,''
\emph{Nuclear Physics B} \textbf{484} (1997), 697--723.
DOI: \nolinkurl{10.1016/S0550-3213(96)00631-1}.
arXiv:\nolinkurl{hep-th/9605153}.

\bibitem{Vassilevich}
\dag D.~V.~Vassilevich,
``Heat kernel expansion: user's manual,''
\emph{Physics Reports} \textbf{388} (2003), 279--360.
DOI: \nolinkurl{10.1016/j.physrep.2003.09.002}.
arXiv:\nolinkurl{hep-th/0306138}.

\bibitem{DLMF}
\dag F.~W.~J.~Olver et al., editors,
\emph{NIST Digital Library of Mathematical Functions},
Chapter 10, especially Sections 10.27, 10.29, 10.32, and 10.43.
Modified-Bessel relations, recurrences, and integral transforms.
\url{https://dlmf.nist.gov/10}.

\bibitem{DerezinskiRichard}
\dag J.~Derezi\'nski and S.~Richard,
``On Schr\"odinger operators with inverse square potentials on
the half-line,''
\emph{Annales Henri Poincar\'e} \textbf{18} (2017), 869--928.
DOI: \nolinkurl{10.1007/s00023-016-0520-7}.
arXiv:\nolinkurl{1604.03340}.

\bibitem{KayStuder}
\dag B.~S.~Kay and U.~M.~Studer,
``Boundary conditions for quantum mechanics on cones and fields
around cosmic strings,''
\emph{Communications in Mathematical Physics}
\textbf{139} (1991), 103--139.
DOI: \nolinkurl{10.1007/BF02102731}.

\bibitem{Stockinger}
\dag D.~St\"ockinger,
``Regularization by dimensional reduction:
Consistency, quantum action principle, and supersymmetry,''
\emph{Journal of High Energy Physics} \textbf{2005}(03), 076.
DOI: \nolinkurl{10.1088/1126-6708/2005/03/076}.
arXiv:\nolinkurl{hep-ph/0503129}.

\bibitem{HSR}
\dag M.~Horodecki, P.~W.~Shor, and M.~B.~Ruskai,
``Entanglement breaking channels,''
\emph{Reviews in Mathematical Physics} \textbf{15} (2003),
629--641.
DOI: \nolinkurl{10.1142/S0129055X03001709}.
Preprint title: ``General Entanglement Breaking Channels,''
arXiv:\nolinkurl{quant-ph/0302031}.

\bibitem{NielsenAverage}
\dag M.~A.~Nielsen,
``A simple formula for the average gate fidelity of a quantum
dynamical operation,''
\emph{Physics Letters A} \textbf{303} (2002), 249--252.
DOI: \nolinkurl{10.1016/S0375-9601(02)01272-0}.
arXiv:\nolinkurl{quant-ph/0205035}.

\bibitem{ReebWolf}
\dag D.~Reeb and M.~M.~Wolf,
``An improved Landauer principle with finite-size corrections,''
\emph{New Journal of Physics} \textbf{16} (2014), 103011.
DOI: \nolinkurl{10.1088/1367-2630/16/10/103011}.
arXiv:\nolinkurl{1306.4352}.

\end{thebibliography}

\end{document}